\documentclass[12pt,a4paper]{article}
\usepackage{graphicx}
\usepackage{slashed}
\usepackage{tikz}
\usepackage{colortbl}
\usepackage{calc}


% Current II-D reference root for PapersMinimal v3.063; mathematical macros otherwise unchanged.
% Macros-v2_188: re-based current paper identities and A--G family; no change to geometry.
% Macros-v2_184: active Paper III-E v1.200/III-F v1.267 binding for framework v3.031; mathematical macros unchanged.
% Macros-v2_183: updated shared cross-reference bindings for Paper1B/III-E/III-F v3.030; mathematical macros unchanged.
% Macros-v2_182: active Paper II-F/V-A cross-references for framework v3.029; definitions unchanged.
% Paper-series v2.381: new I-B torsion/modular results, bibliography and current source roots.
% Paper-series v2.380: new exceptional bibliography and synchronized paper references.
% Paper-series v2.379 current 31-paper macro/aux registry.
% Paper-series v2.378: current 31-paper AUX/bibliography registry updated for Paper I-B v1.2 literature additions.
% Framework v2.371: current-source reference registry refreshed; old theorem labels retained.
% Framework v2.370: current 30-paper source/AUX root rebinding; native mathematical macros unchanged.
% Framework v2.368: synchronized merged 30-root registry; mathematical notation unchanged.
% Framework v2.358: active 30-paper source registry synchronized; existing mathematical macros unchanged.
% Framework v2.357: selective new source-owned certificates; active reference roots updated for five revised papers.
% Framework v2.356: active 30-root registry synchronized to the 30 papers in this release.
% Framework v2.275: exact positive dense-twin word cubature, finite repeatability and eighth-moment classification, nested sampling frames, joint action and POVM-jet exports; preserves v2.274.
% Framework v2.274: filename and root-source synchronization.
% Framework v2.273: restored missing Problem 3 theorems from v2.268; exact exceptional frames, twin geometry, convolution repeatability and projector/occurrence firewalls; merged on v2.272 without overwriting independent branches.
% Framework v2.272: merged v2.270 covariant native-readout mathematics with v2.271 Problem 4 finite-action Hopf hypercharge/holonomy corrections; no universal bulk-gauge uniqueness claim.
% Independent branches: v2.270 readout/guard exports and v2.270 K16 action-consumed charge closure; all prior scoped assumptions retained.
% Framework v2.269: synchronized 30 active AUX roots; mathematics and bibliography preserved from Macros-v2_161.
% Framework v2.266: Paper II-C converse-derived outcome resources; II-E, III-C, IV-A/B, V-A and VI-B downstream mathematical closures; 30-root registry synchronized.
% Framework v2.265: II-C coordinate-stage constitution and exact first-symbol host inheritance; prior mathematical domains retained.
% Framework v2.264: paper II-B reverse/audit and IV-A observable-minimal exchange; synchronized external references to 30 active roots.
% Framework v2.263: Paper I-C adds bridge-controlled root-carry parity and source-relative old-cap cancellation; other root mathematical bodies retained; active 30-root references synchronized.
% Framework v2.262: integrates Paper II-F radius, Casimir and readout results; active 30-root references synchronized.
% Framework v2.261: merged Paper II-B GP2-level/geometric repairs; preserves v2.260 II-D repairs; active references synchronized.
% Framework v2.260: Paper II-D signed RS, finite-symbol typing, and conditional curvature/Jordan repairs; references synchronized.
% Framework v2.259: conditional native SU(2) atlas and explicit later discharge; v2.258 mathematical bodies retained elsewhere.
% Framework v2.258: preserves parallel II-A/dynamics update and adds complete readout/frequency closures.
% Framework v2.257: exact minimal complete position readout and arbitrary-preparation frequency law.
% Framework v2.256: exact native inertia and sharp bounded ordered-source closure; preserves concurrent v2.255 work.
% Framework v2.255: preserves parallel detector integration and adds exact finite readout closures.
% Framework v2.254: Pasterski bibliography and synchronized 30-paper external-reference registry.
% Framework v2.253: typed raw coordinate defect in I-A and synchronized VI-A import; active references synchronized.
% Framework v2.252: synchronized source registry; optional III-G shell-frame calculation.
% Framework v2.251: forward common-action execution, joint source jets and scoped coupled control.
% Framework v2.251: exact full-heat variations, fermion source jets and explicit physical completion frontier.
% Framework v2.251: certified interacting source jets, pole stability and scoped completion proofs.
% Framework v2.251: synchronized active source registry; mathematical notation unchanged.
% Framework v2.251: synchronized active source registry; mathematical notation unchanged.
% Framework v2.251: synchronized active source registry; mathematical notation unchanged.
% Framework v2.251: synchronized source registry; no mathematical notation changes.
% Framework v2.244: synchronized 30-paper AUX registry and rooted phase-source action-order tests.
% Framework v2.242: signed rooted phase lift, complete transverse return and exact family cubic.
% Framework v2.241: proved rooted conormal geometry, conditional cap realization and scoped action tests.
% Framework v2.239: all 30 active source roots synchronized; no mathematical macro changes.
% Framework v2.238: Paper IV A/B/C split and 30-root reference registry.
% Framework v2.237: Paper VII-D technical split and 28-root reference registry.
% Framework v2.236: Paper VI-C technical-volume reference added; all current root bindings advanced.
% Framework v2.235: 26-root AUX and bibliography bindings synchronized; mathematical macros unchanged.
% BEGIN INLINE SOURCE: Paper00-TeXBuildCompatibility-v1_0.tex
% Rendering-only repairs for inherited long appendices and bitmap text fonts.
% Integer counters through 26 retain the original labels; later values use AA,
% AB, ... . This file changes no action coefficient, state or theorem premise.
\RequirePackage{alphalph}
\makeatletter
\let\BHPrimitiveAlph\@Alph
\newalphalph{\BHExtendedAlph}{\BHPrimitiveAlph}{26}
\let\@Alph\BHExtendedAlph
\AtBeginDocument{%
  \@ifpackageloaded{microtype}{\microtypesetup{expansion=false}}{}%
}
\makeatother

% END INLINE SOURCE: Paper00-TeXBuildCompatibility-v1_0.tex

% Framework v2.164: current reference/bibliography profile; mathematical definitions retained.
% Framework v2.162: finite metric, functional Hilbert and scoped Ward integration.
% Framework v2.161: inherits Macros-v2_57.tex exactly except the current external-reference binding.
% Macros v2.57 / framework v2.156: current reference registry only; mathematical definitions unchanged.
% Framework v2.143 / Macros v2.55: current external-reference bindings; mathematical notation unchanged.
% Framework v2.110 / Macros v2.54: current external-reference registry and bibliography; mathematical notation inherited.
% Framework v2.67 / Macros-v2_53: adds Born-action-cell and joint-EH readout symbols; absolute unit realization is metrological; MATCH_{J/scal}^{EH} remains the sole G24 normalization bridge.
% Framework v2.35 / Macros-v2_52: grouped paper numbering, active AUX registry, and Born-v6_20 bibliography binding synchronized; notation unchanged.
% Framework v2.26: active-source AUX and bibliography bindings synchronized; notation unchanged.
% Framework v2.09: reviewed early-reciprocity integration; scoped physical imports.
% Framework v2.08: Born--Euler mathematics and conditional complete-action potential criterion.
% Framework v2.07: active AUX registry includes Paper I-B v1.2; shared bibliography Born-v6_18.
% Macros-v2_46 / framework v2.04: mathematical atlas and downstream realization correction.
% Macros-v2_46 / framework v2.03: scoped structure integration and current dependency bindings.
% Macros-v2_46 / framework v2.01: shared source bindings; prior revision comments are historical.
% v2.42 / framework v1.99: Synchronized RG/AS owner and bibliography bindings.
% v2.41 / framework v1.99: active manuscript AUX bindings and geometric-closure bibliography synchronized.
% Release-v1.80 centralized source synchronization for scalar-circle/comparison-band integration.
% Release v1.79: Hurwitz hierarchy and kinetic-normalized action response integration; physical mass matching remains conditional.
% Release v1.78: approved M8 closure-batch integration; physical and analytic domains retained.
% Current v1.78 registry: all active TeX roots are synchronized below.
% Earlier release comments are historical scope records, not current status.
% Release-v1.77 synchronization: native same-owner Noether closure, generator--flow/history typing, analytic-certification firewalls, scalaron/K16, LF/QCD, quantum-Ward, and RAR boundary-owner transport are registered; shared mathematical notation is unchanged.
% Release-v1.76 synchronization: registry points to the v1.76 active paper roots.
% Release-v1.76 synchronization: registers the local A+Phi+g Noether/constraint closure and dynamics-frontier propagation; shared mathematical notation is unchanged.
% Release-v1.75 synchronization: registry points to the v1.75 active paper roots.
% Release-v1.75 synchronization: integrates the joint-Dirac Einstein bulk real form and the local odd-scalar heat-coefficient action on explicit analytic first-jet/interior domains; v1.74 Yang--Mills closure is retained. Fermionic reality, pure-metric curvature-squared/exact-nonlocal terms, noncanonical boundaries, physical vacuum selection, and actual solution-history tokens remain open.
% Release-v1.74 synchronization: registry points to the v1.74 active paper roots and retains the v1.73 theorem labels.
% Release-v1.74 synchronization: canonical Macros-v2_34.tex; generated local nonlinear Yang--Mills action, constrained connection-field phase space, source-free Gauss momentum map, and Legendre closure are integrated on their proved domains. The remaining non-gauge ACT_cl sectors and actual solution-history selection remain open.
% Release-v1.71 synchronization: closes the former JOINT_x/JOINT_5
% connection-occurrence gap by canonical structure-group extension of the
% nonflat Born connection, faithful associated representation, and the local
% cotangent--Clifford operator constructor.  Actual fluctuation members,
% actions, vacua, solutions, calibrations, and analytic QFT domains remain
% separately typed.  The active-AUX registry is advanced once.
% Release-v1.71 canonical macro and active-AUX registry.
% ======================================================================
% Macros-v2_39.tex -- canonical publication macro file
% Release v1.71 retains the v1.70 null-ray/QCD boundaries and additionally
% removes the duplicate structural connection actuality premise.  It does not
% claim global bundle topology, select a gauge/Higgs member, or promote the
% finite construction to an interacting continuum theory.
% Release v1.68 constructs the finite light-front total-fraction operators and
% parton number/momentum operator measures, proves their simplex and Hausdorff-limit
% theorems, and separates positive cut measures from scheme-renormalized
% light-ray distributions.  It also separates finite real Schur poles, stable
% norm-resolvent/Riesz pole-residue limits, and continued resonance
% pole/residue limits.  The Paper-VII-C scattering package is claim-indexed:
% local QFT and spectral claims import no S-matrix or infrared premise, while
% asymptotic amplitudes, LSZ, S-matrix poles, completeness, and cross sections
% import only their named massive or massless/long-range clauses.
% Release v1.67 factorizes BND48_H into a same-action Hessian/Schur
% descendant and still-open typed boundary-trace, critical-lift, and
% refinement-limit conditions.  It proves the finite-collar Dirichlet-to-
% Neumann law sqrt(y)tanh(L sqrt(y)), its nonuniform critical limit, and the
% finite-return ultra-infrared turnover.  The Paper-I-A Hardy constant mode is
% an exact geometric near-hit, not yet the physical metric boundary mode.
% The current local f(R) Hessian remains massive and therefore still refutes
% a RAR prediction on the minimal action branch.
% Release v1.66 derives the transverse depth-six normalizer bridge from the
% already generated C6 Reynolds image and Paper-I-A binary-dihedral normalizer.
% The bridge is the unique least 2D3-stable scalar-containing realization of
% the faithful quotient 2D3/C6=C2, has trace defect six and hence K=48, but is
% transverse to the nested active carrier.  Exact K=44 and K=52 contraction,
% semigroup, and energy-midpoint identities are registered as structural
% readouts rather than projector events.  The active E48 carrier gap remains
% unchanged.  A same-action boundary-symbol/Calderon intertwiner is isolated
% as the necessary and still unproved owner bridge to the static DtN sort;
% representation equivalence alone neither derives RAR nor creates a K=48
% action junction, pole, occupation, or cosmological transition.
% Release v1.65 derives the native E_48-centered carrier mass gap from the
% GP1--GP3-generated active spectrum, distinguishes its carrier-free interval
% from a particle located at E_48, and withdraws the unsupported promotion to
% a positive continuum Yang--Mills mass gap.  The latter remains a valid
% conditional transfer theorem on an enlarged continuum branch, but it is not
% an active framework prediction or an unfinished native-gap proof.
% Release v1.64 synchronizes Derivation Batches 14--22.  It registers the
% GP1--GP3 heat-profile orbit, same-Hessian electroweak response and pole
% ownership, K16/K48 action-response corrections, perturbation state and
% non-Gaussianity typing, finite-vacuum/gap/QCD limit repairs, factorized
% position-area ownership, the actual static metric response, and the ordered
% Wilson/string-tension limits.  Strict continuum convergence, positive
% physical gap, surviving higher cumulants, and positive tension remain open.
% Release v1.61 closes TIMEWIDTH16 as an independent direct-production datum.
% JUNC_{2|7} continues to own the finite intrinsic microscopic profile/width,
% while the actual same-action background/production history determines its
% physical-time pullback, support, and duration.  The current Paper-I-A--VII-C AUX
% registry is advanced once; shared mathematical notation is unchanged.
% Release v1.60 derives the final-octet primitive-remnant depletion, canonical
% Hopf-current representative, exact Hopf descent to dOmega/(4pi), and
% closed-FLRW a^{-3} neutral-current shape.  It retypes K=48 as the onset
% of the final one-unit geometric transfer and records cDM/RAR as distinct
% physical projections of that common neutral parent while preserving
% occurrence, Hilbert-owner, stress-matching, and static-kernel premises.
% Release v1.59 transports the v1.58 exceptional-shell/K16 closure through
% Papers III-F and VI-B--VII-C.  It replaces Paper VI-B's obsolete mixed-Hessian
% output typing by the third-variation pair vertex and Bogoliubov occupation
% support, derives the post-seed Higgs-pair feedback vertex and its exact
% alpha--delta kernel factors, propagates the unique K16/K48 action-junction
% split, identifies the QCD delta count with the upstream 2O shell atom only
% as a common-scalar numerical compatibility, and narrows THERM16 to dynamical
% approach/beta selection because the finite Gibbs/KMS target is already
% unique on a supplied simple observable algebra once its Hamiltonian and
% inverse temperature are supplied; direct-sum central weights remain typed.
% Release v1.58 registers the exceptional nearest-shell closure, the derived
% transverse K=16 action event, the unique interior action junction, the
% type-correct ordered pair vertex, the thin-layer squeeze law, and the
% radiation-like/matter-like kinematic split.  REG_0, transfer magnitude,
% physical layer width, and thermalization remain explicitly conditional.
% Release v1.57 registers the active-depth-six obstruction, the native
% K=48 carrier desert, the typed 7|9 owner reduction, and the dynamical
% reheating handoff across Papers III-B--VI-B.
% Release v1.56 registers the least-support n_2 inflation dressing of the
% persistent n_7 Laurent response, the analytic global-convexity and positive
% endpoint-measure proofs, the derived inflation core and K=16 coefficient
% freeze, and the no-particle dark-sector split into static RAR plus the
% conditional homogeneous effective-curvature response at reflected K=48.
% There is no n_6 action, scalaron-DM mass, or second Einstein--Hilbert term.
% Release v1.55 registers the protected-octet resolved/residual reflection
% q_sp+Z_W=8, its j<->8-j and K<->64-K scale covariance, and the derived
% rank-one-exit quadratic grade C_1^2/U=C_2 in Papers III-C, III-E, and VIII.
% The reflection is a typed trace/scale theorem, not a time reversal or a
% physical event-occurrence theorem; reheating transfer remains conditional.
% Release v1.54 registers the exact positive-part Laurent indices, the
% active-domain/terminal-endpoint action split, the no-second-EH firewall, and
% the q_DM=6/K_DM=48 scalaron readout typing across Papers III-D/E/F, VI,
% VIII, and IX.
% Release v1.53 registers the common nonnegative Laurent obstacle projection
% for n_1^+ and n_7^+, discharges the three former exit/freeze assumptions,
% and advances the Paper-III-D/E/F, VIII, and IX consumers.  Shared notation
% and bibliography ownership are unchanged.
% Release v1.51 registers the PQTI derivation batch and its updated
% Paper-III-A/C/D/E/F, VIII, and IX consumers.  Shared notation and
% bibliography ownership are unchanged.
% Release v1.50 advances the Paper-III-B AUX target to v1.40 after the
% rigorous audit repair; shared notation and bibliography ownership are unchanged.
% Release v1.49 advances the Paper-II-E/F AUX target to v3.14 after the
% minimal audit repair; shared notation and bibliography ownership are unchanged.
% Release v1.46 synchronizes Paper II-A's minimal frozen constructor-origin
% theorem and its five dependent AUX targets.  Shared notation is unchanged.
% Release v1.45 updates the Paper-II-B AUX target to v2.23 after the
% GP3-rigidity and appendix-placement repair; shared notation is unchanged.
% Release v1.44 advances the centralized Paper-I-A--VII-C AUX registry after the
% Paper-I-A v5.16 protected-completion scope repair.  Every current
% paper imports this one file; no paper-level opt-in or \externaldocument list
% is permitted.  The bibliography remains centralized here through Born-v6_13.
% Shared bibliography, theorem infrastructure, colors, notation alignment,
% and current Paper-I-A--VII-C external-reference registration.
% This file centralizes common imports, bibliography sorting, colors, theorem
% infrastructure, and the notation shared by the complete series.
% v2.9 (release v1.43) made the Paper-I-A--VII-C AUX registry automatic and
% centralized the bibliography and all external-reference ownership.
% v2.8 release-registry refresh (v1.33) registers the rooted generated-parent
% cascade and the POL/REG/free-law/11+4/Born-group closure revisions.  No
% shared mathematical notation or Paper-II-A fundamental principle is changed.
% v2.8 release-registry refresh (v1.32) registers the universal generated-
% continuation revisions in Papers II-A, III-A, III-C, IV, and VIII.  No
% shared mathematical notation or Paper-II-A fundamental principle is changed.
% v2.8 release-registry refresh (v1.30) registers the CRP4/projective-state
% source revisions and their direct registry-only consumers.  No shared
% mathematical notation or Paper-II-A axiom is changed.
% v2.7 release-registry refresh (v1.27) updates only the current external
% Paper-I-A--VII-C basenames after the placement/assumption synchronization; no
% shared mathematical notation or Paper-II-A axiom is changed.
% v2.7 synchronizes the geometric defect-repair revision: Paper-II-A v3.42,
% Paper-III-C v1.26, Paper-III-D v1.27, Paper-III-E v1.32, and Paper-III-F
% v1.34.  It registers the common trace-defect/spectral-filtration theorem,
% its owner-typed response pullback, and the reduced GRED and TERM interfaces.
% No shared mathematical notation is changed.
% v2.6 synchronizes Paper-III-E v1.31's cone-type repair, minimal-response
% classification, and explicit classical action-role boundary, together with
% Paper-III-F v1.32's downstream response-atlas synchronization.  No shared
% mathematical notation is changed.
% v2.5 synchronizes the Paper-III-D v1.26 faithful-trace-duality proof,
% exact ordered-owner joint locus, and typed alternative-presentation
% bookkeeping repair, together with its Paper-III-E and closure-frontier
% references.  No shared mathematical notation is changed.
% v2.4 synchronizes the Paper-III-C v1.25 orbit/process/owner repair, its
% downstream Papers III-B--F, V, VI, VIII, and IX revisions, and Born-v6_13.
% v2.3 synchronizes the Paper-II-A v3.41 witness/support/principal inheritance
% layer, the Paper-III-A v1.24 joint-realizability/readout firewall, all dependent
% host-rate revisions, and the current Paper-I-A--VII-B source registry.
% v2.2 synchronizes the generated active-realization, determinant-depth,
% dimension-weighted hierarchy, and typed cosmological-event interfaces.
% v2.1 synchronizes the exceptional-support interface revisions: Paper I-A
% v5.9, Paper II-A v3.30, Paper II-E/F v3.7, and Paper IV v4.16.
% v2.0 introduces the stage-safe G0 occurrence notation, synchronizes the
% active Paper-I-A--IV source registry (current rigorous Paper-III closure:
% III-A v1.17, III-B v1.22, III-C--E v1.17, III-F v1.14, IV v4.14), and makes
% external-reference self-skip depend on logical paper role rather than a
% versioned job name.
% ======================================================================

% ======================================================================
% macros.tex — Shared macros for the Born Geometry series
% ======================================================================

% ======================================================================
% 1. PACKAGES, BIBLIOGRAPHY, COLORS, AND UNIT DEFINITIONS
% ======================================================================
% Core packages shared across the complete Paper-I-A--VII-C series.  Loading them here makes
% the macro file a genuine series preamble while remaining compatible with existing
% paper preambles that load the same packages before \input.
\usepackage{amsmath}
\usepackage{amssymb}
\usepackage{amsthm}
\usepackage{mathrsfs}
\usepackage{array}
\usepackage{booktabs}
\usepackage{xcolor}
\makeatletter
\@ifpackageloaded{geometry}{}{\usepackage{geometry}\geometry{margin=1in}}
\makeatother
\usepackage{mdframed}

% Shared support packages required by the macros below and by most papers.
\usepackage{ragged2e}
\usepackage{longtable}
\usepackage{tabularx}
\usepackage{enumitem}
\usepackage[most]{tcolorbox}

% Series-wide numeric bibliography convention using the portable BibTeX path.
% All Papers I-A--VII-C use the ordinary \cite interface.
\usepackage[numbers,sort&compress]{natbib}
\providecommand{\printbibliography}{%
  \bibliographystyle{plainnat}%
  \bibliography{Born-v6_59}%
}

% Shared colors.  \providecolor avoids clobbering a paper-specific override.
\providecolor{postulatecolor}{RGB}{230,242,255}
\providecolor{compare}{RGB}{230,245,230}
\providecolor{definitioncolor}{RGB}{230,242,255}

\providecolor{openriskcolor}{RGB}{255,235,235}
\providecolor{bridgecolor}{RGB}{240,240,255}
\providecolor{gold}{RGB}{255,215,0}

% ======================================================================
% 2. CUSTOM ENVIRONMENTS
% ======================================================================
% Common table column types used in all papers.  Guarded so that a paper may
% define its own versions before loading this file.
\makeatletter
\@ifundefined{NC@find@L}{%
  \newcolumntype{L}[1]{>{\RaggedRight\arraybackslash}p{#1}}%
}{}
\@ifundefined{NC@find@C}{%
  \newcolumntype{C}[1]{>{\Centering\arraybackslash}p{#1}}%
}{}
\@ifundefined{NC@find@Y}{%
  \newcolumntype{Y}{>{\RaggedRight\arraybackslash}X}%
}{}
\makeatother

% ======================================================================
% 3. THEOREM STYLES AND ENVIRONMENTS
% ======================================================================
% Custom underlined theorem style (useful for remarks or notes)
\newtheoremstyle{underlined}
  {3pt}   % Space above
  {3pt}   % Space below
  {}      % Body font (empty means normal/roman text, typical for remarks)
  {}      % Indent amount
  {}      % Theorem head font (can be \itshape, \bfseries, etc.)
  {.}     % Punctuation after theorem head
  { }     % Space after theorem head
  {\underline{\thmname{#1}\thmnumber{ #2}}\thmnote{ (#3)}} % Theorem head formatting

% Standard plain theorem style environments.  Each environment is guarded so
% the shared macro file can be loaded after a paper-specific theorem setup.
\theoremstyle{plain}
\makeatletter
\@ifundefined{theorem}{\newtheorem{theorem}{Theorem}[section]}{}
\@ifundefined{lemma}{\newtheorem{lemma}{Lemma}[section]}{}
\@ifundefined{proposition}{\newtheorem{proposition}{Proposition}[section]}{}
\@ifundefined{corollary}{\newtheorem{corollary}{Corollary}[section]}{}
\makeatother

% Definition style environments (typically non-italicized body)
\theoremstyle{definition}
\makeatletter
\@ifundefined{definition}{\newtheorem{definition}{Definition}[section]}{}
\@ifundefined{postulate}{\newtheorem{postulate}{Postulate}[section]}{}
\@ifundefined{closure}{\newtheorem{closure}{Closure}[section]}{}
\makeatother

% Applying the custom underlined style
\theoremstyle{underlined}
\makeatletter
\@ifundefined{remark}{\newtheorem{remark}{Remark}[section]}{}
\@ifundefined{convention}{\newtheorem{convention}{Convention}[section]}{}
\@ifundefined{principle}{\newtheorem{principle}{Principle}[section]}{}
\makeatother

% Framed proof environment.
% Every \begin{proof} ... \end{proof} is automatically framed/shaded.
% This avoids mdframed/trivlist nesting, which can drop the top border after
% theorem-like environments.
\makeatletter
\renewenvironment{proof}[1][\proofname]{%
  \par\addvspace{6pt}%
  \noindent
  \begin{tcolorbox}[%
    enhanced,
    breakable,
    colback=black!4,
    colframe=black!35,
    boxrule=0.35pt,
    arc=1pt,
    left=6pt,
    right=6pt,
    top=4pt,
    bottom=4pt,
    boxsep=0pt,
    before skip=0pt,
    after skip=6pt,
    parbox=false
  ]%
  \pushQED{\qed}%
  \normalfont\itshape #1\@addpunct{.}\enspace
  \normalfont\ignorespaces
}{%
  \unskip\nobreak\hfill\popQED
  \end{tcolorbox}%
  \par\addvspace{2pt}%
}
\makeatother

\mdfdefinestyle{postulatebox}{
  backgroundcolor=yellow!10,
  linecolor=black,
  linewidth=0.4pt,
  skipabove=\baselineskip,
  skipbelow=\baselineskip,
  nobreak=true
}

% ======================================================================
% PROOF-READABILITY HELPERS
% ======================================================================

\newcommand{\ProofPlan}[1]{%
  \par\smallskip
  \noindent\textbf{Proof plan. }#1%
  \par\smallskip
}

\newlist{proofsteps}{enumerate}{2}
\setlist[proofsteps,1]{
  label=\textbf{(\alph*)},
  ref=\alph*,
  leftmargin=2.2em,
  itemsep=0.45\baselineskip,
  topsep=0.35\baselineskip
}
\setlist[proofsteps,2]{
  label=\textbf{(\roman*)},
  ref=\roman*,
  leftmargin=2.0em,
  itemsep=0.25\baselineskip,
  topsep=0.2\baselineskip
}


% ======================================================================
% 4. NUMBER SYSTEMS AND ALPHABETS
% ======================================================================
\providecommand{\RR}{\mathbb{R}} % Real numbers
\providecommand{\PP}{\mathbb{P}} % Projective space / Probability
\providecommand{\CC}{\mathbb{C}} % Complex numbers
\providecommand{\ZZ}{\mathbb{Z}} % Integers
\providecommand{\NN}{\mathbb{N}} % Natural numbers
\providecommand{\HH}{\mathbb{H}} % Quaternions / Hyperbolic space
\providecommand{\LFRayGeom}{\mathsf{LF\!\mbox{-}RAY\!\mbox{-}GEOM}}
\providecommand{\LFRep}{\mathsf{LF\!\mbox{-}REP}}
\providecommand{\LFState}{\mathsf{LF\!\mbox{-}STATE}}
\providecommand{\LFLim}{\mathsf{LF\!\mbox{-}LIM}}
\providecommand{\LFScat}{\mathsf{LF\!\mbox{-}SCAT}}
% Number systems / projective spaces
\providecommand{\CP}{\mathbb{CP}}
\providecommand{\RP}{\mathbb{RP}}

% ======================================================================
% 5. GROUPS AND LIE THEORY
% ======================================================================
\providecommand{\SO}{\ensuremath{\mathrm{SO}}}     % Special Orthogonal group
\providecommand{\SU}{\ensuremath{\mathrm{SU}}}     % Special Unitary group
\providecommand{\Spin}{\ensuremath{\mathrm{Spin}}} % Spin group
\providecommand{\CS}{\ensuremath{\mathrm{CS}}}     % Chern-Simons (capitalized)
\providecommand{\cs}{\ensuremath{\mathrm{cs}}}     % Chern-Simons (lowercase)

% ======================================================================
% 6. COMMON MANIFOLDS AND TOPOLOGY
% ======================================================================
\providecommand{\Sone}{S^1}                 % Standard circle

% ======================================================================
% 7. MATH OPERATORS (Declared via \DeclareMathOperator)
% ======================================================================
\DeclareMathOperator{\End}{End}   % Endomorphisms
\DeclareMathOperator{\Aut}{Aut}   % Automorphisms
\DeclareMathOperator{\spec}{spec} % Lowercase spectrum
\DeclareMathOperator{\tr}{tr}     % Lowercase trace (if you already have \Tr)
\DeclareMathOperator{\supp}{supp} % Support
\DeclareMathOperator{\TrA}{Tr_{A}} % Partial trace over subsystem A
\DeclareMathOperator{\TrB}{Tr_{B}} % Partial trace over subsystem B
% Lie/algebra operators
\DeclareMathOperator{\Ad}{Ad}
\DeclareMathOperator{\Sym}{Sym}
\DeclareMathOperator{\su}{su}
% ======================================================================
% 8. STANDARD OPERATOR-LIKE SYMBOLS (Safe definitions via \providecommand)
% ======================================================================
\providecommand{\vol}{\operatorname{vol}}   % Volume form
\providecommand{\Vol}{\operatorname{Vol}}   % Total volume
\providecommand{\Ric}{\operatorname{Ric}}   % Ricci tensor
\providecommand{\Rm}{\operatorname{Rm}}     % Riemann tensor
\providecommand{\Tr}{\operatorname{Tr}}     % Uppercase trace
\providecommand{\Spec}{\operatorname{Spec}} % Uppercase spectrum
\providecommand{\diag}{\operatorname{diag}} % Diagonal matrix
\providecommand{\rank}{\operatorname{rank}} % Rank
\providecommand{\scal}{\mathcal R} % Scalar curvature
\providecommand{\Cl}{\operatorname{Cl}}     % Clifford algebra
\providecommand{\diam}{\operatorname{diam}} % Diameter

% ======================================================================
% 9. FUNCTIONS, TRIGONOMETRY, AND CALCULUS
% ======================================================================

% Codifferential naming conventions
% (Avoid \ddag: commonly defined by physics.sty and symbol packages)
\providecommand{\codiff}{\mathrm{d}^{\dagger}} % Codifferential operator

% Quantum mechanics bra-ket notation (needed for basis tables)
\providecommand{\ket}[1]{\lvert #1 \rangle}
\providecommand{\bra}[1]{\langle #1 \rvert}
\providecommand{\braket}[2]{\langle #1 \mid #2 \rangle}

% ======================================================================
% 10. PHYSICAL SCALES AND CONSTANTS (Paper 1c / Born Geometry)
% ======================================================================
\providecommand{\Rp}{\ensuremath{R_{\mathrm P}}}     % matched common EH curvature scale (not defined from observed G)
\providecommand{\hBorn}{\ensuremath{h_{\mathrm B}}} % physical representative 2\pi\hbar of one primitive Born scalar action cell
\providecommand{\UEHjoint}{\ensuremath{U_{J}}} % joint-Dirac EH curvature-coordinate scale 8 kappa_J
\providecommand{\Gread}{\ensuremath{G_{\rm scal}}} % scalar-action coordinate c_op^3/(h_B U); physical matching explicit
\providecommand{\GreadJoint}{\ensuremath{G_{J,\mathrm{raw}}}} % raw ordinary-trace coordinate = G_scal/96
\providecommand{\Rds}{R_{\mathrm{dS}}}    % de Sitter Curvature
\providecommand{\Mpl}{M_{\mathrm{Pl}}}    % Planck mass
\providecommand{\lp}{\ell_{\mathrm{Pl}}}  % Planck length
\providecommand{\Pmax}{P_{\mathrm{max}}}  % Maximum momentum
\providecommand{\Mplbar}{\bar{M}_{\mathrm{Pl}}} % Reduced Planck mass
\providecommand{\eff}{\mathrm{eff}}       % Subscript for effective quantities
\providecommand{\Cosmo}{\ensuremath{\Lambda}}                % Cosmological constant (use \Cosmo in source)
\providecommand{\LCDM}{\ensuremath{\Cosmo\mathrm{CDM}}}  % \LambdaCDM (typesets with \Cosmo)
\providecommand{\Lcut}{\ensuremath{\Lambda_{\mathrm{cut}}}}  % Cutoff scale (brace when used as script)
\providecommand{\Rcs}{\ensuremath{R_{\mathrm{cs}}}}      % Configuration-space radius
\providecommand{\CosmoEff}{\ensuremath{\Cosmo_{\eff}}}   % Effective cosmological constant

% Disambiguation: \Rmom is the momentum-space (S^3) radius. 
% Reserve plain R for the 4D curvature scalar in f(R).
\providecommand{\Rmom}{\ensuremath{R_{\mathrm{mom}}}} 
\providecommand{\Punit}{P_{0}} % Momentum unit set by \Rmom (= \hbar/\Rmom)
\providecommand{\Rvec}{\ensuremath{R_{\mathrm{vec}}}} % invariant vector-space radius in the reduced round sector
% Specific numeric values and expressions
\providecommand{\Vlogcap}{L_{\mathrm{Hopf}}^{(1)}}
\providecommand{\Hopf}{\operatorname{Hopf}}
% Scalar phase (S^1) sector + universal cover notation (Paper I-A)
\providecommand{\tper}{T_\star}      % Time period/scale
\providecommand{\Eper}{E_\star}      % Energy period/scale
\providecommand{\ttilde}{\tilde t}   % Universal cover time
\providecommand{\Etilde}{\tilde E}   % Universal cover energy

% ======================================================================
% 10A. SERIES-WIDE LOGICAL / NOTATIONAL DISAMBIGUATIONS
% ======================================================================
% Use these symbols to avoid collisions among: (i) the Paper--III
% no-insertion/geometric-closure class, (ii) Paper--I-A trimer diagonals,
% (iii) Paper--V-A/V-C family-space diagonals, (iv) scalar-index running,
% and (v) protected shell-clock labels.
\providecommand{\Rprim}{\ensuremath{\mathsf R_{\mathrm{prim}}^{\Sigma}}} % primitive finite structural readouts
\providecommand{\Rloop}{\ensuremath{\mathsf R_{\circ}^{\Sigma}}} % closed-loop readout role
\providecommand{\Rbdry}{\ensuremath{\mathsf R_{\partial}^{\Sigma}}} % boundary-conditioned readout role
\providecommand{\Rgen}{\ensuremath{\mathsf R_{\mathrm{fin}}^{\Sigma}}} % finite generated closure of primitive roles
\providecommand{\PhysOcc}{\operatorname{PhysOcc}}
\providecommand{\Actual}{\operatorname{Actual}}
\providecommand{\GeomLic}{\operatorname{GeomLic}}
\providecommand{\InvDist}{\operatorname{InvDist}}
\providecommand{\Dcl}{\ensuremath{\mathsf{D}_1}}      % unmarked positive Dirac-even closed spectral-readout form class
\providecommand{\Dheat}{\ensuremath{\mathsf D_{\mathrm{heat}}}} % one-parameter heat-semigroup subclass
\providecommand{\DtriI}{\ensuremath{D^{\mathrm{tri}}_1}} % Paper--I-A cyclic-trimer diagonal
\providecommand{\DtriII}{\ensuremath{D^{\mathrm{tri}}_2}}
\providecommand{\DfamI}{\ensuremath{D^{\mathrm{fam}}_1}} % Family-space diagonal selector
\providecommand{\DfamII}{\ensuremath{D^{\mathrm{fam}}_2}}
\providecommand{\alpharun}{\ensuremath{\alpha_{\mathrm{run}}}} % scalar spectral-index running
\providecommand{\Ksh}{\ensuremath{K_{\mathrm{sh}}}} % protected shell-clock coordinate
\providecommand{\ddef}{\ensuremath{\delta_{\mathrm{def}}}} % universal defect scalar after Paper--III realization

% ======================================================================
% 10B. SERIES-WIDE LOGICAL STATUS TAGS
% ======================================================================
% Upright small-caps wrapper for audit/status labels.  It prevents LaTeX from
% trying to use unavailable small-caps italic font shapes inside theorem bodies
% such as TU/lmr/m/scit.
\providecommand{\AuditLabel}[1]{}
\renewcommand{\AuditLabel}[1]{{\normalfont\scshape #1}}

% These tags are bookkeeping labels for consistency audits. They do not
% create new mathematical hypotheses; they record how a datum is being used
% in the proof chain.
\providecommand{\StatusPrimitive}{}
\renewcommand{\StatusPrimitive}{\AuditLabel{Primitive}}
\providecommand{\StatusDerived}{}
\renewcommand{\StatusDerived}{\AuditLabel{Derived}}
\providecommand{\StatusImported}{}
\renewcommand{\StatusImported}{\AuditLabel{Imported}}
\providecommand{\StatusConditional}{}
\renewcommand{\StatusConditional}{\AuditLabel{Conditional}}
\providecommand{\StatusBridge}{}
\renewcommand{\StatusBridge}{\AuditLabel{Bridge-derived}}
\providecommand{\StatusComparison}{}
\renewcommand{\StatusComparison}{\AuditLabel{Comparison-only}}

% ======================================================================
% 10C. SERIES-WIDE CLAIM-WORDING AUDIT LABELS
% ======================================================================
% These are textual audit labels used to keep modal language consistent
% across the series.  They introduce no new mathematical assumptions.
\providecommand{\ClaimForced}{}
\renewcommand{\ClaimForced}{\AuditLabel{forced}}
\providecommand{\ClaimUnique}{}
\renewcommand{\ClaimUnique}{\AuditLabel{unique}}
\providecommand{\ClaimCanonical}{}
\renewcommand{\ClaimCanonical}{\AuditLabel{canonical}}
\providecommand{\ClaimMinimal}{}
\renewcommand{\ClaimMinimal}{\AuditLabel{minimal}}
\providecommand{\ClaimClosed}{}
\renewcommand{\ClaimClosed}{\AuditLabel{closed}}
\providecommand{\ClaimPrediction}{}
\renewcommand{\ClaimPrediction}{\AuditLabel{prediction}}

% ======================================================================
% 10D. SERIES-WIDE THEOREM-CHAIN AUDIT LABELS
% ======================================================================
% These labels mark the direction of load-bearing theorem chains. They are
% audit annotations only; they introduce no new mathematical assumptions.
\providecommand{\ChainInput}{}
\renewcommand{\ChainInput}{\AuditLabel{Chain input}}
\providecommand{\ChainOutput}{}
\renewcommand{\ChainOutput}{\AuditLabel{Chain output}}
\providecommand{\ChainGuard}{}
\renewcommand{\ChainGuard}{\AuditLabel{Chain guard}}
\providecommand{\ChainNoBackflow}{}
\renewcommand{\ChainNoBackflow}{\AuditLabel{No backflow}}

% ======================================================================
% 10E. SERIES-WIDE LOCAL THEOREM-AUDIT LABELS
% ======================================================================
% These labels mark a statement-level dependency audit around individual
% theorem/definition nodes. They are bookkeeping labels only.
\providecommand{\LocalAudit}{}
\renewcommand{\LocalAudit}{\AuditLabel{Local audit}}
\providecommand{\UsesOnly}{}
\renewcommand{\UsesOnly}{\AuditLabel{Uses only}}
\providecommand{\Exports}{}
\renewcommand{\Exports}{\AuditLabel{Exports}}
\providecommand{\DoesNotUse}{}
\renewcommand{\DoesNotUse}{\AuditLabel{Does not use}}

% ======================================================================
% 10F. SERIES-WIDE NUMERICAL-ROW AUDIT LABELS
% ======================================================================
% These labels classify individual rows in numerical, benchmark, status,
% and comparison tables.  They are audit annotations only and do not change
% any theorem, value, formula, or comparison number.
\providecommand{\RowAudit}{}
\renewcommand{\RowAudit}{\AuditLabel{Numerical-row audit}}
\providecommand{\RowTheorem}{}
\renewcommand{\RowTheorem}{\AuditLabel{Theorem-fixed row}}
\providecommand{\RowBridge}{}
\renewcommand{\RowBridge}{\AuditLabel{Bridge-fixed row}}
\providecommand{\RowRepresentative}{}
\renewcommand{\RowRepresentative}{\AuditLabel{Representative row}}
\providecommand{\RowComparison}{}
\renewcommand{\RowComparison}{\AuditLabel{Comparison row}}
\providecommand{\RowCalibrated}{}
\renewcommand{\RowCalibrated}{\AuditLabel{Calibrated row}}
\providecommand{\RowEstimate}{}
\renewcommand{\RowEstimate}{\AuditLabel{Estimate row}}
\providecommand{\RowOpen}{}
\renewcommand{\RowOpen}{\AuditLabel{Open row}}
\providecommand{\RowNoSelector}{}
\renewcommand{\RowNoSelector}{\AuditLabel{No selector}}

% ======================================================================
% 10G. TABLE HELPERS
% ======================================================================
% Table helpers used by Paper V-C.  The unused readout-label preset has been
% removed; only active table helpers remain.
\providecommand{\gpsection}[1]{%
  \addlinespace[0.55em]
  \multicolumn{3}{@{}l}{\textsc{#1}}\\[-0.25em]
  \midrule
}
\providecommand{\mathcell}[1]{%
  \(\begin{gathered}#1\end{gathered}\)%
}

\newcommand{\DerBy}[1]{%
  \begin{tabular}[t]{@{}l@{}}#1\end{tabular}%
}

\newcommand{\TableNotes}[1]{%
  \par\medskip
  \begin{minipage}{0.94\linewidth}
  \footnotesize\emph{Notes.} #1
  \end{minipage}\par
}
% ======================================================================
% 10H. PAPER VI-A / QCD NOTATION
% ======================================================================
% Shared symbols used by Paper VI-A and downstream QCD/Yang--Mills handoffs.
\providecommand{\Pqcd}{\mathfrak P_{\rm VIII}^{\rm QCD}}
\providecommand{\Iqcd}{\mathfrak I_{\rm QCD}^{(7)}}
\providecommand{\muQCD}{\mu_{\rm QCD}^{\rm int}}
\providecommand{\LambdaQCD}{\Lambda_{\rm QCD}}
\providecommand{\LambdaBH}{\Lambda_{\rm BH}^{\rm QCD}}
\providecommand{\HQCD}{\mathcal H_{\rm QCD}^{\rm BH}}
\providecommand{\HQCDfin}{H_{\rm QCD}^{\rm fin}}
\providecommand{\Hraw}{\mathcal H_{\rm raw}}
\providecommand{\HphysQCD}{\mathcal H_{\rm phys}^{\rm QCD}}
\providecommand{\Vcol}{V_C}
\providecommand{\Ps}{P_{\rm sing}}
\providecommand{\PiOct}{\Pi_8}
\providecommand{\MSbar}{\overline{\rm MS}}
\providecommand{\YM}{\mathrm{YM}}
\providecommand{\Ktri}{K_{\triangle}}
\providecommand{\KYM}{K_{\YM}^{\rm BH}}
\providecommand{\aBH}{a_{\rm BH}}
\providecommand{\epsBH}{\varepsilon_{\rm BH}}
\newcommand{\RowConstructive}{\textnormal{\textsc{Constructive}}}
\newcommand{\RowObligation}{\textnormal{\textsc{Bridge obligation}}}
\newcommand{\RowTarget}{\textnormal{\textsc{Target row}}}
\newcommand{\RowValidation}{\textnormal{\textsc{Validation row}}}
\providecommand{\slashed}[1]{\not{#1}}

% ======================================================================
% 10I. PAPER VI-B / CONTINUOUS-QFT NOTATION
% ======================================================================
% Shared symbols used by Paper VI-B and any downstream continuous-QFT or
% Yang--Mills handoff paper.  All definitions are guarded so individual
% papers may override them locally before loading this file if needed.
\providecommand{\A}{\mathfrak A}
\providecommand{\Hphys}{\mathcal H_{\rm phys}}
\providecommand{\QCD}{\mathrm{QCD}}
\providecommand{\BH}{\mathrm{BH}}
\providecommand{\loc}{\mathrm{loc}}

\providecommand{\RBH}{\mathcal R_{\rm BH}}
\providecommand{\DG}{\mathfrak D_G}
\providecommand{\PW}{\mathcal{P\!W}} % Peter--Weyl spectral-symbol subspace
\providecommand{\Toe}{\operatorname{Toe}} % Toeplitz compression/representative
\providecommand{\Hol}{\operatorname{Hol}} % Holonomy

% ======================================================================
% 11. LATE-LOADED PACKAGES (Must be at the end of the preamble)
% ======================================================================
% Load late packages only if a paper has not already loaded them, avoiding
% option clashes in paper-specific preambles.
\makeatletter
\@ifpackageloaded{hyperref}{}{\usepackage{hyperref}}
\makeatother
% --- Series-wide disambiguations (Pass 3) ---
\providecommand{\diracdelta}{\delta_{\rm D}} % Dirac delta distribution (delta reserved for hierarchy)
\providecommand{\Hg}{H_{\mathrm g}} % Hodge harmonic projector (keep H for Hubble)
\providecommand{\Gr}{G_{\mathrm r}} % Green operator (keep G for Newton)
\providecommand{\FD}{\mathrm{D}} % functional/variational derivative symbol (delta reserved)
\providecommand{\Id}{\mathbf{1}} % identity/Kronecker tensor (avoid \delta_{ij})


% BEGIN INLINE SOURCE: Paper00-ActiveExternalReferences-v2_130.tex
% Framework v2.191: current reference binding; existing definitions retained.
% Framework v2.190: current22-root registry with native clock-evolution proofs.
% Framework v2.189: exact current22-root external registry, including retained v2.188 hardening owners.
% Framework v2.187: current registry synchronization only; source release v2.186 preserved.
% Framework v2.186: authoritative current 22-paper registry.
% Framework v2.182: all current 22 logical paper AUX bindings.
% Framework v2.170: current I-B v1.43 external AUX binding.
% Framework v2.167: exact current 22-root registry after CP--BL2 integration.
% Framework v2.166: exact current 22-root compact-release AUX registry.
% Framework v2.165: authoritative current root bindings.
% Framework v2.164: exact current 22 manuscript AUX bindings.
% Framework v2.162: finite metric, functional Hilbert and scoped Ward integration.
% Framework v2.161: current 22 manuscript AUX bindings; source only.
% Framework v2.156 / registry v2.113: synchronized current 22 manuscript roots; source-only registry.
% Framework v2.143 / registry v2.111: active manuscript AUX basenames; no generated AUX or PDF included.
% ======================================================================
% Central external-reference imports, numbers only, no PDF jumping
% ======================================================================

% xr-hyper is loaded centrally so that papers may retain their existing
% \usepackage{xr-hyper} lines without redefining \externaldocument after this
% registry is installed.  If plain xr was loaded earlier, retain it.
\makeatletter
\@ifpackageloaded{xr-hyper}{}{%
  \@ifpackageloaded{xr}{}{\usepackage{xr-hyper}}%
}
\makeatother

% Make external-document registration idempotent.  The registry identifies a
% logical paper by its Paper<number><optional letter> prefix.  Prefix-qualified
% imports remain independent.  Current paper roots contain no local
% \externaldocument calls; this guard also makes accidental duplicates harmless.
\ExplSyntaxOn
\prop_new:N \g_bh_external_refs_prop
\cs_new_eq:NN \BHOriginalExternalDocument \externaldocument

\cs_new_protected:Npn \bh_external_document:nnnn #1#2#3#4
 {
  \tl_set:Nn \l_tmpa_tl {#3}
  \regex_replace_once:nnN
   { \A (?: xref/ )? (Paper[0-9]+[A-G]?) .* \Z }
   { \1 }
   \l_tmpa_tl
  \tl_set:Nx \l_tmpb_tl { #1 | \tl_use:N \l_tmpa_tl }
  \prop_if_in:NVTF \g_bh_external_refs_prop \l_tmpb_tl
   {
    \PackageInfo{BornHopf}
     {Skipping~duplicate~external~refs~for~#3}
   }
   {
    \prop_gput:NVn \g_bh_external_refs_prop \l_tmpb_tl {#3}
    \IfNoValueTF{#4}
     { \BHOriginalExternalDocument[#1][#2]{#3} }
     { \BHOriginalExternalDocument[#1][#2]{#3}[#4] }
   }
 }

\RenewDocumentCommand{\externaldocument}{O{} O{} m o}
 { \bh_external_document:nnnn {#1}{#2}{#3}{#4} }

\NewDocumentCommand{\BHExternalRefs}{m}
 {
  % Do not import the aux file of the logical paper currently being compiled,
  % even when the job and registry use different version suffixes.
  \tl_set:Nx \l_tmpa_tl { \c_sys_jobname_str }
  \regex_replace_once:nnN
   { \A (Paper[0-9]+[A-G]?) .* \Z }
   { \1 }
   \l_tmpa_tl
  \tl_set:Nn \l_tmpb_tl { #1 }
  \regex_replace_once:nnN
   { \A (Paper[0-9]+[A-G]?) .* \Z }
   { \1 }
   \l_tmpb_tl
  \str_if_eq:eeTF { \tl_use:N \l_tmpa_tl } { \tl_use:N \l_tmpb_tl }
   {
    \PackageInfo{BornHopf}
     {Skipping~external~refs~for~#1~because~this~is~the~current~document}
   }
   {
    \IfFileExists{xref/#1.aux}
     { \externaldocument[][nocite]{xref/#1}[#1.pdf] }
     {
      \IfFileExists{#1.aux}
       { \externaldocument[][nocite]{#1} }
       {
        \PackageInfo{BornHopf}
         {External~aux~file~#1.aux~not~found;~no~labels~imported}
       }
     }
   }
 }
\ExplSyntaxOff

% External references are automatic.  This file is the sole owner of both the
% bibliography and the current Paper-I-A--VII-D AUX registry (30 roots including I-C).  Missing auxiliaries
% are informational on the first pass and are populated by the source-only
% two-pass series build.  Basenames deliberately omit temporary download/copy
% suffixes such as ``(1)'' or ``-1''.
  % Paper I-A: round S3 compendium.
  \BHExternalRefs{Paper1A-Geometry-v5_180}

  % Paper I-B: complete exceptional geometry and finite spectral structures.
  \BHExternalRefs{Paper1B-ExceptionalGeometry-v1_22}
  % Paper I-D: technical companion of I-A through I-C.
  \BHExternalRefs{Paper1D-TechnicalMaterial-v1_69}
  % Paper I-C: reciprocal M8 compendium.
  \BHExternalRefs{Paper1C-M8Geometry-v1_144}

\ifdefined\BHMathOnly
  % Papers I-A and I-B import only the mathematical atlas labels.
\else
  % Paper II: five physical-mathematical main papers plus one technical proof volume.
  \BHExternalRefs{Paper2A-GeometricConstitution-v3_227}
  \BHExternalRefs{Paper2B-MomentumCarrier-v2_147}
  \BHExternalRefs{Paper2C-ReciprocalLocalization-v2_174}
  \BHExternalRefs{Paper2D-SymmetrySpectra-v2_145}
  % The former II-F export labels remain in II-E; current II-F is the technical volume.
  \BHExternalRefs{Paper2E-KinematicExports-v3_139}
  % Paper II-F: complete technical proof and calculation owner for II-A--II-E.
  \BHExternalRefs{Paper2F-TechnicalMaterial-v1_70}

  % Paper III: Dynamics and Intrinsic Geometric Constitution (six components).
  \BHExternalRefs{Paper3A-GeneratedDynamics-v1_180}
  \BHExternalRefs{Paper3B-SpectralAction-v1_191}
  \BHExternalRefs{Paper3C-ProtectedHierarchy-v1_193}
  \BHExternalRefs{Paper3D-RunningResponse-v1_167}
  \BHExternalRefs{Paper3E-GravityAction-v1_203}
  \BHExternalRefs{Paper3F-CompleteDynamics-v1_276}
  % Paper III-G: complete preserved technical proof owner.
  \BHExternalRefs{Paper3G-TechnicalMaterial-v1_57}

  % Paper IV-A: generated finite quantum mechanics.
  \BHExternalRefs{Paper4A-FiniteQuantumMechanics-v4_216}
  % Paper IV-B: complete finite quantum processes and physical response.
  \BHExternalRefs{Paper4B-QuantumResponse-v1_54}
  % Paper IV-C: extended proofs and finite technical calculations.
  \BHExternalRefs{Paper4C-TechnicalMaterial-v1_50}

  % Paper V-A: The Standard Model.
  \BHExternalRefs{Paper5A-FiniteMatterGeometry-v3_218}

  % Paper V-B: Finite Spectrum and Flavor Closure.
  \BHExternalRefs{Paper5B-TerminalSpectrum-v5_202}

  % Paper V-C: Finite Interaction Boundaries of the Standard Model.
  \BHExternalRefs{Paper5C-InteractionResponse-v2_182}
  % Paper V-D: preserved technical proof and calculation owner.
  \BHExternalRefs{Paper5D-TechnicalMaterial-v1_58}

  % Paper VI-A: Gravity and Closed Cosmology.
  \BHExternalRefs{Paper6A-GravityCosmology-v2_267}

  % Paper VI-B: Geometric Dark Sector and Cosmological Readouts.
  \BHExternalRefs{Paper6B-CosmologyReadouts-v1_214}

  % Paper VI-C: complete technical proof and comparison owner.
  \BHExternalRefs{Paper6C-TechnicalMaterial-v1_53}

  % Paper VII-A: Finite QCD Kernel and Nonperturbative Closure.
\BHExternalRefs{Paper7A-QCDKernel-v1_180}

  % Paper VII-B: Hadronic, Thermal, and Nuclear QCD Readouts.
  \BHExternalRefs{Paper7B-QCDReadouts-v1_165}

  % Paper VII-C: Continuous Quantum Field Theory.
  \BHExternalRefs{Paper7C-ContinuousQFT-v2_208}

  % Paper VII-D: complete technical material for VII-A through VII-C.
  \BHExternalRefs{Paper7D-TechnicalMaterial-v1_60}

\fi

% END INLINE SOURCE: Paper00-ActiveExternalReferences-v2_130.tex


\hypersetup{hidelinks}


\usetikzlibrary{positioning, calc} 

\usepackage{xr-hyper}

\title{\textbf{Born--Hopf Geometry II-D: Symmetries and Spectral Kinematics}\\[0.5em]
\large Casimirs, Hopf geometry, higher spin, and finite recovery}

\author{Peter Kugelmann\thanks{\texttt{pkugelmann@proton.me}}}
\date{Version 2.145 --- 20 November 2026 \\ Framework v3.063}

\setlength{\emergencystretch}{3em}


\begin{document}

\maketitle

\begin{abstract}
Paper~II-D develops the spectral and algebraic-symmetry consequences of
Born--Hopf's constituted round $SU(2)$ parent and finite reciprocal
kinematics.  The structural momentum role is the upstream
$\mathcal H_p\cong\mathcal H_1$ with $m_p=1$ and
$\dim\mathcal H_1=15$; larger $\mathcal H_m$ are auxiliary finite
comparison families, not new momentum-selection stages.  The central
$\Spin(4)$ grading separates an eleven-dimensional geometric sector
from a four-dimensional spinorial sector.  Their even and odd
endomorphism spaces have complex dimensions $137$ and $88$;
no nonzero invariant odd bridge is generated by the even operator
calculus.

Smooth Laplace, heat, zeta, Hodge and higher-spin formulas are
statements about the round parent under their stated operator hypotheses.
Scalar fixed-band compression has an explicit comparison map and zero
error after admission of the band; guarded nonlinear words and the
local round-metric defect have distinct, explicit conditions.  In
particular these are not a large-momentum-level limit.  The canonical
retained Levi--Civita Hodge/spin-square result is presently proved
in Paper~I-B, while the general round $S^3$ mathematical atlas is
owned by Paper~I-A.  Same-principal-symbol operators need not have
the canonical lower-order spectrum or ground.

Hopf connections, character charges, parity and time reversal give
geometric transport and fixed-locus results, not an independently
selected physical axis, charge, energy or dynamics.  On marked finite
exceptional comparison packets we derive conditional Petersen sectors,
exact order-ten Hopf Fourier sampling, charge-residue rank and
cubature certificates, and a finite correlated-memory comparison.
The latter is not a particle-decay clock.  All physical preparation,
operator identification, occurrence and calibration interfaces retain
their separate upstream or downstream obligations.  Extended spectral,
selector and higher-spin calculations are contained in the Paper~II-F
technical companion.
\end{abstract}

\tableofcontents
\listoftables
\newpage


\section{Constituted inputs and spectral domains}
\label{paper2d:sec:organized-scope}

\subsection{Introduction}
\label{p2d:sec:introduction}
\paragraph{Geometric involutions and occurrence.}
Associated-character, Hopf-normalization, spectral, and \(q\leftrightarrow-q\)
outputs retain the typed graphs that generate them.  An explicitly generated
involution is an invertible bridge on its declared orbit; it is not thereby
physical reversal of an elapsed-time process.  Spectral inclusions,
projections, and charge-pair graphs retain their own bridge types, while a
scale, dispersion law, state, or dynamics token remains separately rooted.

\paragraph{Occurrence condition.}
All physical interfaces below are role-specific specializations of
Paper~II-A's stage-safe G0 certificate.  An asserted occurrence must exhibit a
contingent token of a geometrically licensed, claim-relative invariant
distinction with result-independent provenance; its tag alone supplies no
witness.  No instance reopens frozen upstream geometry, and a new role enters
only through a declared downstream dependency frontier.  Analytic and recovery
hypotheses remain separate.  The imported parent and carrier are generated
descendants of the upstream G0 operation occurrence; this paper adds no new
momentum token.


\paragraph{Full technical derivations.}
The active auxiliary classifications, exact calculations, and negative
results supporting this paper are collected without abridgment in
Paper~II-F, Section~\ref{p2ftech:sec:spectra}.
That companion is a proof owner, not an additional constructor,
physical occurrence, state, or calibration assumption.

Paper~II-D uses the derived compact parent and canonical Killing metric from
Paper~II-A, the finite equivariant envelope from Paper~II-B, and the operator
conventions from Paper~II-C.  Its native task is to collect consequences that
depend on those established structures but do not participate in their
derivation.  This ordering prevents spectral or topological data from being
used retroactively to select the carrier or its global form.

\subsubsection{Logical scope and ownership}
Spectral formulas and topological sectors are proved or imported with their
stated round-sphere hypotheses.  The geometric parity and time-reversal
operations below are stated only at the level justified by the compact
kinematic algebra.  Gravity-sector use of these results is exported through
Paper~II-E, and internal commutants are handled in Paper~II-F.  Physical threshold interpretation and the discharge of a free-dispersion law
belong to a later dynamical or phenomenological theorem.  Such a theorem is
used below only as a typed antecedent for oscillator formulas; it is not a
premise of any Paper-II kinematic theorem.

\paragraph{Constitutional inputs.}
Paper~II-A derives the unique least noncommutative parent-role skeleton
\(SU(2)\cong S^3\) and its round negative-Killing geometry, while GP2 and GP3
are respectively the Exact/Faithful and Irreducible/Least stages of the
constitutional calculus.  Paper~II-D needs only the resulting parent formulas;
it does not strengthen Paper~II-A's separately stated lift-stability boundary.
No center-odd test, FD, LR, or OR1 premise is introduced here.  The rank-one
statement remains role-local and does not exclude later \(T^2\) or higher-rank
reciprocal products.

\begin{remark}[Inherited geometry and operator domains]
\label{p2:rem:iid-dependency-record}
Paper~II-D consumes the current upstream chain rather than reopening the
momentum constitution: Paper~II-A supplies the round parent, generated
role-level co-realization, joint-family exhaustion, and recovery interfaces;
Paper~II-B supplies
\(\mathcal H_p\cong\mathcal H_1\) and \(m_p=1\); Paper~II-C supplies
the canonical retained coordinate/symbol branches, finite operator packages,
and the unmarked level-one completion
quotient
\[
 \mathfrak M_{\rm coord}(\mathcal H_1)/{\sim}
 \cong \mathbb R_{\ge0}^{3},
 \qquad
 [\widehat X]\longleftrightarrow
 (r_{1,0},r_{0,1},\varrho_{\rm sp}).
\]
Standard elliptic and higher-spin spectral results are used only under their
stated operator and normalization hypotheses.  Auxiliary TWR/SPEX, lower-ideal and $\mathcal E_9$ comparison
conditions are not independent GP2 momentum-constructor premises;
see Paper~II-F, Section~\ref{p2d:app:legacy-selector-diagnostics},
for the separate mathematical comparison domains.
\end{remark}

\paragraph{Mathematical claims and physical occurrence.}
The objects used below are kept in four logically distinct classes.  Derived
mathematical/process structure includes functional-calculus laws such as
\(e^{-tL}\), the central grading, and exact scalar fixed-band recovery.  G0
occurrences include an actually performed readout, an actual sharp spectral
frontier, or an actual reciprocal-position frontier history.  Unresolved
bridges include nonlinear or non-scalar parent/finite comparison maps together
with their domains, topologies, and remainders.  Downstream interpretation or
calibration includes a physical Hamiltonian, energy or threshold assignment,
field-speed identification, and numerical scale calibration.  Mathematical
availability in one class does not constitute occurrence or select a member of
another. The positive-comparison, kinetic-pair and lower-order-cone certificates in Section~\ref{p2d:sec:v2210-lit-owner-transfers} now discharge their stated transfer steps after the actual maps/forms/solder are supplied; the remaining obligation is that supplied application and its quantitative error or kinetic bound.


\subsubsection{Parent and finite-level operator domains}
\label{p2d:subsec:v2-recovery-contract}
The smooth round \(S^3\) spectra below are parent-grammar statements.  They
remain indispensable because differential order, locality, and principal
symbol cannot be recovered from one static finite spectrum alone.  The
structural momentum-role carrier has already been fixed upstream as
\(\mathcal H_p\cong\mathcal H_1\), with \(m_p=1\).  By contrast, generic
\(\mathcal H_m\) and scalar truncations are auxiliary finite families used for
representation bookkeeping, exact finite formulas, recovery statements, or
later roles whose frozen signatures explicitly constitute them.  They are not
automatic GP2-generated momentum stages.  Hence the level-one momentum-role
specialization sets \(m=1\) directly; it is not a large-\(m\) approximation.

For a reciprocal position role, a recovery claim is typed by
\[
 (B,\mathcal O_{\rm par},\mathcal O_m,J_{m,B},
   \mathcal T,r_{B,\mathcal O}),
\]
where \(B\) is the parent band, \(J_{m,B}\) is the comparison embedding,
\(\mathcal T\) is the stated topology, and \(r_{B,\mathcal O}\) is the
remainder in that topology.  With \(P_B\) the parent band projector and
\(Q_{x,m}\) the parent-to-finite compression, the claim must be either
\[
 J_{x,m,B}^*Q_{x,m}(\mathcal O_{\rm par})J_{x,m,B}
 =P_B\mathcal O_{\rm par}P_B
\]
above a finite threshold, or an explicit norm estimate with a remainder
\(r_{B,\mathcal O}(m)\to0\).  Peter--Weyl density by itself proves neither
estimate.  For the scalar spectral multipliers proved below,
\(\mathcal T\) is operator norm on the fixed finite band and the remainder is
exactly zero once \(m\ge K\).
Likewise, spectral resolution and local flatness are independent: large
\(m_x\) may recover smooth round \(S^3_x\) behavior on a band without proving
that a region of fixed size is approximately \(\mathbb R^3\).  On the exact
round branch the basic metric defect is derived explicitly below; nonround,
higher-derivative, and operator-flatness claims retain their own geometric
estimates.

\subsubsection{Spectral Functions}
\label{p2:sec:spectral-functions}

\paragraph{Finite-level family.}
Throughout the spectral sections, \(\mathcal H_1\) denotes the derived
structural momentum-role carrier.  A generic \(\mathcal H_m\) denotes only the
auxiliary/later-role finite module family unless a separate frozen signature
constitutes that level physically.

\begin{theorem}[Rigidity of the finite heat-process representative]
\label{p2d:thm:heat-semigroup-rigidity}
Let \(L=L^*\ge0\) act on a finite-dimensional Hilbert space and let
\((V_t)_{t\ge0}\) be a strongly continuous semigroup of self-adjoint positive
contractions with \(V_0=1\) and generator
\[
 L=-\left.\frac{d}{dt}V_t\right|_{t=0}.
\]
Then
\[
 \boxed{V_t=e^{-tL}.}
\]
If the occurring readout is precisely the canonical duality closure of
this process endomorphism, with no retained insertion or additional process,
its scalar representative is
\[
 \boxed{Z_L(t)=\operatorname{Tr}e^{-tL}.}
\]
The first box is a process-law theorem; the second additionally consumes the
closed-loop readout theorem of Paper~II-A.  Neither statement asserts that
this process is the physical action readout.
\end{theorem}

\begin{proof}
In finite dimension strong continuity is norm continuity.  The semigroup law
gives
\[
 \frac{V_{t+h}-V_t}{h}
 =V_t\frac{V_h-1}{h}\longrightarrow -V_tL.
\]
Differentiating on the other side gives
\(\dot V_t=-LV_t\), so \(L\) commutes with \(V_t\).  The initial-value
problem
\(\dot V_t=-LV_t,\ V_0=1\) has the unique solution \(e^{-tL}\).  Equivalently,
after diagonalizing \(L\), every scalar component satisfies the continuous
Cauchy exponential equation with derivative \(-\lambda\) at zero and hence
equals \(e^{-t\lambda}\). When the occurring closed-readout word is
the canonical duality closure of this resulting endomorphism,
Lemma~\ref{p2:lem:canonical-finite-duality-trace} and
Theorem~\ref{p2:thm:unmarked-natural-closure} give its ordinary trace.
Any retained insertion or additional process changes the specified word.
\end{proof}

\begin{corollary}[Atomic heat subclass]
\label{p2d:cor:atomic-heat-readout}
If the occurring physical closed readout is explicitly the canonical
duality closure of the same resolution semigroup evaluated once, and \(t=\Lambda^{-2}\), then
\[
 \boxed{S_{\rm heat}(L)=\operatorname{Tr}e^{-L/\Lambda^2}.}
\]
Thus the one-parameter class \(\Dheat\) has no arbitrary profile beyond the
calibration of semigroup time. Its block coefficients are the dimensions,
by Corollary~\ref{p2a:cor:v233-loop-coefficients}, without a prepared-state
assumption. Mathematical existence of the heat flow alone
does not establish this occurrence equality.
\end{corollary}

\begin{remark}[Other exact process representatives]
Once the indicated process role and required domains are supplied, finite
functional calculus gives
\[
\begin{array}{c|c|c}
 \text{process role}&F(\lambda)&\text{closed scalar}\\ \hline
 \text{heat semigroup}&e^{-t\lambda}&\operatorname{Tr}e^{-tL}\\
 \text{unitary group}&e^{-it\lambda}&\operatorname{Tr}e^{-itH}\\
 \text{sharp frontier}&\mathbf1_{[0,\Lambda]}(\lambda)&
   \operatorname{Tr}\mathbf1_{[0,\Lambda]}(L)\\
 \text{resolvent response}&(\lambda-z)^{-1}&
   \operatorname{Tr}(L-z)^{-1}\\
 \text{positive determinant}&\log\lambda&
   \operatorname{Tr}\log L,\qquad L>0.
\end{array}
\]
The sharp row needs an actual spectral-frontier occurrence; the resolvent row
excludes poles; and the displayed logarithmic row needs strict positivity
(otherwise it is restricted to the kernel complement, with zero modes
separately recorded) and, on an infinite parent, a regularization.  These are distinct process laws
inside one closed-loop outer role, not additional primitive readout roles.
\end{remark}

\paragraph{Explicit heat trace}
\label{p2:sec:heat-trace}

The spectral data in this subsection are the Paper~I-A harmonic package for the
derived round \(S^3(\Rmom)\), with the finite sums obtained by imposing the
Paper~I-A \(\Spin(4)\)-equivariant module family \(\mathcal H_m\)
\cite{Kugelmann2026a}.  We recall them here only to fix the kinematic
mode-counting and regularization conventions used later.

On scalar diagonal harmonics \(j_L=j_R=k/2\), the Laplacian eigenvalues on the round \(S^3(\Rmom)\) are
\[
\lambda_k=\frac{k(k+2)}{\Rmom^2},
\qquad k=0,1,2,\dots,
\]
with multiplicity \((k+1)^2\).  On the declared full
\(\Spin(4)\)-equivariant module package, the operator used below is instead
the chosen symmetric Casimir representative, with pair labels
\[
\lambda_{j_L,j_R}=\frac{2}{\Rmom^2}\bigl[j_L(j_L+1)+j_R(j_R+1)\bigr],
\qquad
\mathrm{mult}(j_L,j_R)=(2j_L+1)(2j_R+1).
\]
For \(j_L\ne j_R\), this is not the scalar Laplace--Beltrami spectrum of
\(L^2(S^3)\).

\subparagraph{\texorpdfstring{Smooth parent $S^3$ (presentation grammar):}{Smooth parent S3 (presentation grammar):}}
The exact heat trace is
\[
\mathrm{Tr}\, e^{-t\Delta}\big|_{\text{classical}}
= \sum_{k=0}^{\infty} (k+1)^2 e^{-k(k+2)t/\Rmom^2}.
\]
Shifting the index \(n=k+1\) yields
\[
\mathrm{Tr}\, e^{-t\Delta}
= e^{t/\Rmom^2} \sum_{n=1}^{\infty} n^2 e^{-n^2 t/\Rmom^2}.
\]
This series is known exactly in terms of Jacobi theta functions, but the closed sum is standard on homogeneous spaces \cite{Camporesi1991}.

\begin{definition}[Truncated Degeneracy---Diagonal Sector]\leavevmode\par\noindent
\label{p2:def:truncated-degeneracy-diagonal-sector}
For the \emph{diagonal} sector of finite Toeplitz \(S^3_m\) (i.e.\ \(j_L=j_R\)), the scalar label is \(k=2j_L=2j_R\).  The truncation condition \(j_L+j_R\le m\) becomes \(k\le m\), so
\begin{equation}
d_m^{\mathrm{diag}}(k) =
\begin{cases}
(k+1)^2, & 0\le k\le m,\\
0, & k>m.
\end{cases}
\end{equation}
No parity restriction is present, because \(j_L=j_R\) may be integral or half-integral.

For the \emph{full} finite Toeplitz \(S^3_m\) including non-diagonal representations, states are labeled by pairs \((j_L,j_R)\) with \(j_L+j_R\le m\), eigenvalues
\[
\lambda_{j_L,j_R}=\frac{2}{\Rmom^2}\bigl[j_L(j_L+1)+j_R(j_R+1)\bigr],
\]
and degeneracy \((2j_L+1)(2j_R+1)\).
\end{definition}

\begin{proposition}[Heat trace asymptotics]
\label{p2:prop:heat-trace-asymptotics}
As \(t \to 0^+\), the short-time asymptotics are
\cite{minakshisundaram1949some,berline1992heat}
\[
\mathrm{Tr}\, e^{-t\Delta}
\sim \frac{2\pi^2 \Rmom^3}{(4\pi t)^{3/2}}
\left(1 + \frac{t}{\Rmom^2} + \frac{t^2}{2\Rmom^4}+O(t^3)\right),
\]
consistent with the Minakshisundaram--Pleijel expansion for a compact
\(3\)-manifold with
\[
\operatorname{Vol}=2\pi^2\Rmom^3,\qquad
\scal=\frac{6}{\Rmom^2},\qquad
K=\frac{1}{\Rmom^2}.
\]
Here \(\scal\) denotes scalar curvature and \(K\) denotes sectional curvature.
\end{proposition}

\begin{proof}
This is the standard compact heat-kernel expansion specialized to the Paper~I-A
round-sphere normalization:
\[
  \Vol(S^3(\Rmom))=2\pi^2\Rmom^3,
  \qquad
  \scal=\frac{6}{\Rmom^2}
\]
\cite{Kugelmann2026a,minakshisundaram1949some,berline1992heat}.  The first
correction follows from \(a_1=\frac16\int\scal\,dV\).  For the second displayed
correction, use the scalar Laplacian coefficient
\[
a_2=\frac1{360}\int
\bigl(5\scal^2-2|\operatorname{Ric}|^2+2|\operatorname{Riem}|^2\bigr)\,dV .
\]
On the round three-sphere of radius \(\Rmom\), the curvature invariants are
\[
  \scal=\frac{6}{\Rmom^2},
  \qquad
  |\operatorname{Ric}|^2=\frac{12}{\Rmom^4},
  \qquad
  |\operatorname{Riem}|^2=\frac{12}{\Rmom^4}.
\]
Hence the density in the preceding display is \(1/(2\Rmom^4)\), giving the
coefficient \(t^2/(2\Rmom^4)\) after factoring out the leading volume term.
\end{proof}


\subparagraph{\texorpdfstring{Finite Toeplitz $S^3_m$ (finite-dimensional):}{Finite Toeplitz S^3_m (finite-dimensional):}}
On the truncated space $\mathcal{H}_m$ with $j_L,j_R\in\frac12\ZZ_{\ge0}$
and $j_L + j_R \leq m$, the full \(\Spin(4)\)-module Casimir heat
representative becomes a \textbf{finite sum}:
\[
\boxed{K_m^{\rm full}(t)
= \sum_{\substack{j_L,j_R\in\frac12\ZZ_{\ge0} \\ j_L + j_R \leq m}} (2j_L+1)(2j_R+1) \, 
\exp\left(-\frac{2t}{\Rmom^2}\bigl[j_L(j_L+1) + j_R(j_R+1)\bigr]\right)}
\]
For non-diagonal representations \((j_L\neq j_R)\), the eigenvalue \(\lambda_{j_L,j_R}\) differs from the scalar diagonal formula \(k(k+2)/\Rmom^2\).  Restricting to the diagonal sector gives the scalar finite trace
\[
\mathrm{Tr}_{\mathrm{diag}}\, e^{-t\Delta}\big|_{S^3_m}
=
\sum_{k=0}^{m}(k+1)^2
\exp\!\left(-\frac{k(k+2)t}{\Rmom^2}\right).
\]

\begin{remark}[Typed mode-coupling coefficients]
Before modes are coupled, their homogeneous bundle types and a
\(\Spin(4)\)-equivariant pointwise contraction, equivalently an intertwiner
from the source tensor product to the chosen target harmonic bundle, must be
fixed.  For every such well-typed coupling, projection back to a finite
\(S^3\) harmonic basis reduces to the left--right Clebsch--Gordan and
recoupling coefficients of the Paper-I-A harmonic package.  The resulting
selection rules constrain both \(\SU(2)\) factors.  There is no
bundle-independent coupling coefficient before the contraction and target
bundle have been specified.
\end{remark}


\paragraph{Spectral zeta function}
\label{p2:sec:spectral-zeta}

\subparagraph{\texorpdfstring{Classical $S^3$:}{Classical S^3:}}
The spectral zeta function (excluding the zero mode) is
\[
\zeta_\Delta(s)\big|_{\text{classical}}
= \sum_{k=1}^{\infty} (k+1)^2 \left[\frac{k(k+2)}{\Rmom^2}\right]^{-s}
= \Rmom^{2s} \sum_{k=1}^{\infty} (k+1)^2 [k(k+2)]^{-s}.
\]

\subparagraph{\texorpdfstring{Finite Toeplitz $S^3_m$:}{Finite Toeplitz S^3_m:}}
On the truncated diagonal scalar sector, the spectral zeta function is the finite sum
\[
\boxed{
\zeta_{\Delta,\mathrm{diag}}^{(m)}(s)
= \Rmom^{2s}\sum_{k=1}^{m}(k+1)^2[k(k+2)]^{-s}.
}
\]
On the full \(\Spin(4)\)-equivariant module package, the corresponding finite spectral sum is
\[
\boxed{
\zeta_{\Delta,\mathrm{full}}^{(m)}(s)
=
\sum_{\substack{j_L,j_R\in\frac12\ZZ_{\ge0}\\j_L+j_R\le m\\(j_L,j_R)\ne(0,0)}}
(2j_L+1)(2j_R+1)
\left[
\frac{2}{\Rmom^2}
\bigl(j_L(j_L+1)+j_R(j_R+1)\bigr)
\right]^{-s}.
}
\]
Both are \textbf{entire functions} of \(s\) because the sums are finite.  The
classical poles occur at the usual heat-kernel locations.  In dimension \(3\),
these include
\[
  s=\frac32,\ \frac12,\ -\frac12,\ldots.
\]
They are artifacts of the infinite classical sum and disappear at fixed finite Toeplitz
truncation.

For the classical scalar zeta function, shifting \(n=k+1\) gives
\[
\zeta_\Delta(s)
= \Rmom^{2s}\sum_{n=2}^{\infty} n^2[(n-1)(n+1)]^{-s}
= \Rmom^{2s}\sum_{n=2}^{\infty} n^2(n^2-1)^{-s}.
\]

\begin{proposition}[Spectral zeta continuation]
\label{p2:prop:spectral-zeta-continuation}
The series converges absolutely for \(\Re(s) > 3/2\). It admits an analytic continuation to a meromorphic function on \(\CC\) with a simple pole at \(s = 3/2\) (residue determined by the volume) and possible further simple poles at \(s=3/2-j\), \(j=1,2,\dots\), whose residues are determined by heat-kernel coefficients and may vanish in special cases \cite{berline1992heat}.
\end{proposition}

\begin{proof}
The convergence statement follows from the Paper~I-A scalar spectrum \(\lambda_k\sim k^2/\Rmom^2\) with multiplicity \((k+1)^2\).  The continuation is the standard Mellin-transform consequence of the compact heat-kernel expansion recalled in Proposition~\ref{p2:prop:heat-trace-asymptotics} \cite{berline1992heat}.  Thus no new geometric derivation is being introduced in Paper~II.
\end{proof}


\paragraph{Scalar spectral finite parts and finite oscillator sums}
\label{p2:sec:casimir}

\begin{remark}[Spectral data and the dispersion interface]
\label{p2d:rem:dispersion-speed-boundary}
The results here determine kinematic spectra.
A physical spectral-energy interpretation additionally requires the typed
dispersion, identification of its spatial operator with the Laplacian in
question, and clock/calibration data.
The equality \(c_{\rm field}=c_{\rm op}\) is available only when the
specified field satisfies the principal-ancestry hypotheses of
Theorem~\ref{p2a:thm:undeformed-host-principal-inheritance}.
A finite-block theorem alone supplies no parent-continuum extension.
Paper~II-F, Section~\ref{p2d:app:oscillator-examples} records the conditional
finite sums and continuum finite parts, including these hypotheses and the
required regularization choices.
\end{remark}

\subparagraph{Deferred worked spectral-energy representatives.}
The continuum finite-part prescription and the exact finite oscillator sums
are supplementary conditional examples rather than inputs to the kinematic
chain.  Their hypotheses, regularization boundary, and explicit \(m=1\)
evaluations are given in Paper~II-F, Section~\ref{p2d:app:oscillator-examples}.

\begin{table}[htbp]
\centering
\scriptsize
\begingroup
\setlength{\tabcolsep}{2pt}
\renewcommand{\arraystretch}{1.18}

\begin{tabular}{|
  L{2.8cm}|
  L{4.6cm}|
  L{\dimexpr\linewidth-2.8cm-4.6cm-6\tabcolsep-4\arrayrulewidth\relax}|
}
\hline
\textbf{Quantity}
& \textbf{Classical \(S^3\) Formula}
& \textbf{Finite Toeplitz \(S^3_m\) Formula} \\
\hline

\multicolumn{3}{|c|}{\textbf{\textit{Spectral Functions}}} \\
\hline

Scalar heat trace
& \(\displaystyle
  \sum_{k=0}^{\infty}(k+1)^2
  e^{-k(k+2)t/\Rmom^2}\)
& \(\displaystyle
  \sum_{k=0}^{m}(k+1)^2
  e^{-k(k+2)t/\Rmom^2}\) \\
\hline

Full module finite heat trace
& Not a scalar \(L^2(S^3)\) heat trace
& \(\displaystyle
  \sum_{\substack{
    j_L,j_R\in\frac12\ZZ_{\ge0}\\
    j_L+j_R\leq m
  }}
  d_{j_L,j_R}
  e^{-\lambda_{j_L,j_R}t}\) \\
\hline

Spectral zeta
& \(\displaystyle
  \Rmom^{2s}
  \sum_{k=1}^{\infty}
  \frac{(k+1)^2}{[k(k+2)]^s}\)
& Finite spectral sum; it is entire in \(s\), not a polynomial claim \\
\hline

Mode count \(N(\lambda)\)
& \(\displaystyle
  \frac{\Rmom^3}{3}\lambda^{3/2}+O(\lambda)\)
& Exact with degeneracy:
  \[
    N_m(\lambda)
    =
    \sum_{\substack{
      j_L+j_R\le m\\
      \lambda_{j_L,j_R}\le\lambda
    }}
    (2j_L+1)(2j_R+1).
  \] \\
\hline

\end{tabular}

\caption[Parent and finite Toeplitz \(S^3\) spectral data]{Parent and finite Toeplitz \(S^3\) spectral data.}
\label{p2:tab:spectral-data}
\TableNotes{%
Here \(d_{j_L,j_R}=(2j_L+1)(2j_R+1)\) and \(\lambda_{j_L,j_R} = \frac{2}{\Rmom^2}
[j_L(j_L+1)+j_R(j_R+1)]\).  Conditional spectral-energy representatives and their dispersion,
operator-identification, and regularization requirements are stated in
Paper~II-F, Section~\ref{p2d:app:oscillator-examples}.
}
\endgroup
\end{table}

\begin{remark}[No inverse-spectral input]
No inverse-spectral rigidity theorem is used in the Paper-II derivation.  The
round metric is fixed by the upstream geometric construction before any
spectrum is evaluated; none of the spectral tables in this paper is used to
infer roundness or metric uniqueness from spectral data.
\end{remark}


\section{Parent-to-finite spectral recovery and error bounds}
\label{paper2d:sec:organized-recovery}

\subsection{Typed Parent-to-Finite Spectral Recovery}
\label{p2d:sec:typed-parent-finite-recovery}

For the scalar parent, let
\[
 \mathcal B_K
 :=\bigoplus_{0\le k\le K}V_{k/2}\otimes V_{k/2}^{*}
 \subset L^2(S^3),
 \qquad
 \mathcal H_m^{\rm sc}
 :=\bigoplus_{0\le k\le m}V_{k/2}\otimes V_{k/2}^{*},
\]
let \(P_K\) be the parent orthogonal projector, let
\(\Pi_m^{\rm sc}:L^2(S^3)\to\mathcal H_m^{\rm sc}\) be the scalar truncation
projector, and for \(m\ge K\) let
\(J_{m,K}:\mathcal B_K\hookrightarrow\mathcal H_m^{\rm sc}\) be the canonical
coefficient inclusion.  Define the parent-to-finite spectral compression on
the declared functional calculus by
\[
 Q_m^\Delta(F(\Delta))
 :=\Pi_m^{\rm sc}F(\Delta)\Pi_m^{\rm sc}
    \big|_{\mathcal H_m^{\rm sc}}.
\]
This is the scalar spectral component of the paired realization map
\(Q_{R,m}\) of Paper~II-A.

\begin{theorem}[Exact fixed-band recovery of scalar parent operators]
\label{p2d:thm:scalar-fixed-band-recovery}
For every finite \(K\), every \(m\ge K\), and every spectral multiplier
\(F(\Delta)\) defined on the parent spectrum at least through the finite level
\(m\),
\[
 J_{m,K}^{*}Q_m^\Delta(F(\Delta))J_{m,K}
 =P_KF(\Delta)P_K\big|_{\mathcal B_K}.
\]
This first identity is a statement about nested projections and uses no
spectral property of the compressed operator.  More generally, if a possibly
unbounded parent operator \(\mathcal O\) satisfies
\[
       \mathcal H_m^{\rm sc}\subseteq D(\mathcal O),
\]
then \(\mathcal B_K\subseteq\mathcal H^{\rm sc}_m\) gives
\(P_K\Pi^{\rm sc}_m=P_K\), and every composition below is defined on the
finite band, so
\[
 J^{*}_{m,K}\Pi^{\rm sc}_m\mathcal O\,\Pi^{\rm sc}_mJ_{m,K}
 =P_K\mathcal OP_K\big|_{\mathcal B_K}.
\]
The operator content lies in the next statement, which does use the hypothesis:
because \(F(\Delta)\) is diagonal in the Peter--Weyl decomposition it
\emph{preserves} \(\mathcal B_K\), so that after identifying \(\mathcal B_K\) with
its image in the finite auxiliary module, \(P_KF(\Delta)P_K=F(\Delta)\big|_{\mathcal B_K}\)
and the parent and finite operators agree exactly on that band.  The
restricted spectrum, heat factors, resolvents away from their poles, and every
other declared scalar spectral multiplier therefore stabilize at the finite
threshold \(m=K\).
\end{theorem}

\begin{proof}
The scalar Laplacian is diagonal in the Peter--Weyl decomposition and preserves
\(\mathcal B_K\).  For \(m\ge K\), \(\Pi_m^{\rm sc}\) is the identity on
that band.  Both sides consequently act on each retained eigenspace by the
same scalar \(F(\lambda_k)\).
\end{proof}

\begin{corollary}[Exact fixed-band recovery of process readouts]
\label{p2d:cor:fixed-band-process-readout-recovery}
Under the hypotheses of
Theorem~\ref{p2d:thm:scalar-fixed-band-recovery}, every declared scalar
process representative \(F(\Delta)\)---including heat factors, sharp
projectors, resolvents away from poles, and logarithms on their declared
domains---is recovered exactly on \(\mathcal B_K\) for \(m\ge K\).  If a
closed-loop or boundary-conditioned G0 readout occurrence is supplied, closing
or evaluating the recovered finite operator gives the same scalar as first
restricting the parent process to \(\mathcal B_K\) and then applying that
readout.
\end{corollary}

\begin{proof}
The recovery theorem gives equality of the two operators on the finite band.
The algebraic trace and every fixed boundary coefficient are functions of that
finite operator equality, so applying the same generated closure or the same
generated boundary data preserves equality.  This proves recovery of the
readout value without turning its G0 occurrence into a consequence of spectral
compression.
\end{proof}

\begin{corollary}[Structural momentum-role and conditional position applications]
\label{p2d:cor:physical-level-recovery-handoff}
The theorem has the following Paper-II applications.
\begin{enumerate}[label=(\roman*),leftmargin=2em]
\item Paper~II-B already gives the structural momentum-role seed
\(\mathcal H_p\cong\mathcal H_1\) and \(m_p=1\).  Therefore
Theorem~\ref{p2d:thm:scalar-fixed-band-recovery} applies exactly to every
scalar band with \(K\le1\).  No momentum-side \(m_p\to\infty\) statement is
used or available.  If a later branch retains the same frozen momentum
object, the Paper-II-A chronological-persistence theorem gives
\((\mathrm{FRZ}_p)\), and this same level remains fixed along that branch.
\item Let \(F_x(\widetilde t)\) be an occurring finite frontier in the sense
of \((\mathrm{HIST}_x^{\rm fin})\) of Paper~II-C.  The minimal physical
condition for the fixed-band claim is
\[
 (\mathrm{BAND}_x[K,\widetilde t]):
 \qquad \mathcal P_x(K)=\mathcal B_K\subseteq F_x(\widetilde t).
\]
Under this direct inclusion the scalar parent operators on \(\mathcal B_K\)
are recovered exactly.  On a complete-band history this condition is
equivalent to the derived statistic \(m_x(\widetilde t)\ge K\); no tower or
free integer-trajectory premise is needed.
\item This spectral conclusion does not imply local flatness.  A local
Euclidean metric approximation on an exact round branch additionally requires
the explicit defect bound of
Proposition~\ref{p2d:prop:exact-round-flatness}.  The canonical cotangent or
principal-symbol variable needs no compact-operator pairing.  Only a claim of
compact transvection provenance uses Paper~II-C's
first-symbol pairing theorem; \((\mathrm{OBS}^{(1)}_{xp})\) labels only a
particular performed comparison token, while a stronger finite-band
comparison must supply its own map and error topology.
\end{enumerate}
\end{corollary}

\begin{proof}
For the structural momentum-role seed, substitute the upstream identity \(m_p=1\)
into Theorem~\ref{p2d:thm:scalar-fixed-band-recovery}; this proves (i) for
all \(K\le1\).  If a later branch retains the same momentum object,
chronological persistence supplies \((\mathrm{FRZ}_p)\) as a theorem
conclusion.  For position, the direct band inclusion
places the entire input band in the occurring frontier, and the same theorem
proves (ii).  Paper~II-C derives the complete mathematical tower and the
resolution statistic from that frontier.  Finally, the recovery theorem is
only an exact compression identity on a fixed round-parent band.  It contains
neither a geometric defect estimate on a region of size \(L\) nor a physical
cross-carrier observable occurrence, which are distinct claim-specific
conditions in (iii).
\end{proof}

\begin{remark}[Boundary of the exact theorem]
\label{p2d:rem:recovery-boundary}
Theorem~\ref{p2d:thm:scalar-fixed-band-recovery} is finite-exact and
band-stabilized.  It is not an operator-norm convergence theorem for the full
matrix algebra, a product theorem, or a dynamics theorem.  If a finite
coordinate product, propagator, or interacting evolution reaches the terminal
shell, a valid smooth-parent statement must compare the finite operator with
the parent band through the declared inclusion, for example
\[
 \bigl\|J_{m,B}^{*}\mathcal O_mJ_{m,B}
       -P_B\mathcal O_{\rm par}P_B\bigr\|
 \le r_B(m),
 \qquad r_B(m)\longrightarrow0,
\]
with \(J_{m,B}\), its domain, and the remainder proved separately.
\end{remark}

\begin{corollary}[Guard-banded nonlinear recovery]
\label{p2d:cor:guard-banded-nonlinear-recovery}
Let the input band be \(\mathcal B_K=\mathcal P_x(K)\), and let a
multiplication word contain symbols
\(f_j\in\mathcal P_x(s_j)\).  If
\[
 m\ge K+\sum_j s_j,
\]
then the finite compression of that word agrees exactly with the parent
multiplication word on \(\mathcal B_K\).  Hence its recovery error
\(\varepsilon_{\rm rec}\) is zero on the declared band.  This statement does
not extend to inputs on the terminal shell, where the output may leave the
carrier.
\end{corollary}

\begin{proof}
Paper~II-C, Corollary~\ref{p2c:cor:nonlinear-position-guard-band}, places the
range of the word in
\(\mathcal P_x(K+\sum_js_j)\).  At every truncation level containing that
guard band, both comparison projections act as the identity on the entire
calculation.
\end{proof}


\subsubsection{Explicit recovery budgets and the fixed-momentum boundary}

\begin{theorem}[Finite word and scalar-symbol recovery certificates]
\label{p2d:thm:v233-recovery-certificates}
Let \(\Pi_m\) project onto the complete scalar position band
\(\mathcal P_x(m)\), let \(u\in\mathcal P_x(K)\), and let
\(f_j\in\mathcal P_x(s_j)\). Then
\[
 m\ge K+\sum_{j=1}^r s_j
 \quad\Longrightarrow\quad
 (\Pi_mM_{f_r}\Pi_m)\cdots(\Pi_mM_{f_1}\Pi_m)u
       =M_{f_r}\cdots M_{f_1}u .
\]
On the stage history of Theorem~\ref{p2c:thm:v233-stagewise-band-growth},
the same exact certificate is \(n+1\ge K+\sum_j s_j\).

For the positive scalar Berezin kernel of the complete band, let
\(\mathcal B_m\) denote its averaging map. If \(f\) has Lipschitz
constant \(L_f\) for physical chord distance on \(S^3(R_x)\), then
\[
 \|\mathcal B_m f-f\|_\infty
       \le L_fR_x\sqrt{\frac6{2m+3}}.
\]
For any \(\varepsilon>0\), it is sufficient that
\[
 m\ge
 \max\left\{0,\left\lceil
 \frac{6L_f^2R_x^2/\varepsilon^2-3}{2}
 \right\rceil\right\}.
\]
For position-only tests, tensoring with the fixed momentum carrier
\(m_p=1\) adds no momentum error term.
\end{theorem}
\begin{proof}
Each intermediate multiplication increases polynomial degree by
at most the symbol's degree. Under the displayed guard, every
intermediate vector remains in the band, so each projection is the
identity. This is the exact guard of
Corollary~\ref{p2d:cor:guard-banded-nonlinear-recovery}, with
symbol degrees stated separately.

For probability Haar measure, set
\[
 C_m(s)=\sum_{j=0}^m(j+1)U_j(s),\quad
 d_m=\frac{(m+1)(m+2)(2m+3)}6,\quad
 k_m(q,q')=\frac{|C_m(q\cdot q')|^2}{d_m}.
\]
The reproducing identity gives \(\int k_m(q,q')\,d\sigma(q')=1\).
Multiplying the finite character weight sums gives
\(U_jU_k=\sum_{r=0}^{\min(j,k)}U_{j+k-2r}\).
Thus the coefficient of \(U_1\) in \(C_m^2\) is
\[
 2\sum_{k=0}^{m-1}(k+1)(k+2)
       =\frac{2m(m+1)(m+2)}3,
\]
and the first harmonic multiplier is \(2m/(2m+3)\).
Consequently
\[
 \int|q-q'|^2 k_m(q,q')\,d\sigma(q')
        =\frac6{2m+3}.
\]
Positivity, the Lipschitz inequality and Cauchy--Schwarz now give
the uniform bound. Solving it for the integer \(m\) gives the
certificate. These are the one-factor specialization of Paper~I-C's
Eqs.~\eqref{p1b:der:g09:eq-G09-28} and
\eqref{p1b:der:g09:eq-G09-50}; all normalizations are displayed
here. Tensoring a position operator with an identity preserves
its norm and equality.
\end{proof}

\begin{remark}[The topology of recovery and the separate curvature error]
\label{p2d:rem:v233-recovery-limits}
These results do not imply whole-carrier operator-norm convergence
of coordinate relations. For unit-sphere compressed coordinates
\(C_a=\Pi_mM_{q_a}\Pi_m\), with \(Q_m\) the top-shell projection,
Paper~I-C, Theorems~\ref{p1b:thm:scalar-cutoff-identities} and
\ref{p1b:thm:all-level-boundaries}, gives
\[
 I-\sum_a C_a^2=\frac{m+2}{2(m+1)}Q_m,\qquad
 [C_a,C_b]=\frac{1}{2(m+1)}
       (q_a\partial_b-q_b\partial_a)Q_m.
\]
Their norms tend to \(1/2\), whereas both defects vanish strongly
and the normalized radius-defect trace is \(3/(2m+3)\).
The norm statement uses the rotation weight \(m\) on the top shell;
strong convergence follows first on fixed finite bands, then by
density and uniform boundedness. Thus the preceding symbol estimate
has its stated uniform-symbol topology.
For a general projection \(P\), \(Q=I-P\), the exact product defect
is \(PAP\,PBP-PABP=-PAQBP\).
An isometric derivation commuting with \(\Pi_m\) instead satisfies
\[
 [\Pi_mD\Pi_m,\Pi_mM_f\Pi_m]=\Pi_mM_{Df}\Pi_m
\]
without a single-commutator degree guard.

A genuinely two-factor scalar observable retains the separate
momentum concentration term in the product kernel; holding \(m_p=1\)
does not make that term tend to zero.
For a round geodesic patch \(0\le r\le L\le R_x\),
\[
 0\le1-\left(\frac{\sin(r/R_x)}{r/R_x}\right)^2
       \le\frac{L^2}{3R_x^2}.
\]
The continuous value at \(r=0\) is understood. The lower estimate
\(\sin t\ge t-t^3/6\) proves the upper bound.
Increasing \(m\) does not change this curvature error.
Actual frontier occurrence and physical time assignment are still
needed to say when either accuracy certificate is reached.
\end{remark}


\begin{proposition}[Exact round \(C^0\) flatness defect]
\label{p2d:prop:exact-round-flatness}
Let \(p\in S^3(R_x)\) and \(0<L<\pi R_x\).  On the Euclidean ball
\(B_L(0)\subset T_pS^3\), the exponential chart obeys
\[
 \left(\frac{\sin(L/R_x)}{L/R_x}\right)^2\delta
 \le \exp_p^*g_{R_x}\le\delta .
\]
Equivalently, in the Euclidean operator norm,
\begin{equation}
\label{p2d:eq:exact-round-flatness-defect}
 \left\|\exp_p^*g_{R_x}-\delta\right\|_{\delta,{\rm op}}
 \le
 \varepsilon_g(L,R_x)
 :=1-\left(\frac{\sin(L/R_x)}{L/R_x}\right)^2.
\end{equation}
On the open ball the right-hand side is the sharp supremum, approached in a
tangential direction as \(r\uparrow L\); it is not attained there.  Thus the
round-branch predicate
\((\mathrm{FLAT}^{(0)}_x[p;L,\varepsilon_g])\) is an explicit ratio test,
not an independent hypothesis.  For \(L/R_x\le1\), the sufficient estimate
\(\varepsilon_g(L,R_x)\le L^2/(3R_x^2)\) holds.  Higher \(C^k\), connection,
Laplacian, or dynamical approximations require their corresponding norms and
derivative estimates.
\end{proposition}

\begin{proof}
In geodesic polar coordinates,
\[
 \delta=dr^2+r^2d\Omega_2^2,
 \qquad
 \exp_p^*g_{R_x}=dr^2+R_x^2\sin^2(r/R_x)d\Omega_2^2.
\]
Relative to \(\delta\), the radial eigenvalue is one and the two tangential
eigenvalues are
\((\sin(r/R_x)/(r/R_x))^2\).  This factor decreases for
\(0\le r<\pi R_x\), giving the sharp estimate.  The small-window bound
follows from \(\sin\rho/\rho\ge1-\rho^2/6\) for \(0\le\rho\le1\).
\end{proof}


\label{p2:part:symmetries-limits}

This part collects structures that depend on the finite kinematic algebra but
are not new constitutional assumptions: chiral and discrete algebra symmetries,
Casimirs, the classical limit, scalar Born-pair status, and an explicit
classification of what is derived, chosen, imported, or deferred.  The POVM
and uncertainty constructions remain owned by II-C and are not re-derived in
this part.


\section{Central grading, Casimirs and discrete symmetries}
\label{paper2d:sec:organized-symmetries}

\subsection{Symmetries of the Kinematic Algebra}
\label{p2:sec:symmetries}
The geometric parent is still the single sphere
$S^3\cong SU(2)$.  Its connected geometric isometry group is
$\SO(4)\cong(\SU(2)_L\times\SU(2)_R)/\Delta\ZZ_2$.  On the finite module
package, this isometry action is lifted to
$\Spin(4)\cong\SU(2)_L\times\SU(2)_R$, because the protected sectors include
half-integer spinorial representations that do not descend to honest
$\SO(4)$ representations.  Thus $\Spin(4)$ is a representation symmetry, not
the carrier and not a replacement of $S^3$ by $S^3\times S^3$.


\subsubsection{Chiral Decomposition}

The isometry algebra of $S^3$ admits a canonical decomposition into commuting
left and right chiral symmetry sectors.  These are Killing-symmetry sectors,
not two independent operation-carrier sectors.

\begin{proposition}[Chiral Factorization]
\label{p2:prop:chiral}
The Lie algebra $\mathfrak{so}(4)$ decomposes as a direct sum:
\begin{equation}
\mathfrak{so}(4) \cong \mathfrak{su}(2)_L \oplus \mathfrak{su}(2)_R.
\end{equation}
The dimensionless chiral generators are:
\begin{equation}
L^+_i := \frac{1}{2}\left(\frac{\hat{J}_i}{\hbar} + \frac{\Rmom\hat{P}_i}{\hbar}\right), \quad 
L^-_i := \frac{1}{2}\left(\frac{\hat{J}_i}{\hbar} - \frac{\Rmom\hat{P}_i}{\hbar}\right).
\end{equation}
Inverting:
\begin{equation}
\hat{J}_i = \hbar(L^+_i + L^-_i), \quad \hat{P}_i = \frac{\hbar}{\Rmom}(L^+_i - L^-_i).
\end{equation}
\end{proposition}

\begin{proof}
Direct computation from the three commutation relations
\(
[\hat{J}_i,\hat{J}_j]=i\hbar\epsilon_{ijk}\hat{J}_k
\),
\(
[\hat{J}_i,\hat{P}_j]=i\hbar\epsilon_{ijk}\hat{P}_k
\)
and
\(
[\hat{P}_i,\hat{P}_j]
=\frac{i\hbar}{\Rmom^2}\epsilon_{ijk}\hat{J}_k
\).
All three enter: expanding,
\[
 4\hbar^{2}[L^{+}_i,L^{+}_j]
 =[\hat J_i,\hat J_j]
 +\Rmom[\hat J_i,\hat P_j]
 +\Rmom[\hat P_i,\hat J_j]
 +\Rmom^{2}[\hat P_i,\hat P_j],
\]
in which the rotation bracket appears with coefficient one.
\end{proof}

\begin{corollary}[Lifted representation labels and descent to \texorpdfstring{$\SO(4)$}{SO(4)}]
\label{p2:cor:spin4-labels-so4-descent}
Irreducible sectors of the lifted representation group
\[
\Spin(4)\cong \SU(2)_L\times \SU(2)_R
\]
are labeled by pairs $(j_L,j_R)$ with $j_L,j_R\in\frac12\ZZ_{\ge0}$, and have dimension $(2j_L+1)(2j_R+1)$ \cite{fulton1991representation,Hall2015lie}.  Such a sector descends to an honest representation of the quotient
\[
\SO(4)=\Spin(4)/\{\pm(1,1)\}
\]
if and only if the common central element acts trivially, equivalently
\[
j_L+j_R\in\ZZ.
\]
Consequently the protected spinor sectors $(\tfrac12,0)$ and $(0,\tfrac12)$ are \(\Spin(4)\)-sectors rather than ordinary \(\SO(4)\)-representations.
\end{corollary}

\begin{proof}
The product decomposition of \(\Spin(4)\) gives irreducibles
$V_{j_L}\otimes V_{j_R}$ with $j_L,j_R\in \frac12\ZZ_{\ge0}$, and the stated dimension is the product of the two \(SU(2)\) dimensions.  The element $(-1,-1)\in SU(2)_L\times SU(2)_R$ acts on this tensor product by
\[
(-1)^{2j_L}(-1)^{2j_R}=(-1)^{2(j_L+j_R)}.
\]
It acts trivially precisely when $2(j_L+j_R)$ is even, i.e., when $j_L+j_R\in\ZZ$.  This is exactly the descent condition for the quotient by \(\{\pm(1,1)\}\).
\end{proof}

\begin{remark}[Derived lift and operator-level descent]
\label{p2:rem:sector-grammar-sync}
The half-integral labels above are the finite unitary sector grammar of the
derived Killing algebra established in Paper~II-B; each sector integrates to
the derived universal lift \(\Spin(4)\).  This is representation structure of
the single \(S^3\) parent, not a covering parent inserted by hand.  Although
a protected spinor sector does not itself descend to an \(SO(4)\)
representation, conjugation on its endomorphism algebra does descend because
the diagonal central element acts scalarly.  On a direct sum, arbitrary maps
between opposite central-character classes need not descend.  The Killing
generators, equivariant projectors, and coordinate channels used here preserve
the central character, so all operator-level covariance statements are made
inside the precise descent domain of Paper~II-B.
\end{remark}

\begin{theorem}[Canonical diagonal-center grading at level one]
\label{p2:thm:diagonal-center-grading}
On
\[
 \mathcal H_1=
 V_{0,0}\oplus V_{\frac12,0}\oplus V_{0,\frac12}
 \oplus V_{\frac12,\frac12}\oplus V_{1,0}\oplus V_{0,1},
\]
let \(z_c=(-I,-I)\in\Spin(4)\) and define
\[
 \Gamma_c:=\rho_1(z_c)
 =P_{0,0}-P_{\frac12,0}-P_{0,\frac12}
  +P_{\frac12,\frac12}+P_{1,0}+P_{0,1}.
\]
Then \(\Gamma_c=\Gamma_c^*=\Gamma_c^{-1}\), and
\[
 P_\pm=\frac{1\pm\Gamma_c}{2},
 \qquad
 \mathcal H_+=P_+\mathcal H_1=\mathcal H_{\rm geo},
 \qquad
 \mathcal H_-=P_-\mathcal H_1=\mathcal H_{\rm sp},
\]
with
\[
 \boxed{\dim\mathcal H_+=11,\qquad\dim\mathcal H_-=4.}
\]
The constituted rotation/transvection generators, their generated algebra,
the equivariant sector projectors, and every current coordinate channel
commute with \(\Gamma_c\).  Hence every operator generated by the current
kinematic grammar is central-even and
\[
 \mathcal A_{\rm nat}
 \subseteq
 \End(\mathcal H_+)\oplus\End(\mathcal H_-)
 \cong M_{11}(\CC)\oplus M_4(\CC).
\]
\end{theorem}

\begin{proof}
On \(V_{j_L,j_R}\), the common central element acts as
\((-1)^{2(j_L+j_R)}\), giving the displayed signs and dimensions.  The
Killing generators commute with every central representation element.  A
coordinate vector has type \((\tfrac12,\tfrac12)\), on which \(z_c\) acts
trivially, so equivariance gives
\(\Gamma_cX_a\Gamma_c=X_a\).  Sums, products, adjoints, and generated
functional calculus preserve this commutation relation.
\end{proof}

The projectors $P_+=P_g$ and $P_-=P_s$ also supply the central tags of
Paper~II-C, Theorem~\ref{p2c:thm:v2211-central-sector-refinement}.
Corollary~\ref{p2c:cor:v2211-tagged-intrinsic-populations} proves that
refining the quaternion or Hurwitz localization rule by these tags preserves
its exact coordinate moments and separates their invariant populations.
The untagged first-moment rule does not separate that direction.  This is a
claim-relative measurement construction, not a preparation theorem, complete
tomography statement, or identification of this central grading with physical
spacetime chirality.

\begin{corollary}[Central-even and central-odd operator spaces]
\label{p2:cor:central-even-odd-operator-splitting}
For \(\Theta(T):=\Gamma_cT\Gamma_c\),
\[
 \End(\mathcal H_1)^{\bar0}
 =\End(\mathcal H_+)\oplus\End(\mathcal H_-),
 \qquad \dim_{\CC}=137,
\]
whereas
\[
 \End(\mathcal H_1)^{\bar1}
 =\operatorname{Hom}(\mathcal H_-,\mathcal H_+)
 \oplus\operatorname{Hom}(\mathcal H_+,\mathcal H_-),
 \qquad \dim_{\CC}=88.
\]
Thus \(137+88=225=15^2\).
\end{corollary}

\begin{proof}
Relative to \(\mathcal H_1=\mathcal H_+\oplus\mathcal H_-\), conjugation by
\(\Gamma_c=\operatorname{diag}(I_{11},-I_4)\) fixes the two diagonal blocks
and negates the two rectangular blocks.  Their dimensions are
\(11^2+4^2=137\) and \(2\cdot11\cdot4=88\), respectively.
\end{proof}

\begin{lemma}[No invariant central-odd bridge]
\label{p2:lem:no-invariant-central-odd-bridge}
If \(Y\in\End(\mathcal H_1)\) is \(\Spin(4)\)-invariant and central-odd,
then \(Y=0\).
\end{lemma}

\begin{proof}
Invariance gives \(\Gamma_cY\Gamma_c=Y\), whereas oddness gives
\(\Gamma_cY\Gamma_c=-Y\).
\end{proof}

\begin{remark}[Distinction from other parity and grading operators]
\label{p2:rem:diagonal-center-grading-distinction}
The involution \(\Gamma_c\) is not spatial parity, antipodal harmonic parity,
the protected chiral grading \(\Gamma_{\rm sp}\), or the Cartan weight
grading.  A protected Clifford coordinate may be chiral-odd while remaining
central-even.  The two central-character eigenspaces are not called bosonic
and fermionic here; that interpretation would require a later spin--statistics
theorem.
\end{remark}


\begin{theorem}[Antipodal lifts versus the diagonal center]
\label{p2d:thm:antipodal-lifts-center}
In \(\Spin(4)=SU(2)_L\times SU(2)_R\), set
\[
 z_L=(-I,I),\qquad z_R=(I,-I),\qquad z_c=(-I,-I)=z_Lz_R.
\]
Under the standard action \((a,b):g\mapsto agb^{-1}\) on \(S^3\cong SU(2)\),
\[
 z_L:g\mapsto-g,\qquad z_R:g\mapsto-g,\qquad z_c:g\mapsto g.
\]
Thus \(z_L\) and \(z_R\) are the two spin lifts of the geometric antipodal
isometry, whereas \(z_c\) is the kernel element of
\(\Spin(4)\to SO(4)\).  Consequently antipodal harmonic parity and the
level-one diagonal-center \(11\oplus4\) grading are distinct generated
involutions.
\end{theorem}

\begin{proof}
The three formulas follow directly from \(agb^{-1}\) and centrality of
\(-I\in SU(2)\).  The last statement follows because
Theorem~\ref{p2:thm:diagonal-center-grading} represents \(z_c\), whereas the
antipodal action is represented by either \(z_L\) or \(z_R\).
\end{proof}

\begin{corollary}[Oriented Hopf generators from rotation and transvection]
\label{p2d:cor:hopf-rotation-transvection}
Use the Pauli parametrization
\(g=y_4I+i y_i\sigma_i\), with \(y\in S^3(1)\), for the round parent.
For an oriented unit Cartan axis \(n\), put
\(h_n(\alpha)=\exp(i\alpha n_i\sigma_i)\), \(\alpha\sim\alpha+2\pi\).
Define the scalar pullbacks
\[
 (U_L(\alpha)f)(g)=f(h_n(\alpha)g),\qquad
 (U_R(\alpha)f)(g)=f(gh_n(\alpha)),
 \qquad U_{L/R}(\alpha)=e^{i\alpha Q_{H,L/R}}.
\]
With the existing generators
\[
 \widehat J_n=\hbar(L_n^++L_n^-),\qquad
 \widehat P_n=\frac{\hbar}{\Rmom}(L_n^+-L_n^-),
\]
their geometric action generators satisfy
\[
 \boxed{\hbar Q_{H,L}=\Rmom\widehat P_n-\widehat J_n
       =-2\hbar L_n^-,\qquad
        \hbar Q_{H,R}=\Rmom\widehat P_n+\widehat J_n
       =2\hbar L_n^+.}
\]
The right action here is \(g\mapsto gh_n(\alpha)\); replacing it by
\(g\mapsto gh_n(\alpha)^{-1}\) reverses its generator.
The algebraic labels on \(L^\pm\) are retained, with this explicit
dictionary to the geometric multiplication sides.  An atlas that instead
labels irreducibles by geometric left/right multiplication factors is
translated by \((j_{\rm geom,L},j_{\rm geom,R})=(j_R,j_L)\).
Every operation descending from the same marked Hopf occurrence inherits
its axis and oriented action.
\end{corollary}

\begin{proof}
Write
\[
 \mathcal R_i=\epsilon_{ijk}y_j\partial_{y_k},\qquad
 \mathcal B_i=y_4\partial_{y_i}-y_i\partial_{y_4}.
\]
The coordinate commutators give
\(\widehat J_i/\hbar=-i\mathcal R_i\) and
\(\Rmom\widehat P_i/\hbar=-i\mathcal B_i\).
Direct Pauli multiplication gives the infinitesimal fields
\(\mathcal V_i^L=\mathcal B_i-\mathcal R_i\) and
\(\mathcal V_i^R=\mathcal B_i+\mathcal R_i\).
Their Hermitian generators are \(-i\mathcal V_i^{L/R}\).
Contracting with \(n_i\) proves the formulas.
\end{proof}

\begin{remark}[Chirality Symmetry]
The exchange $L \leftrightarrow R$ (equivalently $\hat{P}_i \to -\hat{P}_i$) is an outer automorphism of $\mathfrak{so}(4)$. The round metric and spinor bundle are invariant under this exchange.  A selected Hopf $U(1)$, contact form, or Hopf-charge convention is additional structure and may be exchanged or relabeled rather than fixed.
\end{remark}

\subsubsection{Casimir Operators}

The quadratic center is generated by the two chiral Casimirs below.
Arbitrary invariant operators are not thereby unique; on a finite
representation they may be functions of the central spectral data and may
also act on multiplicity spaces.  The stronger uniqueness statement needed
for scalar second-order principal symbols is proved separately in
Theorem~\ref{p2:thm:operator-rigidity}.

\begin{proposition}[Casimir Operators]
\label{p2:prop:casimir}
The $\mathfrak{so}(4)$ algebra has two independent Casimir operators: \cite{Hall2015lie}
\begin{equation}
C_L = \sum_i (L^+_i)^2, \quad C_R = \sum_i (L^-_i)^2.
\end{equation}
On the representation \((j_L,j_R)\):
\begin{equation}
C_L = j_L(j_L+1), \quad C_R = j_R(j_R+1).
\end{equation}
\end{proposition}

\begin{proof}
Each $\mathfrak{su}(2)$ factor has quadratic Casimir $C=\sum_i J_i^2$ commuting with the algebra.
For the dimensionless generators used here, the spin-$j$ irreducible representation has
\[
C=j(j+1)\,\mathbf 1.
\]
Equivalently, physical angular-momentum generators $\hbar J_i$ would give the eigenvalue $\hbar^2j(j+1)$.
Since the left and right factors commute, $C_L$ and $C_R$ are independent and act with the stated dimensionless eigenvalues on $(j_L,j_R)$.
\end{proof}

\begin{theorem}[Invariant scalar second-order operator]
\label{p2:thm:operator-rigidity}
Let $S^3(\Rmom)\cong G/K$, where $G=\Spin(4)$ and
$K=\SU(2)_{\mathrm{diag}}$.  On scalar functions, the principal
second-order part of a $G$-invariant differential operator is
one-dimensional.  Hence the Laplace--Beltrami operator is unique up to scale
and an additive zeroth-order constant.  With the standard radius
normalization,
\begin{equation}
\Delta_{S^3}=\frac{2}{\Rmom^2}(C_L+C_R).
\end{equation}
\end{theorem}

\begin{proof}
For a homogeneous space $G/K$, invariant scalar second-order principal symbols correspond to $K$-invariant elements of $S^2(\mathfrak p)$. Here $K=\SU(2)_{\mathrm{diag}}$ acts on $\mathfrak p$ through its effective quotient $K/\{\pm I\}\cong SO(3)$, and $\mathfrak p$ is the irreducible $3$-dimensional vector representation of that quotient, so the invariant subspace is one-dimensional.  After subtracting the multiple of \(\Delta_{S^3}\) having the same principal symbol, the remainder has order at most one.  Its first-order symbol would be a \(K\)-invariant element of \(\mathfrak p\) (equivalently of \(\mathfrak p^*\)); none exists because
\[
\mathfrak p^K=0.
\]
Thus the first-order part vanishes.  A \(G\)-invariant zeroth-order coefficient is constant because \(G\) acts transitively on \(G/K\).  Fixing the Laplacian normalization and the zero mode removes this additive freedom. The coefficient in the displayed normalization is fixed on the scalar diagonal sector: for $(j_L,j_R)=(k/2,k/2)$,
\[
\frac{2}{\Rmom^2}(C_L+C_R)
=\frac{2}{\Rmom^2}\left(2\cdot\frac{k}{2}\left(\frac{k}{2}+1\right)\right)
=\frac{k(k+2)}{\Rmom^2},
\]
which is the standard round-sphere scalar Laplace eigenvalue.
\end{proof}

\begin{remark}[Scope of scalar second-order rigidity]
\label{p2d:rem:scalar-rigidity-first-order-firewall}
Theorem~\ref{p2:thm:operator-rigidity} is a theorem about the principal
second-order part on scalar functions.  It explicitly leaves an additive
constant until the scalar zero-mode normalization is imposed.  It does not
imply rigidity of a first-order operator on the reducible forms-plus-spinors
bundle.  Such an operator admits equivariant bundle-valued zero-order terms,
as proved below; its principal symbol and scalar commutators therefore do not
determine its square.
\end{remark}

\begin{corollary}[Spectral formula for scalar harmonics]
\label{p2:cor:laplace-spectrum}
On the diagonal scalar-harmonic sector $(j_L,j_R)=(k/2,k/2)$, the Laplacian eigenvalue is
\begin{equation}
\lambda_k=\frac{k(k+2)}{\Rmom^2}.
\end{equation}
\end{corollary}

\begin{proof}
Insert $j_L=j_R=k/2$ into the Casimir formula.
\end{proof}

\subsubsection{Equivariant scalars and absence of a preferred basis}
\label{p2:subsec:no-preferred-basis}

\begin{theorem}[Schur factorization of equivariant operators]
\label{p2:thm:no-preferred-basis}
Write the finite kinematic carrier in isotypic form
\[
\mathcal H_m
=
\bigoplus_{\lambda}\bigl(V_\lambda\otimes M_\lambda\bigr),
\]
where the \(V_\lambda\) are pairwise inequivalent irreducible \(\Spin(4)\)-modules and the \(M_\lambda\) are finite multiplicity spaces. Then every complex-linear operator \(T\) commuting with the \(\Spin(4)\)-action has the form
\[
T
=
\bigoplus_{\lambda}\bigl(\mathbf 1_{V_\lambda}\otimes B_\lambda\bigr)
\]
for unique \(B_\lambda\in \End(M_\lambda)\).
\end{theorem}

\begin{proof}
This is the standard commutant factorization for a semisimple representation,
that is, Schur's lemma applied to the isotypic decomposition.
\end{proof}

\begin{corollary}[Self-adjoint blocks and the residual basis freedom]
\label{p2d:cor:selfadjoint-schur-blocks}
In the factorization of Theorem~\ref{p2:thm:no-preferred-basis},
\(T=T^*\) if and only if every \(B_\lambda=B_\lambda^*\).  The representation
itself supplies no basis of \(M_\lambda\).  A particular self-adjoint
\(B_\lambda\) supplies its spectral subspaces, but bases inside degenerate
eigenspaces remain noncanonical.
\end{corollary}

\begin{proof}
Taking adjoints blockwise gives
\(T^*=\bigoplus_\lambda(\mathbf1\otimes B_\lambda^*)\), proving the first
claim.  The spectral theorem gives only the eigenspace decomposition; the
unitary group of each degenerate eigenspace preserves the operator.
\end{proof}

\begin{corollary}[Invariant scalar functionals are spectral]
\label{p2d:cor:schur-scalar-functionals}
Every scalar functional of an equivariant self-adjoint \(T\) that is invariant
under independent unitary conjugation on the multiplicity spaces factors
through the finite eigenvalue multisets of the blocks \(B_\lambda\), including
their multiplicities.  Hence trace polynomials and other symmetric spectral
functions are basis-independent, but the kinematics selects no eigenbasis.
\end{corollary}

\begin{proof}
The unitary-conjugacy class of a finite-dimensional self-adjoint matrix is
exactly its eigenvalue multiset.  Apply this independently to every block.
\end{proof}


\subsubsection{Discrete Symmetries}
\label{p2:sec:discrete-symmetries}

The smooth parent representation on \(S^3\) admits geometric parity and
time-reversal operations, and their sign assignments define automorphisms of
the abstract compact kinematic relations.  A finite coordinate completion is
additional operator data.  At level one its selector-relevant realization is
first the unmarked class
\[
 [\widehat X]\in
 \mathfrak M_{\rm coord}(\mathcal H_1)/{\sim}
 \cong\mathbb R_{\ge0}^{3},
 \qquad
 [\widehat X]\leftrightarrow
 (r_{1,0},r_{0,1},\varrho_{\rm sp}),
\]
and only secondarily a phase-marked representative \(\widehat X\) of that
class.  Discrete symmetry must therefore be tested first by equality of
transported realization classes and, when a literal implementing
representative is required, by the stronger equality of chosen
representatives.  The parent/abstract statements and these finite-completion
conditions are recorded separately.  Parity exchanges the selected left- and
right-fiber presentations, whereas anti-unitary time reversal reverses the
Hermitian left-fiber Hopf momentum.

Spatial parity is a concrete orientation-reversing isometry of the compact
momentum sphere.  Time reversal is used here as an anti-unitary algebra
automorphism, not as the base-space inversion map \(g\mapsto g^{-1}\), since
inversion on \(\SU(2)\) flips the spatial embedding coordinates.

No geometric charge-conjugation or CPT operation is inferred from these two
maps.  An optional two-state internal label can be appended as independent
representation data, but it is not the geometric Hopf momentum.

\paragraph{Parity}

\begin{definition}[Parity on $S^3$]
\label{p2:def:parity}
Identify \(S^3\) with unit quaternions, or equivalently with
\[
g=x_4\mathbf 1+i(x_1\sigma_1+x_2\sigma_2+x_3\sigma_3)\in SU(2).
\]
The \textbf{parity operator} $\mathcal{P}$ is the orientation-reversing isometry of $S^3$ defined by
\begin{equation}
\mathcal{P}: (x_1, x_2, x_3, x_4) \mapsto (-x_1, -x_2, -x_3, x_4),
\end{equation}
equivalently by quaternionic conjugation, or matrix Hermitian conjugation,
\begin{equation}
\mathcal{P}: g\mapsto g^{-1}=g^\dagger=x_4\mathbf 1-i(x_1\sigma_1+x_2\sigma_2+x_3\sigma_3).
\end{equation}
The formula \(g\mapsto\sigma_3g\sigma_3\) is \emph{not} this spatial parity in the standard embedding convention; it fixes the third spatial axis and flips only the first two.
\end{definition}


\begin{theorem}[Parent parity action and abstract kinematic automorphism]
\label{p2:thm:parity-action}
Let \(M_{x_a}\) be the real embedding-coordinate multiplication
operators on the smooth parent.  Under spatial parity
\(\mathcal P:g\mapsto g^{-1}\),
\begin{align}
\mathcal{P}M_{x_i}\mathcal{P}^{-1} &= -M_{x_i}, \quad (i = 1, 2, 3), \\
\mathcal{P}M_{x_4}\mathcal{P}^{-1} &= +M_{x_4}, \\
\mathcal{P}\hat{P}_i\mathcal{P}^{-1} &= -\hat{P}_i, \\
\mathcal{P}\hat{J}_i\mathcal{P}^{-1} &= +\hat{J}_i.
\end{align}
Parity is a \(\ZZ_2\) symmetry of the parent representation:
\(\mathcal{P}^2 = 1\).  The displayed signs also define the corresponding
automorphism of the abstract compact kinematic relations.  They are not, by
themselves, a symmetry theorem for every finite coordinate completion; that
requires the fixed-point condition of
Proposition~\ref{p2d:prop:finite-completion-discrete-symmetry-scope}.
No intrinsic
transformation law for the geometric left-fiber Hopf charge is included
in this statement: inversion sends a left-fiber action to a right action, so
any Hopf-sector sign convention must be specified as an additional
charge-carrier lift.
\end{theorem}

\begin{proof}
\leavevmode\par\noindent
The coordinate transformation is the defining embedding-space action:
\[
  \mathcal{P}(x_i)=-x_i\quad (i=1,2,3),
  \qquad
  \mathcal{P}(x_4)=x_4.
\]
In the standard identification
\[
  g=x_4\mathbf 1+i x_i\sigma_i,
\]
this is exactly \(g\mapsto g^\dagger=g^{-1}\), hence
\(\mathcal P^2=1\).

Under inversion on the group, the two chiral triples are exchanged.  Here
\(L_i^\pm\) denote the dimensionless representation generators of the two
commuting left/right group actions used in
\[
  \hat P_i=\frac{\hbar}{\Rmom}(L_i^+-L_i^-),
\]
not raw left/right invariant differential fields with an independent sign
convention.  The exchange is not a free convention: it is forced by the
embedding-space action.  Writing \(\mathcal P\) on \(\RR^4\) as
\(O=\mathrm{diag}(-1,-1,-1,+1)\) and using the identification of \(\hat J_i\) with
the rotation generators \(M_{jk}\) and of \(\hat P_i\) with the transvection
generators \(M_{i4}\), conjugation gives
\[
  OM_{jk}O^{\mathsf T}=+M_{jk},
  \qquad
  OM_{i4}O^{\mathsf T}=-M_{i4},
\]
so \(\hat J_i\) is parity-even and \(\hat P_i\) parity-odd, and the opposite sign
choice is excluded.  Equivalently, in terms of the chiral generators, inversion
gives
\[
  \mathcal{P}L_i^+\mathcal{P}^{-1}=L_i^-,
  \qquad
  \mathcal{P}L_i^-\mathcal{P}^{-1}=L_i^+.
\]
Equivalently, if one starts from raw invariant vector fields, their usual
inversion signs are absorbed in the definition of the corresponding action
generators before forming \(L_i^\pm\).  Therefore
\[
\begin{aligned}
  \mathcal{P}\hat P_i\mathcal{P}^{-1}
  &=\frac{\hbar}{\Rmom}(L_i^- - L_i^+)
   =-\hat P_i,
  \\
  \mathcal{P}\hat J_i\mathcal{P}^{-1}
  &=\hbar(L_i^-+L_i^+)
   =\hat J_i.
\end{aligned}
\]
For the geometric Hopf generator, the relevant action is instead the chosen
left-fiber Hopf action.  Since
\[
  (h_\theta g)^{-1}=g^{-1}h_\theta^{-1},
\]
With both fibre actions oriented by the same \(h_\alpha\), this gives
\[
 \mathcal P Q_{H,L}\mathcal P^{-1}=-Q_{H,R},\qquad
 \mathcal P Q_{H,R}\mathcal P^{-1}=-Q_{H,L}.
\]
Thus, denoting the two action types explicitly,
\[
 \mathcal P:\mathcal V_q^L\longrightarrow\mathcal V_{-q}^R.
\]
This is transport to the right-fibre atlas; it does not define an intrinsic
charge-conjugation operator within one fixed left-fibre atlas.
\end{proof}

\subparagraph{Antipodal grading.}
The antipodal \(\mathbb Z_2\) grading is distinct from spatial parity; the
standard harmonic statement and proof are deferred to
Paper~II-F, Section~\ref{p2d:app:discrete-supplements}.


\paragraph{Time Reversal}

\begin{definition}[Geometric kinematic time reversal]
\label{p2:def:time-reversal}
On the smooth parent, let \(M_{x_a}\) be the real embedding-coordinate multiplication operators and let \(\mathcal T_0\)
be an anti-unitary kinematic time-reversal lift which fixes these real
embedding-coordinate multiplication operators.  Its action on the basic
generators is
\begin{align}
\mathcal T_0M_{x_i}\mathcal T_0^{-1} &= +M_{x_i}, &
\mathcal T_0M_{x_4}\mathcal T_0^{-1} &= +M_{x_4}, \\
\mathcal T_0\hat{P}_i\mathcal T_0^{-1} &= -\hat{P}_i, &
\mathcal T_0\hat{J}_i\mathcal T_0^{-1} &= -\hat{J}_i, \\
\mathcal T_0 i\mathcal T_0^{-1} &= -i.
\end{align}
For the real left-fiber vector field \(\xi\), fixing real coordinates gives
\(\mathcal T_0\xi\mathcal T_0^{-1}=\xi\).  Therefore the geometric Hopf
operator is necessarily time-reversal odd:
\begin{equation}
\label{p2:eq:geometric-time-reversal-hopf-odd}
\mathcal T_0\hat{\mathfrak h}\mathcal T_0^{-1}
=\mathcal T_0(-i\Rmom\xi)\mathcal T_0^{-1}
=+i\Rmom\xi
=-\hat{\mathfrak h}.
\end{equation}
This is not the base-space inversion map $g\mapsto g^{-1}$ on $SU(2)$; that inversion is quaternionic conjugation and flips the spatial embedding coordinates.
For a chosen phase-marked representative
\(\widehat X=(\widehat X_1,\ldots,\widehat X_4)\) of a finite completion
class, a spin lift is called \emph{representative-compatible} only when it
also satisfies
\(\mathcal T_0\widehat X_a\mathcal T_0^{-1}=\widehat X_a\) for every \(a\).
This literal fixed-point condition is stronger than class-level invariance
\([\widehat X^{\mathcal T}]=[\widehat X]\), because unmarked sector
rephasings may relate the transported and chosen representatives.  Neither
condition follows merely from the existence of a spin lift for the
\(\Spin(4)\)-representation.
\end{definition}

\begin{theorem}[Abstract time-reversal automorphism and compatible implementation]
\label{p2:thm:time-reversal-action}
The assignments in Definition~\ref{p2:def:time-reversal} define an
anti-linear automorphism of the abstract compact
coordinate--transvection--rotation relations and send the geometric Hopf
charge \(q\) to \(-q\).  If a chosen finite coordinate representative admits a
representative-compatible anti-unitary spin lift, that lift implements the
automorphism literally on that representative.  Class-level symmetry requires
only \([\widehat X^{\mathcal T}]=[\widehat X]\).  Whenever the standard
spin lift is representative-compatible,
\(\mathcal T_0^2=+1\) on integer-spin sectors and
\(\mathcal T_0^2=-1\) on half-integer-spin sectors.  The induced algebra
automorphism does not determine this implementing square.  Here a spin sector
means an irreducible spin-\(s\) summand for the rotation algebra generated by
the \(\hat J_i\); ``integer'' and ``half-integer'' refer to
\(s\in\ZZ\) and \(s\in\tfrac12+\ZZ\), respectively.
\end{theorem}

\begin{proof}
\ProofPlan{We verify the time-reversal automorphism in five steps:
(a) record the anti-unitary sign convention;
(b) check rotations and mixed brackets;
(c) check compact momentum curvature;
(d) check coordinate-transvection brackets;
(e) construct the standard lift on each \(\Spin(4)\) summand and compute its
square on every diagonal-spin sector.}

\begin{proofsteps}
\item \textbf{Anti-unitary sign convention.}
Because \(\mathcal T_0\) is anti-unitary, it sends \(i\mapsto -i\).

\item \textbf{Rotation and mixed brackets.}
The rotation brackets are preserved since
\[
[-\hat J_i,-\hat J_j]=[\hat J_i,\hat J_j]
=i\hbar\epsilon_{ijk}\hat J_k
=(-i)\hbar\epsilon_{ijk}(-\hat J_k).
\]
The mixed brackets are preserved because
\[
[-\hat J_i,-\hat P_j]=[\hat J_i,\hat P_j]
=i\hbar\epsilon_{ijk}\hat P_k
=(-i)\hbar\epsilon_{ijk}(-\hat P_k).
\]

\item \textbf{Compact momentum curvature bracket.}
The compact momentum curvature bracket is preserved because
\[
[-\hat P_i,-\hat P_j]=[\hat P_i,\hat P_j]
=\frac{i\hbar}{\Rmom^2}\epsilon_{ijk}\hat J_k
=\frac{(-i)\hbar}{\Rmom^2}\epsilon_{ijk}(-\hat J_k).
\]

\item \textbf{Coordinate-transvection brackets.}
Using the coordinate--transvection relations imported from Paper~II-C,
\[
[\hat P_i,M_{x_j}]=-\frac{i\hbar}{\Rmom}\delta_{ij}M_{x_4},
\qquad
[\hat P_i,M_{x_4}]=+\frac{i\hbar}{\Rmom}M_{x_i},
\]
one obtains
\[
[-\hat P_i,+M_{x_j}]
=+\frac{i\hbar}{\Rmom}\delta_{ij}M_{x_4}
=-\frac{(-i)\hbar}{\Rmom}\delta_{ij}(+M_{x_4}),
\]
and
\[
[-\hat P_i,+M_{x_4}]
=-\frac{i\hbar}{\Rmom}M_{x_i}
=\frac{(-i)\hbar}{\Rmom}(+M_{x_i}).
\]
Thus all defining relations are preserved anti-linearly.

\item \textbf{Geometric Hopf generator and lift square.}
Since \(\xi\) is a real differential vector field fixed by the real-coordinate
action while anti-unitarity sends \(i\) to \(-i\),
\[
\mathcal T_0(-i\Rmom\xi)\mathcal T_0^{-1}
=+i\Rmom\xi
=-\hat{\mathfrak h}.
\]
Thus the geometric charge reversal is forced and is not an optional sign
convention.  The statement about \(\mathcal T_0^2\) belongs to the chosen
spin lift rather than to the abstract sign automorphism.  Because a finite
unitary \(\Spin(4)\)-module is an orthogonal direct sum of irreducibles
\(V_{j_L}\otimes V_{j_R}\), it suffices to construct the lift on one such
summand.  On a standard spin-\(j\) \(\SU(2)\)-module, define
\[
\Theta_j\!\left(\sum_{m=-j}^{j}c_m|j,m\rangle\right)
=
\sum_{m=-j}^{j}(-1)^{j-m}\overline{c_m}\,|j,-m\rangle .
\]
Writing \(S_i^{(j)}\) for the standard dimensionless spin generators, the
ladder relations and one second application of the formula give
\[
 \Theta_j S_i^{(j)}\Theta_j^{-1}=-S_i^{(j)},
 \qquad
 \Theta_j^2=(-1)^{2j}\mathbf 1_{V_j}.
\]
On \(V_{j_L}\otimes V_{j_R}\), take the anti-unitary product lift
\[
 \Theta_{j_L,j_R}:=\Theta_{j_L}\otimes\Theta_{j_R}.
\]
The pure-tensor rule is well defined anti-linearly, and
\[
 \big\langle\Theta_{j_L,j_R}(v\otimes w),
 \Theta_{j_L,j_R}(v'\otimes w')\big\rangle
 =
 \overline{\big\langle v\otimes w,v'\otimes w'\big\rangle},
\]
so the product lift is anti-unitary.
Here
\(L_i^+=S_i^{(j_L)}\otimes\mathbf 1\) and
\(L_i^-=\mathbf 1\otimes S_i^{(j_R)}\), so
\[
 \Theta_{j_L,j_R}L_i^\pm\Theta_{j_L,j_R}^{-1}=-L_i^\pm.
\]
The inverse chiral relations from
Proposition~\ref{p2:prop:chiral} therefore give both required signs,
\[
 \Theta_{j_L,j_R}\hat J_i\Theta_{j_L,j_R}^{-1}=-\hat J_i,
 \qquad
 \Theta_{j_L,j_R}\hat P_i\Theta_{j_L,j_R}^{-1}=-\hat P_i.
\]
Moreover,
\[
 \Theta_{j_L,j_R}^2
 =(-1)^{2(j_L+j_R)}\mathbf 1.
\]
Under the diagonal rotation algebra,
\[
 V_{j_L}\otimes V_{j_R}
 =\bigoplus_{s=|j_L-j_R|}^{j_L+j_R}V_s,
 \qquad j_L+j_R-s\in\ZZ_{\ge0}.
\]
The product lift preserves each \(V_s\), since it fixes the diagonal
quadratic Casimir, and the displayed integrality gives
\[
 (-1)^{2(j_L+j_R)}=(-1)^{2s}.
\]
Taking the orthogonal direct sum of these product lifts proves the asserted
\(+1\) square on every integer-spin summand and \(-1\) square on every
half-integer-spin summand.  This construction fixes the representation-level
lift; its literal compatibility with a chosen coordinate representative is
the separate representative fixed-point premise stated in the theorem.
\end{proofsteps}
\end{proof}

\begin{remark}[Kinematic reversal and line reality]
\label{p2d:rem:reversal-line-reality-scope}
The preceding kinematic reversal retains its own charge and
antiunitary action.  The later quadratic-line construction in
Paper IV gives an antiunitary involution only after a nonzero
symmetric line form is supplied.  These are different input
signatures.  Neither the present reversal nor that line result
constitutes a global physical \(J_F\), a nonzero same-field
Majorana pairing or a quantum PCT theorem.
\end{remark}


\begin{proposition}[Discrete-symmetry scope of a finite coordinate completion]
\label{p2d:prop:finite-completion-discrete-symmetry-scope}
Let \(\widehat X=(\widehat X_1,\ldots,\widehat X_4)\) be a compatible
Hermitian \(\Spin(4)\)-equivariant coordinate vector on a finite auxiliary
module \(\mathcal H_m\).  Let \(U_{\mathcal P}\) implement the chiral
exchange and let \(\Theta\) be a representation-level time-reversal spin
lift.
For transport within the fixed scalar-symbol slice, normalize
\(U_{\mathcal P}\) and \(\Theta\) on the scalar summands to the geometric
parity and complex-conjugation lifts, respectively.  Such a normalization
always exists by multiplicity one and Schur's lemma.  Arbitrary sector-phase
choices yield the same quotient action after a compensating sector
rephasing back to this slice.
With \(s=(-1,-1,-1,+1)\), define the transported completions
\[
 (\widehat X^{\mathcal P})_a
 :=s_aU_{\mathcal P}\widehat X_aU_{\mathcal P}^{-1},
 \qquad
 (\widehat X^{\mathcal T})_a
 :=\Theta\widehat X_a\Theta^{-1}.
\]
Both are again compatible Hermitian equivariant coordinate vectors.  At the
selector-relevant level, parity or time reversal is a symmetry of the finite
realization class precisely when
\[
 [\widehat X^{\mathcal P}]=[\widehat X],
 \qquad\text{or}\qquad
 [\widehat X^{\mathcal T}]=[\widehat X],
\]
respectively, where \({\sim}\) is the unmarked sector-rephasing equivalence of
Paper~II-C.  Literal implementation on the chosen phase-marked representative
is the stronger condition
\[
 \widehat X^{\mathcal P}=\widehat X,
 \qquad\text{or}\qquad
 \widehat X^{\mathcal T}=\widehat X.
\]

At \(m=1\), with
\([\widehat X]\leftrightarrow(r_{1,0},r_{0,1},\varrho_{\rm sp})\), parity
acts on the realization quotient by
\[
 (r_{1,0},r_{0,1},\varrho_{\rm sp})
 \longmapsto
 (r_{0,1},r_{1,0},\varrho_{\rm sp}),
\]
so
\[
 \boxed{
 [\widehat X^{\mathcal P}]=[\widehat X]
 \Longleftrightarrow
 r_{1,0}=r_{0,1}.}
\]
Hence the level-one parity-fixed realization family is
\[
 \mathfrak M_{\rm coord}^{\mathcal P}(\mathcal H_1)/{\sim}
 \cong\mathbb R_{\ge0}^{2},
 \qquad (r,\varrho_{\rm sp}).
\]
At the same level, every completion class is time-reversal invariant:
\[
 \boxed{
 [\widehat X^{\mathcal T}]=[\widehat X]
 \quad\text{for every}\quad
 [\widehat X]\in
 \mathfrak M_{\rm coord}(\mathcal H_1)/{\sim}.}
\]
This class-level statement does not assert that an arbitrary prechosen
phase-marked representative is literally fixed by an arbitrary prechosen
anti-unitary lift.

At a finite auxiliary level, choose normalized forward intertwiners
\(T_{e,a}\) in a real coordinate-vector basis, with reverse frames
\(T_{\bar e,a}=T_{e,a}^*\), and write each nonzero edge block as
\(c_eT_{e,a}\).
Define
\[
 \Theta_vT_{e,a}\Theta_u^{-1}=\kappa_eT_{e,a},
 \qquad |\kappa_e|=1,\qquad e:u\to v.
\]
For an oriented cycle \(C\) of the nonzero-edge graph put
\(H_C=\prod_{e\in C}c_e\) and \(\kappa_C=\prod_{e\in C}\kappa_e\),
using adjoint reverse data for backward traversal.
Then
\[
 \boxed{
 [\widehat X^{\mathcal T}]=[\widehat X]
 \quad\Longleftrightarrow\quad
 H_C=\kappa_C\overline{H_C}
 \ \text{for every fundamental cycle}.}
\]
Time-reversal-real normalized intertwiner frames exist; in those frames
\(\kappa_e=1\) and the criterion is \(H_C\in\mathbb R\).
Negative real products are allowed.  A non-real raw product in an arbitrary
reference frame is not by itself an obstruction.
Lemma~\ref{p2d:lem:time-reversal-cycle-criterion} gives the frame conversion
and the constructive sufficiency proof.
\end{proposition}

\begin{proof}
Unitary conjugation preserves Hermiticity, while anti-unitary conjugation
preserves it after complex conjugating scalar coefficients.  The sign vector
\(s\) is the real \(O(4)\) parity action on the coordinate-vector index.
Conjugating the equivariance relations by \(U_{\mathcal P}\) or \(\Theta\)
therefore produces another compatible equivariant vector.  Equality in the
unmarked rephasing quotient is the selector-relevant fixed-point condition;
literal equality is the stronger representative-level implementation
condition.

On \(\mathcal H_1\), chiral exchange interchanges the sectors
\(V_{1,0}\) and \(V_{0,1}\), and therefore interchanges the reduced
magnitudes \(r_{1,0}\) and \(r_{0,1}\) while preserving
\(\varrho_{\rm sp}\).  This proves the parity action and the iff fixed-family
criterion.

For time reversal, anti-linearity acts on both factors of an edge block:
\[
 \Theta(c_eT_e)\Theta^{-1}
 =\overline{c_e}\,\kappa_eT_e.
\]
The edge magnitudes are unchanged.  A sector rephasing
\(c_e\mapsto z_v\overline{z_u}c_e\) relates the transformed and original
coefficients exactly when
\(\kappa_e\overline{c_e}/c_e=z_v/z_u\) on each nonzero edge.
At level one the graph is a forest, so these equations can always be solved
recursively and every quotient class is fixed, including the zero-edge
boundary strata.  At higher levels, solving on a spanning forest leaves
exactly the displayed fundamental-cycle compatibility equations.
The prescribed scalar backbone is fixed by the geometric scalar lift and
can be included in the spanning forest.
The details, including independence of the reference phases, are given in
Lemma~\ref{p2d:lem:time-reversal-cycle-criterion}.
\end{proof}

\begin{remark}[Class symmetry versus implementation gauge]
The quotient fixed-point statements above answer the physical selector
question.  A concrete calculation may additionally choose phases and a
specific implementing lift so that a representative itself is fixed.  Such a
choice is implementation data inside an already fixed realization class; it
must not be promoted to a second physical completion selector.
\end{remark}

\begin{corollary}[Joint completion/discrete-symmetry loci]
\label{p2d:cor:joint-completion-symmetry-loci}
Let
\[
 \mathfrak A_{\rm coord}
 =\{(r_L,r_R,\varrho):
       r_L,r_R,\varrho\in\mathbb R_{\geq0}\}
\]
denote the unmarked level-one completion atlas.
\begin{enumerate}[label=(\roman*),leftmargin=2.2em]
\item If parity is only required to act covariantly on the completion family,
the licensed output is the graph
\[
 \Gamma_{\mathcal P}
 =\{(c,\mathcal P c):c\in\mathfrak A_{\rm coord}\},
 \qquad
 \mathcal P(r_L,r_R,\varrho)=(r_R,r_L,\varrho).
\]
Projection onto \(c\) is a bijection, so covariance selects no completion.
\item If an independently generated claim requires the \emph{same} unmarked
completion class to be parity fixed, the joint zero-defect locus is the fibre
product of \(\Gamma_{\mathcal P}\) with the diagonal:
\[
 \boxed{\mathfrak Z_{\mathcal P}
 =\{(r,r,\varrho):r,\varrho\in\mathbb R_{\geq0}\}.}
\]
The crossed classes \(r_L\neq r_R\) survive the coordinate classification
separately but have nonzero joint relation defect for this same-completion
claim.
\item At level one, time reversal fixes every unmarked completion class.
Consequently its same-completion locus is
\[
 \boxed{\mathfrak Z_{\mathcal T}=\mathfrak A_{\rm coord}.}
\]
Requiring literal equality of a phase-marked representative with a chosen
anti-unitary lift is a stronger implementation relation and cannot be inserted
into this quotient-level locus.
\end{enumerate}
Intersecting the parity-fixed locus with Paper~II-C's least-closure graph
removes the two one-adjoint-channel branches: for \(r=0\) the least closure of
\(\mathcal E_9\) is \(\mathcal E_9\), while for \(r>0\) it is
\(\mathcal H_1\).  This is a joint compatibility statement, not a new carrier
or completion selector.
\end{corollary}

\begin{proof}
The proposition above gives the two actions on quotient coordinates.  A
generated readout \(c\mapsto\mathcal Pc\) produces its graph, whose first
projection has inverse \(c\mapsto(c,\mathcal Pc)\).  Imposing equality of its
two entries gives \(r_L=r_R\) and no condition on \(\varrho\).  Time reversal
acts trivially on the level-one quotient, giving the full atlas.  Finally,
Paper~II-C's coordinate/closure theorem adds one adjoint target for each
nonzero one of \(r_L,r_R\); on the diagonal they vanish together or occur
together, proving the last statement.
\end{proof}

\subparagraph{Anti-unitary phase redefinitions.}
A scalar phase cannot change an anti-unitary square; the elementary proof is
deferred to Paper~II-F, Section~\ref{p2d:app:discrete-supplements}.


\subparagraph{No geometric charge-conjugation theorem.}
No internal two-state exchange follows from the geometric kinematics.  An
explicit independently supplied doublet construction, included only to make
this boundary concrete, is deferred to
Paper~II-F, Section~\ref{p2d:app:discrete-supplements}.


\paragraph{Geometric discrete-symmetry summary}

\begin{table}[htbp]
\centering
\small
\begingroup
\renewcommand{\arraystretch}{1.3}
\begin{tabular}{|c|c|c|c|c|c|c|}
\hline
\textbf{Symmetry} & $M_{x_i}$ & $M_{x_4}$ & $\hat{P}_i$ & $\hat{J}_i$ & $\hat{\mathfrak h}$ & Type \\
\hline
$\mathcal{P}$ (parity) & $-$ & $+$ & $-$ & $+$ & --- & Unitary \\
\hline
$\mathcal{T}_0$ (geometric T) & $+$ & $+$ & $-$ & $-$ & $-$ & Anti-unitary \\
\hline
\end{tabular}
\caption[Geometric discrete-symmetry transformation rules]{Geometric discrete-symmetry transformation rules.}
\label{p2:tab:discrete-symmetries}
\TableNotes{%
Signs indicate eigenvalue under conjugation:
\(\mathcal{S}\hat{O}\mathcal{S}^{-1}=\pm\hat{O}\).
The coordinate signs are parent/abstract signs.  On a selected finite
completion they hold only under the fixed-point conditions of
Proposition~\ref{p2d:prop:finite-completion-discrete-symmetry-scope}.
The dash records that parity converts the selected left-fiber action into a
right-fiber action rather than acting intrinsically on the fixed left-fiber
grading: \(\mathcal P Q_{H,L}\mathcal P^{-1}=-Q_{H,R}\).
The independent exchange
\(\mathcal E_{\mathrm{int}}\) of
Proposition~\ref{p2:prop:independent-internal-doublet}
is not a geometric symmetry row: it fixes every displayed geometric operator
and reverses only the externally supplied \(Q_{\mathrm{int}}\).
}
\endgroup
\end{table}


\subsection{Covariant odd bridges and the even-word obstruction}
\label{p2d:sec:v352-odd-port}
\begin{corollary}[Canonical two-component bridge family]
\label{p2d:cor:v352-minimal-covariant-port}
On the canonical $G=E_0\oplus E_1\oplus K_+\oplus K_-$ and
$S=S_+\oplus S_-$ decomposition, the odd bridge representation has
exactly the irreducible content of
Theorem~\ref{p2b:thm:v264-odd-bridge-types}.  It contains no
trivial summand, and its smallest nonzero irreducible dimension is
two.  The full two-dimensional subspace
\[
 \mathcal B_+=\operatorname{Hom}(E_0,S_+)
 =\bigl\{\,T_+(s)=|s\rangle\langle e_0|:s\in S_+\,\bigr\}
 \cong (\tfrac12,0)
\]
is $\Spin(4)$-covariant and independent of any chosen nonzero
spinor $s$ (apart from the immaterial phase of the invariant line
$E_0$).  Adjoining the entire family and its adjoints to
$\mathcal A_D$ generates $M_{15}(\mathbb C)$, while no
nonzero individual member is $\Spin(4)$-invariant.
\end{corollary}
\begin{proof}
The cited Clebsch--Gordan decomposition contains only spinorial
half-integral types, whose smallest dimension is two.
Since $E_0$ is a trivial line, $\mathcal B_+$ transforms as
$S_+$.  Every nonzero $T_+(s)$ connects the existing
$M_{11}$ and $M_4$ blocks, so
Theorem~\ref{p2c:thm:v352-rankone-full-bridge} applies.
The invariant subspace of $S_+$ is zero.
\end{proof}

\begin{theorem}[Parity-even word recursion cannot generate an odd field]
\label{p2d:thm:v352-even-word-no-go}
With $\Theta(T)=\Gamma T\Gamma$ and
$\Gamma=I_{11}\oplus(-I_4)$,
$\operatorname{Fix}(\Theta)=M_{11}\oplus M_4$.
Any finite construction from central-even operator seeds using
sums, scalar products, matrix products, commutators, adjoints,
inverses when defined, and continuous functional calculus remains
in $\operatorname{Fix}(\Theta)$.  It cannot generate a nonzero
operator in
$\operatorname{Hom}(G,S)\oplus\operatorname{Hom}(S,G)$.
\end{theorem}
\begin{proof}
$\Theta$ is a unital involutive $*$-automorphism; its fixed-point
algebra is a unital $C^*$-subalgebra and is closed under all the
listed operations.  A nonzero off-block operator has
$\Theta(B)=-B$, so cannot belong to that algebra.
\end{proof}

\begin{remark}[Process parity versus odd observable]
\label{p2d:rem:v352-process-not-event}
A parity-preserving completely positive \emph{map} of the even
observable algebra can mix its central populations; this does not
make its possible odd Kraus operators accessible as field effects.
The conditional transfer generator in
Theorem~\ref{p4a:thm:v352-two-rate-transfer} explicitly shows the
difference.  It requires a separately generated and occurring
process; it is not an even-word proof of a new odd operator.
\end{remark}


\section{Hopf connections and charge-sector geometry}
\label{paper2d:sec:organized-hopf}

\subsection{Hopf grading and connection structure}
\subsubsection{Finite diagonal grading and Hopf fibration data}
\label{p2:subsec:finite-diagonal-grading}

For formulas, fix one representative maximal torus \(T\subset SU(2)\) from
the unmarked conjugacy class and use the \emph{left-fiber convention}
\[
        h_\alpha:=\exp(i\alpha\sigma_3),
        \qquad
        h_\alpha\cdot g:=h_\alpha g,
        \qquad
        S^2\simeq T\backslash SU(2),
        \qquad
        \alpha\sim\alpha+2\pi .
\]
The equivalent right-fiber presentation \(SU(2)/T\) is obtained by inversion
but is not mixed with this convention in the formulas below.  When the
\(4\pi\)-periodic \(SU(2)\) angle \(\psi=2\alpha\) is used, the same
principal generator is
\(\xi_0=2\partial_\psi\), and the metric-normalized Reeb field is
\(\xi=\xi_0/\Rmom\).

The chosen left-fiber Hopf torus and the diagonal Cartan axis used below are
presentation representatives inside conjugacy classes, not generated physical
directions.  On the finite non-diagonal \(\Spin(4)\)-module the Cartan
representative provides a useful weight atlas.  These two structures are
related only after the representative and its left/right identification have
been fixed; they are not the same operator in the Peter--Weyl convention used
for scalar Hopf charge below.

\begin{definition}[Finite diagonal spin weight and geometric Hopf charge]
\label{p2:def:finite-diagonal-spin-weight-hopf-charge}
On \(\mathcal H_m\), after fixing a diagonal Cartan axis, define
\[
 \widehat{\mathfrak s}:=L_3^++L_3^-=\widehat J_3/\hbar,\qquad
 \mathfrak s=m_L+m_R\in\tfrac12\ZZ,\qquad
 -m\le\mathfrak s\le m.
\]
Here \(m_L,m_R\) denote the weights of the existing algebraic generators
\(L_3^+,L_3^-\), respectively.

For the chosen geometric left-fibre action
\(g\mapsto h_\alpha g\), \(h_\alpha=\exp(i\alpha\sigma_3)\), set
\[
 \widehat{\mathfrak h}:=Q_{H,L}
 =-i\Rmom\xi=-2i\partial_\psi=-2L_3^-,
 \qquad \psi=2\alpha,\qquad \psi\sim\psi+4\pi.
\]
Equivalently,
\[
 U_L(\alpha)f(g)=f(h_\alpha g),\qquad
 U_L(\alpha)=e^{i\alpha Q_{H,L}}=e^{-2i\alpha L_3^-}.
\]
For \(D^j_{ab}(g)=\langle a|\rho_j(g)|b\rangle\),
\[
 D^j_{ab}(h_\alpha g)=e^{2ia\alpha}D^j_{ab}(g),\qquad
 D^j_{ab}(gh_\alpha)=e^{2ib\alpha}D^j_{ab}(g).
\]
Consequently
\[
 \boxed{m_L=b,\qquad m_R=-a,\qquad q=q_L=-2m_R=2a,}
 \qquad q_R=2m_L=2b.
\]
The diagonal weight on these scalar matrix coefficients is
\(\mathfrak s=b-a\).  It generates the diagonal subgroup, represented
geometrically by conjugation, and is distinct from the chosen left
Hopf-fibre charge.  The integer character is \(e^{iq\alpha}\), or
\(e^{iq\psi/2}\) in the \(4\pi\)-periodic parameter.
\end{definition}

\begin{proposition}[Coordinate Hopf-charge steps and the superselection boundary]
\label{p2d:prop:coordinate-hopf-charge-steps}
Fix the oriented left-fibre representative of
Definition~\ref{p2:def:finite-diagonal-spin-weight-hopf-charge} and write
\(Q_H:=Q_{H,L}=-2L_3^-\).  The complexification of any equivariant
coordinate tuple of type \((\tfrac12,\tfrac12)\) admits components
\(X_{\pm,\alpha}\), \(\alpha=1,2\), such that
\[
 \boxed{[Q_H,X_{\pm,\alpha}]=\pm X_{\pm,\alpha},\qquad
 P_qX_{\pm,\alpha}P_r=0\quad\text{unless }q=r\pm1,}
\]
where \(P_q\) are the spectral projectors of \(Q_H\) on the retained
carrier.  The same identities hold under every equivariant finite
compression.  These coordinate components preserve the diagonal central
parity \(\Gamma_c\) of Theorem~\ref{p2:thm:diagonal-center-grading}.
Thus Hopf grading and central parity impose different selection rules.

In particular, the nonzero intrinsic coordinate tuple does not lie in
\(\{Q_H\}'\).  A Hopf-weight decomposition alone therefore does not
establish charge superselection for an accessible algebra containing that
tuple.  The charge-invariant accessible algebra and dynamical stability
required by Paper~IV, Corollary~\ref{p4:cor:hopf-reeb-superselection},
remain the appropriate additional conditions.
\end{proposition}

\begin{proof}
The right spin-\(\tfrac12\) factor of the coordinate representation has
\(-2L_3^-\)-weights \(\pm1\); its other spin-\(\tfrac12\) factor
supplies the index \(\alpha\).  Differentiating coordinate covariance
gives the displayed commutator.  Multiplication on the left and right by
\(P_q,P_r\) gives
\((q-r\mp1)P_qX_{\pm,\alpha}P_r=0\), proving the selection rule.
An equivariant compression projector commutes with \(Q_H\), so compression
preserves the commutator exactly.  Central-evenness is the existing
coordinate-type calculation in Theorem~\ref{p2:thm:diagonal-center-grading}.
If every component commuted with \(Q_H\), the commutator identities would
force the entire tuple to vanish.  The intrinsic tuple is nonzero, giving
the final claim.
\end{proof}


\begin{remark}[Hopf fibre charge, diagonal weight, and transverse-shell degree]
\label{p2d:rem:hopf-charge-versus-transverse-shell-degree}
Three labels that live on the same marked round geometry have different
representation-theoretic meanings and are not identified.  The geometric
Hopf charge \(q\) above is the character of the selected one-sided fibre
\(U(1)\).  The number \(\mathfrak s\) is a Cartan weight for the diagonal
rotation subgroup acting by conjugation.  By contrast, the transverse-shell
label \(\ell\) in the Paper~I-A/II-C shell-symbol theorem
\[
 E_k^x\simeq\bigoplus_{\ell=0}^{k}\mathscr Y_\ell(S_H^2)
\]
labels the irreducible diagonal-rotation \(SU(2)\) type itself; within one
such \(V_\ell\), its Cartan weights range from \(-\ell\) to \(+\ell\).
Accordingly, neither \(\ell=q\) nor \(\ell=\mathfrak s\) is generated by the
framework.  The shell-symbol map is also not ordinary Hopf pullback:
ordinary fibre-basic pullback obeys
\(\pi^*\mathscr Y_\ell(S_H^2)\subset E_{2\ell}(S^3)\), whereas the
Hopf-marked diagonal symbolization reconstructs the full shell \(E_k\) as the
complete window \(0\le\ell\le k\).  This firewall prevents the transverse
frontier rank from being reinterpreted as a Hopf-charge count.
\end{remark}

\begin{proposition}[Diagonal spin-weight sector decomposition]
\label{p2:prop:diagonal-spin-weight-sector-decomposition}
The Hilbert space decomposes into finite diagonal spin-weight sectors:
\begin{equation}
\mathcal H_m=\bigoplus_{\mathfrak s\in \frac12\ZZ,\ |\mathfrak s|\le m}\mathcal H_m^{(\mathfrak s)},
\qquad
\mathcal H_m^{(\mathfrak s)}
=
\{\psi\in\mathcal H_m:\hat{\mathfrak s}\psi=\mathfrak s\psi\}.
\end{equation}
\end{proposition}

\begin{proof}
Since \(\hat{\mathfrak s}\) is Hermitian, \(\mathcal H_m\) decomposes as the orthogonal direct sum of its eigenspaces:
\[
\mathcal H_m=\bigoplus_{\mathfrak s}\ker(\hat{\mathfrak s}-\mathfrak s\mathbf 1).
\]
This is exactly the stated decomposition.
\end{proof}

For \(m=1\), the finite diagonal spin weight \(\mathfrak s=m_L+m_R\) has sectors
\begin{center}
\begin{tabular}{|c|c|c|c|c|c|}
\hline
\(\mathfrak s\) & \(-1\) & \(-\frac12\) & \(0\) & \(\frac12\) & \(1\) \\
\hline
\(\dim \mathcal H_1^{(\mathfrak s)}\) & 3 & 2 & 5 & 2 & 3 \\
\hline
\end{tabular}
\end{center}
These are diagonal-spin sectors, not scalar left-fiber Hopf-charge
sectors.  The geometric scalar charge sectors are defined separately in
Definition~\ref{p2:def:hopf-charge}.

\subparagraph{Protected diagonal-weight outcome.}
At \(m=1\), the \(\mathfrak s=\pm\tfrac12\) sectors have complex
dimension two each and their sum is the protected four-dimensional spinor
submodule
\[
 \boxed{
 \mathcal H_1^{(-1/2)}\oplus\mathcal H_1^{(+1/2)}
 =V_{\frac12,0}\oplus V_{0,\frac12}.}
\]
Although the individual weight projectors are written relative to a chosen
Cartan axis, this direct sum is the invariant protected spinor carrier and is
independent of that representative.  The projector derivation and detailed
isolation mechanics are in
Paper~II-F, Section~\ref{p2d:app:diagonal-weight-isolation}.

\subsubsection{Momentum-space Hopf-sector connection data}
\label{p2:sec:momentum-hopfions}

The constituted momentum parent is the round group
\(S^3(\Rmom)\simeq SU(2)\), independently of any finite truncation level.
Its maximal tori supply compatible Hopf fibrations, but the unmarked parent
selects no particular torus or axis.

\begin{proposition}[Hopf-fibration existence, conjugacy, and nonselection]
\label{p2:prop:hopf-class-existence-nonselection}
Let \(G\simeq SU(2)\) be the constituted nonabelian transfer parent.  Every
maximal torus \(T\subset G\) satisfies
\[
        T\simeq U(1),
        \qquad
        T\backslash G\simeq S^2,
\]
and the left-orbit quotient map \(G\to T\backslash G\) is a Hopf fibration
(equivalently one may use \(G/T\) after inversion).  All maximal tori are
conjugate, so these fibrations form one equivalence class after a handedness
has been fixed, but no particular \(T\) is invariant under all inner
automorphisms of \(G\).  The parent geometry therefore determines the Hopf
conjugacy class and its choice-independent invariants, but no preferred embedded
fiber circle.
\end{proposition}

\begin{proof}
Existence and \(T\backslash G\simeq S^2\) are the rank-one homogeneous-space
description \(U(1)\backslash SU(2)\simeq S^2\), equivalent under inversion to
\(SU(2)/U(1)\simeq S^2\).  The conjugacy theorem for maximal tori
shows that any two choices are related by an inner automorphism.  If one torus
were selected by the unmarked parent, it would be characteristic.  This is
impossible because conjugation acts transitively and nontrivially on the
maximal tori of \(SU(2)\).  Hence only the orbit of choices is endogenous.
\end{proof}

\begin{remark}[Hopf fiber versus scalar parent]
\label{p2:rem:hopf-fiber-versus-scalar-parent}
The \(U(1)\) appearing as a maximal torus in this proposition is an internal
fiber subgroup of the already derived \(SU(2)\) parent.  It is not the
independent scalar-succession parent \(S^1_A\) classified earlier.  No
identification between them is made in Paper~II; such an identification would
require an additional generated comparison map in the later scalar
constitution.
\end{remark}

\begin{theorem}[Chern unit of a Hopf representative on the simply connected parent]
\label{p2:thm:forced-hopf-chern-unit}
Let
\[
        U(1)\longrightarrow P\longrightarrow S^2
\]
be any principal-circle representative whose total space is the constituted
simply connected parent \(P\simeq S^3\).  Then its first Chern class satisfies
\[
|c_1(P)|=1.
\]
After the Hopf/base orientation convention is fixed, the chosen oriented
representative may be written
\[
c_1(P)=+1.
\]
\end{theorem}

\begin{proof}
Principal \(U(1)\)-bundles over \(S^2\) are classified by
\(c_1(P)=k\in\mathbb Z\).  The homotopy exact sequence contains
\[
 \pi_2(S^2)\cong\mathbb Z
 \xrightarrow{\ \times k\ }
 \pi_1(U(1))\cong\mathbb Z
 \longrightarrow\pi_1(P)\longrightarrow0.
\]
Consequently \(\pi_1(P)\cong\mathbb Z/k\mathbb Z\), with the \(k=0\)
case interpreted as \(\mathbb Z\).  Since \(P\simeq S^3\) is simply
connected, \(|k|=1\).  The two signs are exchanged by reversing either the fibre orientation or
the base orientation, with the other held fixed; simultaneous reversal of
both preserves \(k\).  Indeed,
\(k=\langle c_1(P),[S^2]\rangle\), and these two reversals negate
\(c_1(P)\) and \([S^2]\), respectively.  With the Paper-I-A oriented
representative, \(k=+1\).
\end{proof}

\begin{remark}[Principal Chern unit does not select a matter charge]
The oriented principal-bundle result
\(c_1(P)=+1\) fixes the primitive Hopf bundle, not a physical associated
character sector.  In particular,
\[
 \boxed{c_1(P)=+1\not\Longrightarrow\text{physical Hopf charge }q=+1.}
\]
The complete associated-character atlas is \(\{L^q\}_{q\in\mathbb Z}\), with
\(c_1(L^q)=q\).  A later role may derive \(q=+1\) only if it independently
specifies the primitive oriented associated line \(L\) (or an equivalent
generated charge signature).
\end{remark}

Fix the chosen left-fiber Paper-I-A representative of this conjugacy class.
Let
\[
\eta_0:=\eta/\Rmom
\]
be its dimensionless principal \(U(1)\)-connection.  The bundle is
\[
\Sone\hookrightarrow S^3(\Rmom)\longrightarrow S^2.
\]
Then
\begin{equation}
\label{p2:eq:eta0}
\oint_{S^1}\eta_0=2\pi,
\qquad
 d\eta_0=\pi^*\omega,
\qquad
\frac{1}{2\pi}\int_{S^2}\omega = 1,
\end{equation}
so \(c_1(P)=1\).  We use \(\eta_0\) as a representative of the
homogeneous Reeb-normalized principal-connection gauge class.  Gauge-related
principal forms are the same unmarked connection realization; the displayed
representative is not an additional physical selector.  Every invariant
conclusion below is unchanged under conjugating the maximal torus; axis- or
handedness-dependent formulas are convention-relative.


\subparagraph{Projected doubled-orbit diagnostic.}
Half-integral projected Reeb weights define only local/double-cover lift data,
not a global Berry connection on the principal \(2\pi\)-orbit.  The precise
definition is deferred to Paper~II-F, Section~\ref{p2d:app:berry-technical}.

\subparagraph{Berry technical supplement.}
Global trivialization, orbit-weight integrality, and homogeneous uniqueness of
the Reeb-normalized Berry connection are proved in
Paper~II-F, Section~\ref{p2d:app:berry-technical}.  They are not premises of the main
proposition below.  The appendix packages their conditional global-Berry
consequence separately, after stating all of its hypotheses.
Corollary~\ref{p2d:cor:fibre-basic-berry-cs} extends the quadratic Berry
integral to genuinely fibre-basic ray maps without requiring homogeneity;
the stronger homogeneous connection identity retains its own hypotheses.

\begin{definition}[Abelian Chern--Simons diagnostic on the total space]
\label{p2:def:cs-charge}
For a real one-form \(A\) on \(S^3\), define the quadratic abelian
Chern--Simons diagnostic
\begin{equation}
\label{p2:eq:cs}
Q_{\CS}[A] \;:=\; \frac{1}{4\pi^2}\int_{S^3} A\wedge dA.
\end{equation}
When \(A=\eta_0\), this is evaluated on the principal Hopf connection.  When
\(A=q\eta_0\) with \(q\ne1\), it is evaluated on the pullback of the
connection induced on the associated character line; \(q\eta_0\) is not a
new principal connection on the same circle bundle because its vertical
normalization is \(q\), not \(1\).
\end{definition}


\begin{proposition}[Principal and associated-character normalizations]
\label{p2:prop:hopfion-charge}
Among the homogeneous scalar multiples \(s\eta_0\), the form \(\eta_0\) is
the unique one with the principal normalization
\((s\eta_0)(\xi_0)=1\).  It represents the homogeneous Reeb-normalized principal-connection gauge
class and satisfies
\(
Q_{\CS}[\eta_0]=1.
\)
For the integer character
\(\rho_q(e^{i\alpha})=e^{iq\alpha}\), let
\[
        L^q:=S^3\times_{\rho_q}\mathbb C\longrightarrow S^2.
\]
The induced associated-line connection has, after pullback to \(S^3\),
\begin{equation}
\label{p2:eq:Aq0-q0eta0}
A_q=q\,\eta_0-d\chi,
\qquad q\in\ZZ,
\end{equation}
and its curvature and first Chern class are
\[
        F_q=q\,\omega,
        \qquad
        c_1(L^q)=\frac{[F_q]}{2\pi}=q.
\]
Its total-space quadratic diagnostic is
\(Q_{\CS}[A_q]=q^2\).

The arbitrary-real one-form diagnostic and the conditional global homogeneous
Berry normalization are supplementary and are recorded in
Paper~II-F, Section~\ref{p2d:app:berry-technical} after their exact hypotheses.
\end{proposition}

\begin{proof}
For a homogeneous scalar multiple \(s\eta_0\), the principal vertical
normalization is
\((s\eta_0)(\xi_0)=s\), so it forces \(s=1\).  This uniqueness statement is
only within the homogeneous scalar-multiple family; it does not exclude
distinct gauge-related principal connections such as
\(\eta_0-\pi^*d\bar\chi\) for a smooth real function
\(\bar\chi\) on \(S^2\) with \(d\bar\chi\not\equiv0\).
Using \eqref{p2:eq:eta0} and fiber integration directly in the main proof gives
\[
 \int_{S^3}\eta_0\wedge d\eta_0
 =\int_{S^2}\!\left(\oint_{S^1}\eta_0\right)\omega
 =(2\pi)(2\pi)=4\pi^2,
\]
and hence \(Q_{\CS}[\eta_0]=1\) in the convention
\eqref{p2:eq:cs}.  Applying the character differential
\(d\rho_q:q\) to the principal connection and curvature gives
\(A_q=q\eta_0-d\chi\) and \(F_q=q\omega\); hence
\(c_1(L^q)=q[\omega/(2\pi)]=q\).  For
\(A_s=s\eta_0-d\chi\), the exact gauge term integrates to zero on closed
\(S^3\), so quadratic scaling gives
\(Q_{\CS}[A_s]=s^2\).  Applying this with \(s=q\) proves the associated-line
claim.
\end{proof}

\begin{corollary}[Associated-character readout graph and nonselection]
\label{p2d:cor:associated-character-readout-graph}
Modulo gauge on the pulled-back connection, the complete generated
associated-character package is the graph
\[
 \boxed{
 \Gamma_{\rm char}
 =
 \left\{
 \left(q,\,[q\eta_0],\,q\omega,\,q,\,q^2\right):
 q\in\mathbb Z
 \right\},}
\]
where the entries after \(q\) are respectively the pulled-back connection
class, curvature, first Chern number, and total-space quadratic diagnostic.
Projection \(\Gamma_{\rm char}\to\mathbb Z\) is a bijection.  Thus curvature,
Chern number, and Chern--Simons readouts co-realize with the character and
cannot by themselves select \(q\).  An independently generated associated-line
Chern requirement \(c_1(L^q)=n\) instead forms a fibre product and fixes
\(q=n\); the principal-bundle statement \(c_1(P)=1\) is differently typed and
does not supply that relation.
\end{corollary}

\begin{proof}
Proposition~\ref{p2:prop:hopfion-charge} gives every component as a function of
\(q\), and the first component recovers \(q\), proving that the package is a
graph isomorphic to the integer atlas.  Intersecting it with the
associated-line equation \(c_1=n\) gives \(q=n\).  The principal Chern class
belongs to the principal bundle rather than to a chosen associated line, so
the typed cross-relation firewall forbids replacing the latter equation by
\(c_1(P)=1\).
\end{proof}

\begin{remark}[Normalization and interpretation boundary]
\label{p2:rem:hopf-normalization-conversion}
The integer character \(q\) and the projected ray weight
\(r\in\frac12\mathbb Z\) are independent labels.  The associated-line Chern
class is linear and oriented, whereas a total-space Chern--Simons diagnostic
is quadratic and loses the sign.  A projected half-integral \(r\) is
double-cover/local lift data; only an integer \(r\) can occur in the global
homogeneous Berry form under the present \(2\pi\) principal normalization.
The chosen form \(B_r=r\eta_0\) exists for arbitrary real \(r\), but its
quadratic diagnostic must not be promoted to a global Berry invariant without
a separate global ray-and-connection construction.  No physical charge
identification or superselection statement follows in Paper~II.
\end{remark}


\section{Topological and reduced-phase-space spectra}
\label{paper2d:sec:organized-topology}

\subsection{Reduced phase-space and topological spectral consequences}
\subsubsection{Geometric Prequantization of the Reduced Phase Space}
\label{p2:sec:prequantization}

We use the Hopf/contact normalization established in Paper~I-A \cite{Kugelmann2026a}.  Thus the Hopf fibration \(\Sone \hookrightarrow S^3 \xrightarrow{\pi} S^2\) realizes $S^3(\Rmom)$ as a principal $\mathrm{U}(1)$-bundle over the symplectic manifold $(S^2(\Rmom/2),\omega)$. Let $\vol_{S^2(\Rmom/2)}$ denote the Riemannian area form, so $\int_{S^2(\Rmom/2)}\vol_{S^2(\Rmom/2)}=\pi \Rmom^2$. We take the normalized symplectic form
\[
\omega := \frac{2}{\Rmom^2}\,\vol_{S^2(\Rmom/2)},\qquad \int_{S^2(\Rmom/2)}\omega = 2\pi,
\]
so that the associated complex line bundle $L\to S^2(\Rmom/2)$ has $c_1(L)=[\omega]/(2\pi)=1$.


\begin{proposition}[Hopf connection and prequantum bundle]
\label{p2:prop:hopf-connection-prequantum-bundle}
Let $\eta_0$ denote the \emph{dimensionless} principal $\mathrm{U}(1)$-connection 1-form on the Hopf bundle $S^3\to S^2$, with principal generator $\xi_0$ normalized by $\eta_0(\xi_0)=1$, and set
\[
\eta := \Rmom\,\eta_0,
\qquad
\xi:=\Rmom^{-1}\xi_0.
\]
The metric-scaled contact form is normalized by $\eta(\xi)=1$.
Then the curvature of the principal connection is \cite{blair2010contact,boyer2008sasakian}
\[
d\eta_0 = \pi^* \omega,
\]
so \((S^3, \eta_0)\) is a prequantum circle bundle over \((S^2, \omega)\) with integer Chern class \(c_1 = 1\).
\end{proposition}

\begin{proof}
Paper~I-A proves this degree-one Boothby--Wang normalization
\cite{Kugelmann2026a}.  The displayed equations translate its physical and
dimensionless connections into Paper-II notation.
\end{proof}

\begin{corollary}[Generated principal/metric Hopf normalization graph]
\label{p2d:cor:generated-hopf-normalization-graph}
Fix the occurring round radius \(\Rmom>0\), the \(2\pi\)-periodic principal
generator \(\xi_0\), and its principal connection \(\eta_0\) with
\(\eta_0(\xi_0)=1\).  The metric normalization is the generated transport
graph
\[
 \boxed{
 \xi=\Rmom^{-1}\xi_0,\qquad
 \eta=\Rmom\eta_0,\qquad
 \eta(\xi)=1,\qquad
 d\eta=\Rmom\,\pi^*\omega.}
\]
The metric generator, metric contact form, and their curvature are
nonredundant outputs of the same normalization dictionary and co-realize with
the principal data.  They are not independent members requiring separate
selectors.  Conversely, this graph does not select \(\Rmom\), a maximal-torus
representative, an associated character, or a physical matter charge.
\end{corollary}

\begin{proof}
The first two identities are the declared principal-to-metric transport.
Substitution gives
\(\eta(\xi)=(\Rmom\eta_0)(\Rmom^{-1}\xi_0)=1\), and constancy of
\(\Rmom\) on the round host gives
\(d\eta=\Rmom\,d\eta_0=\Rmom\,\pi^*\omega\).
Omitting any declared output gives a closure defect once the dictionary is
licensed; multiplying either output by an additional free scalar changes the
generated normalization relation.  The remaining quantities occupy distinct
owner or associated-line slots and therefore are not fixed by this graph.
\end{proof}


\subparagraph{Supplementary quantization dictionary.}
The \(\mathbb{CP}^1\) Borel--Weil/geometric-quantization dictionary is useful
for interpretation but is not used in the positive spectral chain; it is
recorded in Paper~II-F, Section~\ref{p2d:app:supplementary-geometry}.

\subparagraph{Supplementary geodesic periodicity.}
Round-parent geodesic closure and its arclength-period interpretation are
recorded in Paper~II-F, Section~\ref{p2d:app:supplementary-geometry}; they do not supply
a physical-time period.

\subsubsection{Topological Spectral Gaps}
\label{p2:sec:kinematic-spectral-gaps}

The non-trivial topology of the Hopf fibration ($c_1 = 1$) induces a positive
\textbf{spectral gap} for charged Hopf sectors.  Topology fixes the
obstruction to a zero mode; the exact numerical gap also uses the round
metric, Laplacian normalization, and the present Hopf-charge convention.  No
energy or mass interpretation is part of this statement.  In particular,
these are lower bounds in the spectra of specified positive spatial
operators; conversion to an energy threshold requires the additional
later finite-dispersion domain and dispersion hypotheses stated in
Remark~\ref{p2d:rem:dispersion-speed-boundary}.

\paragraph{Charge Sectors on the Hopf Bundle}

\begin{definition}[Integer Hopf charge in the $4\pi$ fiber convention]
\label{p2:def:hopf-charge}
Let $\psi\sim\psi+4\pi$ be the standard Hopf fiber angle on
$S^3\cong SU(2)$, so that the chosen left-fiber Hopf action is
\[
 g\longmapsto h_\theta g,
 \qquad
 h_\theta:=\exp(i\theta\sigma_3/2),
 \qquad \theta\sim\theta+4\pi.
\]
The Reeb vector field $\xi = \frac{2}{\Rmom}\partial_\psi$ generates this fiber action on $S^3(\Rmom)$. A field $\phi \in C^\infty(S^3)$ has \textbf{integer Hopf charge} $q \in \ZZ$ if
\begin{equation}
\partial_\psi \phi = i\frac{q}{2}\,\phi,
\qquad
\xi(\phi)=\frac{i q}{\Rmom}\,\phi,
\qquad
\phi(h_\theta g)=e^{i(q/2)\theta}\phi(g).
\label{p2:eq:hopf-charge-convention}
\end{equation}
Equivalently, if one uses the $2\pi$-periodic principal-$U(1)$ parameter $\alpha=\theta/2$, the character is $e^{iq\alpha}$.  The Hilbert space decomposes as $L^2(S^3) = \bigoplus_{q \in \ZZ} \mathcal{V}_q$, where $\mathcal{V}_q$ consists of charge-$q$ fields.
\end{definition}

\begin{remark}[Geometric interpretation and notation]
Charge-$q$ fields are sections of the line bundle $L^q \to S^2$, where
$L$ is the degree-one bundle associated to the chosen left-fiber Hopf
representative. The case $q = 0$ corresponds to functions on the base $S^2$
pulled back to $S^3$; $q \neq 0$ corresponds to genuinely ``twisted'' fields
that cannot be defined on $S^2$ alone.  Throughout this section, $q$ denotes
the integer geometric Hopf charge and $q/2$ is the exponent of its
left-fiber character in the \(4\pi\) parameter.  This exponent is not
identified with the projected Berry label \(r\), and both are distinct from
the finite diagonal spin weight $\mathfrak s$ used in the preceding
finite-module projector notation.
\end{remark}

\begin{remark}[Charge-sector selector boundary]
The decomposition
\(L^2(S^3)=\bigoplus_{q\in\mathbb Z}\mathcal V_q\) is a generated
mathematical charge atlas.  Paper~II-D does not select one \(q\) as an
actually occurring matter sector.  A physical charge channel belongs to the
later role that supplies an independently generated charge or associated-line
criterion; if that criterion is the primitive oriented line \(L\), its
signature fixes \(q=+1\) there rather than retroactively in this paper.
\end{remark}

\begin{corollary}[Joint charge/time-reversal and parity covariance]
\label{p2d:cor:joint-hopf-charge-discrete-symmetry}
For the chosen left-fibre charge atlas
\(\{\mathcal V_q\}_{q\in\mathbb Z}\), time reversal generates
\[
 \mathcal T:\mathcal V_q\longrightarrow\mathcal V_{-q}.
\]
Hence the complete time-reversal-covariant orbit of a nonzero charge is the
pair \(\mathcal V_q\oplus\mathcal V_{-q}\); retaining both generated
constituents is co-realization, not a choice between them.  The joint locus for
a claim that one \emph{single} charge sector is itself time-reversal fixed is
\[
 \boxed{\{q\in\mathbb Z:q=-q\}=\{0\}.}
\]
Consequently the simultaneous requirements \(q\neq0\) and one
time-reversal-fixed charge sector have empty joint zero-defect locus.  GP3
cannot repair that empty locus; a nonzero time-reversal-covariant package must
contain the opposite-charge pair (or additional typed sectors supplied by a
later constructor).

Spatial parity transports the chosen left-fibre action with
\(\mathcal P Q_{H,L}\mathcal P^{-1}=-Q_{H,R}\), so
\(\mathcal P:\mathcal V_q^L\to\mathcal V_{-q}^R\).
It therefore generates a left/right covariance orbit, not
an endomorphism of one fixed left-fibre sector and not a selector of a charge
sign.  A parity-invariant charged claim must supply the corresponding
left/right pair or another independently generated bridge identifying the
relevant typed actions.
\end{corollary}

\begin{proof}
Theorem~\ref{p2:thm:time-reversal-action} gives \(q\mapsto-q\).
The orbit has one member exactly when \(q=-q\) in \(\mathbb Z\), which forces
\(q=0\); its intersection with \(\mathbb Z\setminus\{0\}\) is empty.
The parity clause is the left-to-right action proved in
Theorem~\ref{p2:thm:parity-action}.  Since the target fibre action is differently
typed, the typed cross-relation firewall forbids treating that transport as
charge-sign selection inside the original left-fibre atlas.
\end{proof}

\paragraph{Spectrum in Each Charge Sector}

\begin{lemma}[Integer Hopf-charge selection rule]
\label{p2:lem:hopf-charge-selection-rule}
In the scalar Peter--Weyl sector with matrix coefficients \(D^j_{ab}\), the
chosen left-fiber Hopf action gives integer Hopf charge
\[
q=2a.
\]
Writing \(k=2j\), the charge-\(q\) sector occurs at level \(k\) iff
\[
k\ge |q|,
\qquad
k\equiv q \pmod 2.
\]
For fixed \(q\) and \(k\), the remaining right matrix-coefficient index gives
degeneracy \(k+1\).
\end{lemma}

\begin{proof}
By Peter--Weyl, \(L^2(S^3)\cong L^2(SU(2))\) decomposes into matrix
coefficients \(D^j_{ab}\), where \(a,b=-j,-j+1,\ldots,j\).  With the
convention \(D^j_{ab}(g)=\langle a|\rho_j(g)|b\rangle\), left multiplication by
\[
h_\theta=\exp(i\theta\sigma_3/2),
\qquad \theta\sim\theta+4\pi,
\]
gives
\[
D^j_{ab}(h_\theta g)=e^{ia\theta}D^j_{ab}(g).
\]
Comparing with Definition~\ref{p2:def:hopf-charge}, namely
\(\phi(h_\theta g)=e^{i(q/2)\theta}\phi(g)\), gives \(q=2a\).  Since
\(a\) ranges from \(-j\) to \(j\) in integer steps and \(k=2j\), a fixed
charge \(q\) can occur exactly when
\[
q\in\{-k,-k+2,\ldots,k-2,k\},
\]
which is equivalent to \(k\ge |q|\) and \(k\equiv q\pmod2\).  Once
\(q=2a\) is fixed, the left index is fixed and the right index \(b\) has
\(2j+1=k+1\) possible values.  This gives the stated fixed-charge degeneracy.
\end{proof}

\begin{theorem}[Spectrum decomposition by integer Hopf charge]
\label{p2:thm:spectrum-by-charge}
The nonnegative scalar Laplacian $\Delta_{S^3}$ preserves each integer charge sector $\mathcal{V}_q$. The spectrum of $\Delta_{S^3}|_{\mathcal{V}_q}$ is
\begin{equation}
\mathrm{Spec}(\Delta_{S^3}|_{\mathcal{V}_q}) = \left\{ \frac{k(k+2)}{\Rmom^2} : k \geq |q|, \; k - |q| \in 2\ZZ \right\},
\label{p2:eq:spectrum-charge-n}
\end{equation}
with degeneracy $(k+1)$ in the fixed charge sector.  Here $k=2j$ is the
scalar Peter--Weyl level for the diagonal sector $(j,j)$; this statement is
about the scalar Laplacian and should not be confused with the full
non-diagonal $\Spin(4)$ module Casimir used for the declared finite package.
\end{theorem}

\begin{proof}
The spectrum, parity restriction, and fixed-charge multiplicity are the
Paper-I-A scalar Hopf-charge spectrum theorem~\cite{Kugelmann2026a}.
Paper~II imports it only for kinematic sector classification.
\end{proof}

\paragraph{The Topological Spectral-Gap Theorem}

\begin{theorem}[Exact retained Hopf-charge kinematic spectral gap]
\label{p2:thm:hopf-charge-spectral-gap}
For Hopf charge $q \neq 0$, the spectrum of the nonnegative scalar Laplacian $\Delta_{S^3}|_{\mathcal{V}_q}$ is bounded below by a strictly positive constant:
\begin{equation}
\lambda_{\min}(q) = \frac{|q|(|q|+2)}{\Rmom^2} > 0.
\label{p2:eq:hopf-charge-spectral-gap}
\end{equation}
Equivalently, the charged sector has a purely kinematic gap in the positive spatial operator.  In a finite scalar-diagonal truncation \(S^3_m\), this continuum-sector statement applies only to retained nonempty sectors satisfying \(|q|\le m\) and the charged-mode parity condition \(k\equiv |q|\pmod 2\).
\end{theorem}

\begin{proof}
Paper~I-A proves both the geometric lower bound and its associated-line-bundle
realization~\cite{Kugelmann2026a}.  Imposing the Paper-II retained-sector
condition yields the stated kinematic gap.
\end{proof}


\subparagraph{Monopole-harmonic interpretation.}
The exact charged-sector spectrum above admits the standard monopole-harmonic
interpretation; the comparison is supplementary and is deferred to
Paper~II-F, Section~\ref{p2d:app:supplementary-geometry}.

\paragraph{The Bochner-Weitzenböck Perspective}

\begin{theorem}[Lichnerowicz Formula on $S^3$]
\label{p2:thm:lichnerowicz}
For any spinor $\psi$ on the round 3-sphere $S^3(\Rmom)$,
\begin{equation}
D_{S^3}^{\,2}\psi = \nabla^*\nabla\psi + \frac{\scal}{4}\psi = \nabla^*\nabla\psi + \frac{3}{2\Rmom^2}\psi.
\end{equation}
\end{theorem}

\begin{proof}
This is the round-\(S^3\) Lichnerowicz identity recorded in
Paper~I-A~\cite{Kugelmann2026a}.  It is repeated here only to state the
kinematic no-zero consequence.
\end{proof}

\begin{corollary}[Lichnerowicz no-zero estimate]
\label{p2:cor:fermion-gap}
Every Dirac eigenvalue $\lambda$ satisfies
\[
\lambda^2\ge \frac{3}{2\Rmom^2}.
\]
In particular, the Lichnerowicz identity rules out harmonic spinors on the positively curved round $S^3$. The sharper exact gap is obtained from the round-sphere spectrum in Theorem~\ref{p2:thm:dirac-spectrum}.
\end{corollary}

\begin{proof}
From Theorem~\ref{p2:thm:lichnerowicz}, $D_{S^3}^{\,2}=\nabla^*\nabla+\frac{3}{2\Rmom^2}$.
Since $\nabla^*\nabla$ is non-negative, every Dirac eigenvalue $\lambda$ satisfies the displayed bound.
\end{proof}


\paragraph{Dirac Operator and Spinor Spectral Gap}

\begin{theorem}[Dirac Spectrum on $S^3$]
\label{p2:thm:dirac-spectrum}
The Dirac operator $D_{S^3}$ on $S^3(\Rmom)$ with the unique spin structure has spectrum:
\begin{equation}
\mathrm{Spec}(D_{S^3}) = \left\{ \pm\frac{k + 3/2}{\Rmom} : k = 0, 1, 2, \ldots \right\},
\end{equation}
with multiplicity $(k+1)(k+2)$ for each individual eigenvalue $+(k+3/2)/\Rmom$ and $-(k+3/2)/\Rmom$. Counting both signs together gives $2(k+1)(k+2)$.

The \textbf{spinor Dirac spectral gap} is:
\begin{equation}
|{D_{S^3}}|_{\min} = \frac{3}{2\Rmom} = \frac{3\Punit}{2\hbar}.
\end{equation}
\end{theorem}

\begin{proof}
The eigenvalues, per-sign multiplicities, and exact first gap are the spinor
part of the Paper-I-A round-\(S^3\) harmonic package
\cite{Kugelmann2026a}.  Paper~II imports this only as a kinematic continuum
spectrum.
\end{proof}

\begin{corollary}[No Dirac zero modes]
\label{p2:cor:no-dirac-zero}
On compact \(S^3\) momentum space, there are no zero modes of the Dirac
operator.  The kinematic Dirac spectrum has the positive operator gap
\begin{equation}
|D_{S^3}|_{\min}=\frac{3}{2\Rmom}.
\end{equation}
This does not assert the absence of physically massless fields; such an
interpretation requires a Hamiltonian, dispersion law, and physical field
assignment.
\end{corollary}

\begin{proof}
The Dirac spectrum in Theorem~\ref{p2:thm:dirac-spectrum} is $\pm(k+3/2)/\Rmom$ with $k\ge 0$, so $0$ is not an eigenvalue and the smallest absolute eigenvalue is $3/(2\Rmom)$.
\end{proof}


\begin{remark}[Lichnerowicz lower bound versus the exact gap]
On a closed Riemannian spin manifold \((M,g)\), the Lichnerowicz identity
implies for every Dirac eigenvalue \(\lambda\) that
\[
\lambda^2\ge\frac14\min_M\scal_g.
\]
For the round \(S^3(\Rmom)\), one has
\(\scal_g=6/\Rmom^2\), so
\[
|\lambda|\ge\frac{\sqrt{3/2}}{\Rmom}.
\]
This lower bound proves positivity and excludes zero modes, but it is weaker
than the exact first absolute eigenvalue \(3/(2\Rmom)\), which comes from
Theorem~\ref{p2:thm:dirac-spectrum}, not from the Lichnerowicz estimate alone.
\end{remark}


\section{Higher-spin and canonical operator spectra}
\label{paper2d:sec:organized-spin}

\subsection{Higher Spin Spectra}
\label{p2:sec:higher-spin}


\begin{definition}[Spin-$s$ Fields on $S^3$]
\label{p2:def:spin-s-fields-s3}
Write
\[
S^3\cong G/K,
\qquad
G=\Spin(4),
\qquad
K=\Spin(3)_{\mathrm{diag}}.
\]
A pure spin-\(s\) field on \(S^3\) is a section of the homogeneous vector
bundle associated to the irreducible \(K\)-representation \([s]\):
\begin{itemize}
    \item \textbf{Spin-0} (scalar): functions
    \(\phi\in C^\infty(S^3)\).
    \item \textbf{Spin-1/2} (spinor): sections of the spinor bundle
    \(\Sigma\to S^3\).
    \item \textbf{Spin-1} (vector): one-forms
    \(A\in\Omega^1(S^3)\), equivalently vector fields after using the
    round metric.
    \item \textbf{Spin-3/2} (Rarita--Schwinger): sections of the
    \(\gamma\)-traceless vector-spinor bundle
    \[
      S_{3/2}
      :=
      \ker\!\left(
        \gamma:T^*S^3\otimes\Sigma\longrightarrow\Sigma
      \right).
    \]
    \item \textbf{Spin-2} (tensor): trace-free symmetric two-tensors
    \(h\in\Gamma(S^2_0T^*S^3)\).  The transverse--traceless spectral
    sector used below additionally satisfies \(\nabla^a h_{ab}=0\).
\end{itemize}
Thus an arbitrary spinor-valued one-form or symmetric two-tensor belongs to
an ambient reducible bundle and is not, without the displayed projection, a
pure spin-\(3/2\) or spin-\(2\) field.
\end{definition}

The standard round-\(S^3\) scalar, spinor, vector,
transverse-traceless tensor, and Rarita--Schwinger spectra are owned by the
following labeled statements.  Their operator normalizations and sector
conditions are part of those statements; no later summary table is used as a
premise.  Every ``gap'' in this section means a gap in the spectrum of the
displayed spatial or Dirac-type operator.  None is by itself a particle mass,
rest energy, or physical threshold.

\subsubsection{Spin-1: Vector Fields and 1-Forms}

\begin{theorem}[Spectrum of Hodge Laplacian on 1-Forms]
\label{p2:thm:spin1-spectrum}
On the round 3-sphere $S^3(\Rmom)$ the Hodge Laplacian $\Delta_1 = d d^\dagger + d^\dagger d$ acting on $1$-forms splits under the Hodge decomposition into exact and coexact parts, with spectra
\begin{align}
\mathrm{Spec}(\Delta_1|_{\mathrm{exact}}) &= \left\{ \frac{k(k+2)}{\Rmom^2} : k = 1, 2, 3, \ldots \right\},\\
\mathrm{Spec}(\Delta_1|_{\mathrm{coexact}}) &= \left\{ \frac{(k+1)^2}{\Rmom^2} : k = 1, 2, 3, \ldots \right\}.
\end{align}
The exact eigenspace at level $k$ has dimension $(k+1)^2$, while the coexact eigenspace at level $k$ has dimension $2k(k+2)$. These are standard vector-harmonic results. \cite{besse1978manifolds,Camporesi1991}
\end{theorem}

\begin{proof}
On $S^3$, the Hodge decomposition for $1$-forms is $\Omega^1 = d\Omega^0 \oplus d^\dagger\Omega^2$ since the space of harmonic $1$-forms vanishes ($H^1(S^3)=0$). 

\textbf{Exact branch.}
Let $f$ be a scalar eigenfunction with $\Delta_{S^3} f = \frac{k(k+2)}{\Rmom^2} f$ and $k\ge 1$. Then
\[
\Delta_1(df)=d(\Delta_{S^3} f)=\frac{k(k+2)}{\Rmom^2}\,df,
\]
so $df$ is an eigen-$1$-form with the same eigenvalue. The scalar eigenspace has dimension $(k+1)^2$. For $k\ge1$, no nonzero scalar eigenfunction is constant, so $d:E_k\to\Omega^1(S^3)$ is injective and $\dim(dE_k)=(k+1)^2$. Constants occur only in the $k=0$ scalar eigenspace and do not enter the exact branch listed here.

\textbf{Coexact branch.}
For a coexact $1$-form $\varpi$ with $d^\dagger\varpi=0$, in dimension $3$ the coexact branch is conventionally described by the curl operator.  With the self-adjoint curl convention
\[
C:=\star d\quad\text{on coexact one-forms},
\]
up to the harmless sign fixed by the Hodge-star convention one has \(\Delta_1\varpi=C^2\varpi\).  The curl spectrum on $S^3(\Rmom)$ is $\mathrm{Spec}(C)=\{\pm\frac{k+1}{\Rmom}:k\ge 1\}$, hence the coexact eigenvalues are $\frac{(k+1)^2}{\Rmom^2}$. The stated multiplicities are the standard vector-harmonic multiplicities. \cite{Camporesi1991,besse1978manifolds}
\end{proof}


\begin{theorem}[Imported retained Levi--Civita square and scalar ground]
\label{p2d:thm:canonical-retained-dirac-square}
Paper~I-B, Corollary~\ref{p1:cor:canonical-retained-dirac-square}, proves the
following geometric statement.  Let
\(J_R:\mathcal H_1\to\mathcal L_R\) be the retained smooth realization,
identified with the generated carrier by Paper~II-B, and let
\[
 \mathscr D_{\rm LC}=i(d-d^\dagger)\oplus D_{S^3}
\]
be the canonical Levi--Civita Hodge--de Rham plus spin Dirac representative.
Then the finite squared compression
\[
 K_{\rm LC}:=J_R^*\mathscr D_{\rm LC}^{2}J_R
\]
is diagonal on the retained irreducible sectors, with
\[
\begin{array}{c|c}
\text{retained sector} & K_{\rm LC}\\
\hline
(0,0)\text{, normalized constant} & 0\\
(\tfrac12,\tfrac12)\text{, exact }dE_1 & 3/\Rmom^2\\
(1,0)\oplus(0,1)\text{, Hodge-dual Killing module}
 & 4/\Rmom^2\\
(\tfrac12,0)\oplus(0,\tfrac12)\text{, lowest spinors}
 & 9/(4\Rmom^2)
\end{array}
\]
Hence the zero eigenspace of \(K_{\rm LC}\) is exactly the scalar line.
Write
\[
 \Pi_{\rm sc}:=\mathbf1_{\{0\}}(K_{\rm LC})
\]
for its spectral projector.  This is a canonical geometric ground projector
inside the retained smooth realization; it is not by itself a physically
prepared vacuum state or a physical Hamiltonian ground.
\end{theorem}

\begin{proof}
Paper~I-B derives the four entries from its Hodge and Dirac spectra and proves
that the scalar entry is the only zero.  Paper~II-B supplies the
same-parent identification of the GP2-generated carrier with that retained
smooth model.  The final sentence is a type firewall: neither upstream result
contains a physical preparation or Hamiltonian-owner constructor.
\end{proof}


\begin{corollary}[The minimal first-order closure preserves the retained readout]
\label{p2d:cor:v215-closure-projected-readout}
For any already constituted retained geometric readout
\(\operatorname{Tr}_{15}F(K_{\rm LC})\),
Theorem~\ref{p1:thm:v215-closure-return} provides the exact expression
\[
 \operatorname{Tr}_{15}F(K_{\rm LC})
   =\operatorname{Tr}_{25}\bigl(PF(\widetilde D^2)P\bigr).
\]
Thus the first-order closure can represent that readout without changing
its retained trace or the specified owner of its square. A coefficient
\(1/2\) would scale the geometric heat-trace difference to
\[
 \frac12\operatorname{Tr}_{25}e^{-t\widetilde D^2}
 -\frac12\operatorname{Tr}_{15}e^{-tK_{\rm LC}}
       =2e^{-3t/R^2}+3e^{-4t/R^2}.
\]
This is a comparison of the displayed operators. The native III-B
momentum action has its actual operator \(D_p\); using this particular
formula in that action additionally requires a certificate identifying
its actual squared operator and profile with the displayed
\(K_{\rm LC}\) readout. No such identification follows from the
common carrier, a principal symbol, or the role scalarizer alone.
The nonzero difference does not show that the existing action makes
that replacement. Likewise,
the Schur term \(L/z\) is the exact first-order partner return on its
stated comparison space; interpreting the comparison parameter as a
physical clock or populating its ten partner modes requires the actual
operator, state, and comparison map.
\end{corollary}
\begin{proof}
The first equality is the projected functional-calculus identity. The
second follows from the full-trace difference with \(F(s)=e^{-ts}\)
and multiplication by \(1/2\). No change of trace is used in
the first equality.
\end{proof}


\begin{theorem}[No same-symbol zero-order rigidity]
\label{p2d:thm:same-symbol-zero-order-no-go}
Locality, self-adjointness, \(\Spin(4)\)-equivariance, and the principal symbol
of \(\mathscr D_{\rm LC}\) do not determine its square or its ground space.
Indeed, for every \(m\in\mathbb R\),
\[
 \mathscr D_m
 :=i(d-d^\dagger)\oplus
   \left(D_{S^3}+\frac{m}{\Rmom}I\right)
\]
has the same principal symbol and the same coordinate commutators as
\(\mathscr D_{\rm LC}\).  On the two original retained \(k=0\) spinor
sectors, its squared eigenvalues are
\[
 \frac{(m+3/2)^2}{\Rmom^2},\qquad
 \frac{(m-3/2)^2}{\Rmom^2}.
\]
These are retained-sector values; for general \(m\) they need not be the
lowest values on the full smooth spinor bundle.  At \(m=\pm3/2\), an
additional retained spin sector reaches zero.  Therefore
the spectrum and unique scalar ground of
Theorem~\ref{p2d:thm:canonical-retained-dirac-square} are derived only for the
explicitly typed canonical Levi--Civita representative, not for the entire
same-symbol first-order class.
\end{theorem}

\begin{proof}
The added term is a local self-adjoint equivariant bundle endomorphism and
commutes with every scalar multiplier, so it changes neither the principal
symbol nor \([\mathscr D_m,M_f]\).  It shifts the original retained
\(k=0\) Dirac eigenvalues \(\pm3/(2\Rmom)\) by \(m/\Rmom\).
Squaring gives the displayed
values and the zero modes at \(m=\mp3/2\).  Thus operators satisfying all
listed same-symbol hypotheses can have different retained spectra and ground
multiplicities.  Paper~II-F, Section~\ref{p2d:app:form-zero-order-counterfamily} gives a
second counterfamily that also annihilates constants and retains the operator's
odd grading under exterior-degree parity.
\end{proof}

\begin{remark}[Basic Hopf zero-mode accounting]
\label{p2:rem:basic-hopf-zero-mode-accounting}
In the regular \emph{untwisted} \(q=0\) Hopf-basic Dolbeault complex, the
Paper-I-A theorem gives one scalar ground state and no basic \((0,1)\) harmonic
zero mode.  Thus that untwisted basic complex contains no additional
transverse base zero-mode multiplet.  This statement does not apply
automatically to \(L^q\)-twisted or fixed charged sectors; their zero modes
must be determined from the corresponding twisted operator or
charge-sector spectrum.
\end{remark}


\subsubsection{Spin-2: Transverse-traceless symmetric tensor representatives}

The spin-2 TT formula is a convention-sensitive imported round-sphere spectrum. The spin-3/2 formula is instead independently derived from the round projected Levi--Civita operator in Paper~II-F, Theorem~\ref{p2f:thm:v260-signed-rs-s3}. Neither formula is derived from the finite-mode constitution clauses; operator and radius normalizations remain explicit hypotheses.


\begin{theorem}[Imported TT Lichnerowicz spectrum on the round $S^3$]
\label{p2:thm:spin2-spectrum}
With the positive convention
\[
(\Delta_L h)_{ab}:=-\nabla^c\nabla_c h_{ab}+2R_{c(a}h_{b)}{}^c-2R_{acbd}h^{cd},
\]
using the curvature sign for which the unit round sphere has positive sectional curvature, the Lichnerowicz Laplacian $\Delta_L$ acting on transverse-traceless symmetric 
2-tensors on $S^3(\Rmom)$ has the standard TT tower
\begin{equation}
\mathrm{Spec}(\Delta_L|_{\mathrm{TT}})
=
\left\{
\frac{L(L+2)+4}{\Rmom^2}:L=2,3,4,\ldots
\right\}.
\end{equation}
The minimum eigenvalue is
\begin{equation}
\lambda_{\min}^{TT}=\frac{12}{\Rmom^2}
\qquad\text{at }L=2.
\end{equation}
This is a standard tensor-harmonic spectrum imported with the displayed convention; the convention is part of the statement.  No derivation of the general TT tensor-harmonic classification is claimed here, and the proof below uses the cited standard tensor-harmonic spectrum as an external input rather than as a consequence of the preceding Paper~II kinematic constitution clauses.
\end{theorem}

\begin{proof}
This is the Paper-I-A TT Lichnerowicz spectrum in the same positive convention
\cite{Kugelmann2026a}.  The statement is retained here solely as the
kinematic source of the exact TT gap.
\end{proof}

\begin{proposition}[Exact TT symmetric-tensor mode count]
\label{p2:prop:tt-mode-count}
For the TT tower of Theorem~\ref{p2:thm:spin2-spectrum}, the multiplicity at level $L\ge2$ is
\begin{equation}
 d_L^{\mathrm{TT}}=2(L-1)(L+3).
\end{equation}
For \(\lambda\ge0\), the exact hard counting function is
\begin{equation}
N_{\mathrm{TT}}(\lambda)=0
\qquad
\text{when }
0\le\lambda<\frac{12}{\Rmom^2}.
\end{equation}
For \(\lambda\ge12/\Rmom^2\), define
\begin{equation}
M_{\mathrm{TT}}(\lambda)
:=
\left\lfloor\sqrt{\Rmom^2\lambda-3}-1\right\rfloor,
\end{equation}
so that \(M_{\mathrm{TT}}(\lambda)\ge2\), and then
\begin{equation}
N_{\mathrm{TT}}(\lambda)
=\sum_{L=2}^{M_{\mathrm{TT}}(\lambda)}2(L-1)(L+3)
=\frac{
M_{\mathrm{TT}}(\lambda)
\bigl(M_{\mathrm{TT}}(\lambda)-1\bigr)
\bigl(2M_{\mathrm{TT}}(\lambda)+11\bigr)
}{3}.
\end{equation}
In particular,
\begin{equation}
N_{\mathrm{TT}}(\lambda)=\frac23\,\Rmom^3\lambda^{3/2}+O(\lambda).
\end{equation}
Thus, if the dimensionful Weyl law is written
\[
N_{\mathrm{TT}}(\lambda)
=C_{\mathrm{TT}}\lambda^{3/2}+O(\lambda),
\]
then
\[
C_{\mathrm{TT}}=\frac23\,\Rmom^3.
\]
Equivalently, after the factor \(\Rmom^3\) is extracted explicitly, the
remaining dimensionless numerical coefficient is \(2/3\).
\end{proposition}

\begin{proof}
The multiplicity, exact counting function, and Weyl coefficient are
the Paper-I-A TT counting result~\cite{Kugelmann2026a}.  Paper~II imports
those geometric counts and uses them only for finite kinematic occupancy
comparisons.
\end{proof}

\begin{corollary}[Finite tensor/scalar spectral occupancy ratio]
\label{p2:cor:finite-tt-scalar-ratio}
At the TT level cutoff
\[
\Lambda_L:=\frac{L(L+2)+4}{\Rmom^2},\qquad L\ge2,
\]
the scalar modes are retained through scalar level \(L\), while the TT modes
are retained through TT level \(L\).  Let \(N_0(\lambda)\) denote the scalar
hard counting function and let \(N_{\mathrm{TT}}(\lambda)\) be the hard
counting function of Proposition~\ref{p2:prop:tt-mode-count}.  The finite
tensor/scalar occupancy ratio at the common spectral cutoff is therefore
\begin{equation}
\mathcal R_{T/S}(L)
:=\frac{N_{\mathrm{TT}}(\Lambda_L)}{N_0(\Lambda_L)}
=\frac{2L(L-1)(2L+11)}{(L+1)(L+2)(2L+3)}
=2-\frac{12(4L+1)}{(L+1)(L+2)(2L+3)}.
\end{equation}
In particular,
\begin{equation}
\mathcal R_{T/S}(L)=2-\frac{24}{L^2}+\frac{102}{L^3}+O(L^{-4}).
\end{equation}
\end{corollary}

\begin{proof}
At the cutoff $\Lambda_L$, the scalar level $L$ is included because $L(L+2)/\Rmom^2<\Lambda_L$, while level $L+1$ is excluded since $(L+1)(L+3)/\Rmom^2>\Lambda_L$.  Hence
\[
N_0(\Lambda_L)=\sum_{\ell=0}^L(\ell+1)^2=\frac{(L+1)(L+2)(2L+3)}6.
\]
The TT count through level $L$ is
\[
N_{\mathrm{TT}}(\Lambda_L)=\sum_{\ell=2}^L2(\ell-1)(\ell+3)=\frac{L(L-1)(2L+11)}3.
\]
Dividing gives the stated ratio and simplifying against the asymptotic value $2$ gives the finite-cutoff suppression term.  The expansion is the large-$L$ expansion of the displayed rational function.
\end{proof}


\begin{corollary}[Lowest TT kinematic gap]
\label{p2:cor:lowest-tt-gap}
The lowest transverse--traceless tensor harmonics on $S^3(\Rmom)$ occur at $L=2$ and have eigenvalue
\begin{equation}
\lambda_{\min}^{TT}=12/\Rmom^2.
\end{equation}
Thus the TT tensor sector has a strictly positive kinematic gap.  No
particle or mass interpretation follows from this eigenvalue statement.
\end{corollary}

\begin{proof}
This is immediate from Theorem~\ref{p2:thm:spin2-spectrum}, which gives the TT spectrum and shows that the lowest allowed TT level is $L=2$.
\end{proof}


\subsubsection{Spin-3/2: Rarita-Schwinger Field}

\begin{theorem}[Signed Rarita--Schwinger spectrum on the round $S^3(\Rmom)$]
\label{p2:thm:spin32-spectrum}
Let $\slashed{D}_{3/2}$ denote the Rarita--Schwinger operator acting on $\gamma$-traceless
spinor-valued $1$-forms on $S^3(\Rmom)$, i.e.\ on sections $\psi$ satisfying
$\gamma^i\psi_i=0$. Then the spectrum splits into two series:
\begin{align}
\lambda^{(A)}_k &= \pm \frac{k+\tfrac52}{\Rmom},
& \operatorname{mult}_{A,\pm}(k) &= (k+1)(k+4),
& k&=0,1,2,\dots, \\
\lambda^{(B)}_k &= \pm \frac{k+\tfrac52}{3\Rmom},
& \operatorname{mult}_{B,\pm}(k) &= (k+2)(k+3),
& k&=0,1,2,\dots,
\end{align}
where the A-series belongs to \(\ker P^*\), the B-series belongs to
\(\operatorname{Im}P\), and every displayed multiplicity is a complex
\emph{branch multiplicity per sign}.  It is not the total full-operator
multiplicity when the two branches share a numerical eigenvalue.  Such
collisions occur precisely at
\[
\lambda^{(A)}_a=\lambda^{(B)}_{3a+5},
\qquad a=0,1,2,\ldots,
\]
and at those values the complex full-bundle multiplicity per sign is
\[
\begin{aligned}
\operatorname{mult}_{\Gamma(S_{3/2})}
\!\left(\pm\frac{a+\tfrac52}{\Rmom}\right)
&=(a+1)(a+4)+(3a+7)(3a+8)\\
&=10(a+2)(a+3).
\end{aligned}
\]
The signed spectrum on the smooth round parent is derived independently from its projected Levi--Civita operator in Theorem~\ref{p2f:thm:v260-signed-rs-s3} of Paper~II-F. It is not a consequence of the finite momentum module. Consequently, the gap on the full \(\gamma\)-traceless vector-spinor bundle is
\begin{equation}
|\slashed{D}_{3/2}|_{\min,\,\Gamma(S_{3/2})}=\frac{5}{6\Rmom}.
\end{equation}
If one restricts instead to the transverse Rarita--Schwinger summand \(\ker P^*\), the first gap is
\begin{equation}
|\slashed{D}_{3/2}|_{\min,\,\ker P^*}=\frac{5}{2\Rmom}.
\end{equation}
The sector convention is part of the statement.  The smaller full-bundle
value \(5/(6\Rmom)\) comes from the twistor-image summand; for the transverse
Rarita--Schwinger sector \(\ker P^*\), the first gap is
\(5/(2\Rmom)\).
\end{theorem}

\begin{proof}
Theorem~\ref{p2f:thm:v260-signed-rs-s3} proves the signed eigenvalues
and multiplicities by diagonalizing the projected Levi--Civita operator
in the homogeneous round frame. Proposition~\ref{p2:prop:rs-spectrum-cross-check}
independently identifies the twistor-image and transverse representation
families. The displayed gap and collision formulas are then finite arithmetic. 
Related parent and technical calculations are stated in Papers~I-A and II-F.
\end{proof}

\subparagraph{Independent spectrum cross-check.}
The Rarita--Schwinger sector inventory and multiplicities also admit an
independent representation-theoretic/twistor cross-check; the extended proof
is deferred to Paper~II-F, Section~\ref{p2d:app:rs-crosscheck}.  The transverse A-series is now independently obtained from the smooth
round operator by Theorem~\ref{p2f:thm:v260-signed-rs-s3};
the earlier representation-only cross-check did not itself fix that eigenvalue.


\subsubsection{Finite clock spectra and the scalar chiral Casimir}
\label{p2d:sec:v180-clock-kinematics}

\begin{proposition}[Level-one momentum Hopf weights]
\label{p2d:prop:level-one-hopf-weight-multiplicities}
On the fixed full momentum carrier \(\mathcal H_p=\mathcal H_1\),
use the primitive left generator \(Q_H:=Q_{H,L}=-2L_3^-\).
Then
\[
 \mathcal H_1|_{U(1)_H}
 =7\chi_0\oplus3\chi_1\oplus3\chi_{-1}\oplus\chi_2\oplus\chi_{-2},
 \qquad Q_H(Q_H^2-I)(Q_H^2-4I)=0.
\]
Its multiplicities at \(q=-2,-1,0,1,2\) are \(1,3,7,3,1\).
In this left convention the extreme weights \(q=\pm2\) occur in
\(V_{0,1}\); for the geometric right convention they occur in \(V_{1,0}\).
They are absent from the scalar harmonic submodule \(E_0\oplus E_1\).
\end{proposition}

\begin{proof}
On \(V_{j_L,j_R}\), the charge is \(-2m_R\) with multiplicity \(2j_L+1\).
The six summands contribute, in the order
\((0,0),(\tfrac12,0),(0,\tfrac12),
(\tfrac12,\tfrac12),(1,0),(0,1)\),
\[
 \chi_0,\quad 2\chi_0,\quad
 \chi_{-1}+\chi_1,\quad 2\chi_{-1}+2\chi_1,\quad
 3\chi_0,\quad \chi_{-2}+\chi_0+\chi_2.
\]
Summing gives the stated character and hence the polynomial identity.
\end{proof}

\begin{corollary}[Joint central-grading and Hopf-weight decomposition]
\label{p2d:cor:v260-joint-central-hopf-weights}
On the frozen level-one carrier and the left primitive Hopf generator
$Q_H=-2L_3^-$, the central-even/central-odd charge multiplicities are
\[
\begin{array}{c|ccccc}
 & q=-2&q=-1&q=0&q=1&q=2\\\hline
\mathcal H_+ &1&2&5&2&1\\
\mathcal H_- &0&1&2&1&0
\end{array}.
\]
In particular the $q=\pm2$ eigenspaces are entirely central-even.
This is a joint classification of the \emph{already frozen} carrier,
not a physical charge-character selector or a new momentum constructor.
\end{corollary}
\begin{proof}
Central-even summands are $V_{0,0}$, $V_{1/2,1/2}$,
$V_{1,0}$ and $V_{0,1}$. They contribute respectively
$q=0$ (one), $q=\pm1$ (two each), $q=0$ (three), and
$q=0,\pm2$ (one each). The central-odd summands
$V_{1/2,0}$ and $V_{0,1/2}$ contribute $q=0$ (two)
and $q=\pm1$ (one each). Add the weights.
\end{proof}


\begin{proposition}[Generated finite Hopf-charge algebra]
\label{p2d:prop:generated-hopf-charge-algebra}
For $F_H=\{-2,-1,0,1,2\}$ put
\[
 P_q=\prod_{\substack{r\in F_H\\r\ne q}}
 \frac{Q_H-rI}{q-r}.
\]
Then
\[
 \mathcal C_H=C^*(I,Q_H)=\bigoplus_{q\in F_H}\mathbb CP_q\cong\mathbb C^5.
\]
Every faithful unital representation of this five-projector algebra has
dimension at least five.  The inherited normalized parent trace is weighted
by $(1,3,7,3,1)/15$ and therefore is not replaced by the uniform trace on a
minimal five-dimensional presentation.
\end{proposition}

\begin{proposition}[Complete-position Hopf multiplicities]
\label{p2d:prop:complete-position-hopf-multiplicities}
For
$\mathcal P_x(m_x)=\bigoplus_{k=0}^{m_x}E_k^x$ with
$E_k^x\cong V_{k/2}\otimes V_{k/2}^*$, degree $k$ has weights
$-k,-k+2,\ldots,k$, each with multiplicity $k+1$.  Hence
\[
 \operatorname{spec}(Q_H|_{\mathcal P_x(m_x)})=\{-m_x,\ldots,m_x\}.
\]
For $|q|\le m_x$, writing
$r=\lfloor(m_x-|q|)/2\rfloor$, the exact multiplicity is
\[
 d_{m_x}(q)=(r+1)(|q|+1+r),
\]
and
\[
 D_{m_x}:=\dim\mathcal P_x(m_x)
 =\frac{(m_x+1)(m_x+2)(2m_x+3)}6.
\]
Thus $m_p=1$ does not impose a universal bound on later position-clock
weights.
\end{proposition}

\begin{corollary}[Chiral scalar Casimir identity]
\label{p2d:cor:chiral-scalar-casimir-identity}
Let \(a>0\), \(m\in\ZZ_{\ge0}\), and
\[
 \mathcal P_x(m)=\bigoplus_{k=0}^{m}E_k^x,\qquad
 E_k^x\cong V_{k/2,k/2},
\]
be a complete finite scalar position band on \(S^3(a)\).
Use the left-action convention of
Corollary~\ref{p2d:cor:hopf-rotation-transvection}, and set
\[
 Q_i:=Q_{H,L,i}^{x}=-2L_{i,x}^-.
\]
For the nonnegative scalar Laplacian restricted to this band,
\[
 \boxed{\sum_{i=1}^{3}Q_i^2=a^2\Delta_a.}
\]
For every spectral weight \(W=w(\Delta_a)\) on the band,
\[
 \operatorname{Tr}(WQ_iQ_j)
 =\frac{a^2\delta_{ij}}{3}\operatorname{Tr}(W\Delta_a).
\]
The same identities hold for the explicitly chosen right chiral triple
\(Q_i=2L_{i,x}^+\).  They concern complete scalar bands and chiral
generators; no clock-energy calibration or Hamiltonian premise is used.
\end{corollary}

\begin{proof}
On \(E_k^x\), the chosen chiral factor has spin \(j=k/2\), so
\[
 \sum_iQ_i^2=4j(j+1)I=k(k+2)I=a^2\Delta_a|_{E_k^x}.
\]
The complete band is invariant under the chiral action, so the finite direct
sum proves the operator identity without a compression defect.
The weight \(W\) commutes with that action; hence
\(T_{ij}:=\operatorname{Tr}(WQ_iQ_j)\) is an invariant two-tensor for
the adjoint \(SO(3)\) action and equals \(c\delta_{ij}\).
Summing its diagonal entries gives
\(3c=a^2\operatorname{Tr}(W\Delta_a)\).
\end{proof}

\begin{corollary}[Finite reciprocal-position momentum and its cotangent principal symbol]
\label{p2d:cor:finite-spatial-momentum-principal-symbol}
On the round reciprocal-position sphere $S^3(a)$, let
$\mathcal V_{i,a}^{L}$ be the left-action Killing fields with the Pauli-angle
normalization of Corollary~\ref{p2d:cor:hopf-rotation-transvection}.  They have
round length $a$, so
\[
 e_i:=a^{-1}\mathcal V_{i,a}^{L}
\]
is an orthonormal invariant frame.  Let $D_i=-i\hbar e_i$ be the unbounded first-order operator on
the smooth reciprocal-position parent and let $\mathsf P_{x,m}$ be
the orthogonal projector to the complete scalar band
$\mathcal P_x(m)=\bigoplus_{k=0}^{m}E_k^x$. Define the finite operator
\begin{equation}
 \label{p2d:eq:finite-spatial-momentum-operator}
 \boxed{\Pi_{i,m}^{(a)}:=\mathsf P_{x,m}D_i\mathsf P_{x,m}
                    \big|_{\mathcal P_x(m)}
       =\frac{\hbar}{a}Q_i\big|_{\mathcal P_x(m)}.}
\end{equation}
As $e_i$ is Killing, $[D_i,\Delta_a]=0$; hence the complete band is
invariant and this compression is exactly the restriction of $D_i$.
The dimensionful \emph{inherited-parent} symbol is defined by
\begin{equation}
 \label{p2d:eq:finite-spatial-momentum-symbol}
 \boxed{\sigma_{\rm inherited}(\Pi_{i,m}^{(a)})(x,p)
             :=\sigma_1(D_i)(x,p)=p(e_i(x)),
        \qquad p=\hbar\xi\in T_x^*S^3(a).}
\end{equation}
Thus $\sum_i\sigma_{\rm inherited}(\Pi_{i,m}^{(a)})^2
=\|p\|_{h^{-1}}^2$, whereas the distinct \emph{finite operator}
identity is
\[
 \boxed{\sum_i(\Pi_{i,m}^{(a)})^2
                 =\hbar^2\Delta_a\big|_{\mathcal P_x(m)}.}
\]
In particular its degree-$k$ eigenvalue defines the \emph{spectral
shell scale}
\[
 \boxed{\mathfrak p_k:=\frac{\hbar}{a}\sqrt{k(k+2)}}
 \qquad(0\le k\le m),
\]
not a pointwise cotangent covector assigned to every vector in the shell.
The extension $\mathsf P_{x,m}D_i\mathsf P_{x,m}$ to the
whole smooth parent is finite-rank smoothing, so its \emph{intrinsic}
first-order pseudodifferential principal symbol vanishes: the inherited
symbol is explicitly extra parent-comparison data.
The right-chiral convention gives the same identities in its oriented
frame. Neither the finite reciprocal-position operators nor their
parent symbols identify them with the microscopic compact-carrier
transvections $\widehat P_i$ on $\mathcal H_1$.
\end{corollary}

\begin{proof}
Corollary~\ref{p2d:cor:hopf-rotation-transvection} identifies the
left-action operator with $-i\mathcal V_i^L$; rescaling radius gives
$D_i=(\hbar/a)Q_i$. A Killing field commutes with the round scalar
Laplacian, hence $D_i$ preserves each $E_k^x$, establishing the exact
restriction and its chiral Casimir square. The ordinary parent
first-order symbol of $D_i$ is $\hbar\xi(e_i)=p(e_i)$; frame
orthonormality gives the quadratic symbol identity. Restricting the
finite square to $E_k^x$ yields
$\hbar^2k(k+2)/a^2$.
The finite-rank extension has a smooth kernel (the finite band is a sum
of smooth eigenspaces), and is therefore smoothing. Its intrinsic
order-one symbol is zero, proving the necessity of the inherited label.
\end{proof}


\section{Finite exceptional spectra, Fourier samplers and cubature}
\label{p2d:sec:finite-exceptional-comparisons}
The following operators and finite sampling maps require their stated
binary-group packets, marked axes and scalar comparison bands.  They
are not spectral Hamiltonians or physical event measurements without
separate source and readout constructions.

\subsection{Signed Petersen spectral sectors and operator scope}
\label{p2d:sec:v318-petersen-spectral-types}
\begin{theorem}[Exact finite operator sectors, conditional on the $2I$ packet]
\label{p2d:thm:v318-petersen-commutant-spectrum}
For an occurring embedded $2I$ packet, the signed Hopf-axis matrix
$S_{vw}=\operatorname{sgn}(u_v\cdot u_w)$ on Petersen-adjacent
axes and $0$ otherwise has the purely geometric spectrum
\[
 \operatorname{Spec}(S)=\{(\sqrt5)^{(3)},0^{(4)},
                             (-\sqrt5)^{(3)}\}.
\]
On the real fifteen-edge module
$W_P=\mathbf4\oplus\mathbf5\oplus\mathbf3\oplus\mathbf3'$, every
self-adjoint $A_5$-equivariant endomorphism is necessarily
\[
 \boxed{T=t_4\Pi_4+t_5\Pi_5+t_3\Pi_3+t_{3'}\Pi_{3'},
       \qquad t_a\in\mathbb R.}
\]
Thus its degeneracies are constrained to the sector dimensions
$4,5,3,3$ before accidental coincidences; symmetry does not fix
the four scalar coefficients.  If one also \emph{imposes} the
non-native full abstract graph automorphism group $S_5$, then
$t_3=t_{3'}$, because an odd permutation exchanges those two
$A_5$ representations while preserving the other two sectors.
\end{theorem}
\begin{proof}
The signed-matrix identity $S^3=5S$ and its projector ranks are
Paper-I-B finite identities.  The four irreducible summands of the
edge module are pairwise inequivalent real-type $A_5$ modules.
Real Schur orthogonality therefore makes the self-adjoint
commutant scalar on each summand.  On oriented edges, the
signed edge-permutation matrix of any odd label permutation
interchanges the two rank-three projectors and preserves the
rank-four and rank-five projectors; this gives the last statement.
\end{proof}

\begin{remark}[No physical spectra from incidence alone]
\label{p2d:rem:v318-petersen-energy-firewall}
The finite eigenvalues $\sqrt5,0,-\sqrt5$ and the four free
coefficients $t_a$ belong to an auxiliary marked-group comparison
operator, not to the round momentum Laplacian, the Hopf-charge
operator, a physical kinetic Hamiltonian, or a hadron mass operator.
The $A_5$ quotient acts trivially on the binary center; obtaining
spinorial parity needs the separately constructed binary lift.
No new energy gap, charge, chirality or physical generation follows
from this section.  The existing parent-to-finite spectral recovery
and Hopf-charge assignments are unchanged.
\end{remark}


\subsection{Twelve-fibre binary-icosahedral Fourier resolution}
\label{p2d:sec:cc14-hopf-120}

\begin{theorem}[Finite binary-icosahedral Hopf Fourier sampler]
\label{p2d:thm:cc14-hopf-120}
Fix a binary-icosahedral comparison subgroup $2I\subset SU(2)$,
choose its order-ten element $h=\exp(i\pi\sigma_3/5)$ after
conjugating the left Hopf axis to this element, and choose twelve
representatives $g_r$ of the left cosets $C_{10}g_r$,
$C_{10}=\langle h\rangle$.  Thus
$2I=\bigsqcup_{r=1}^{12}\{h^kg_r:0\le k<10\}$.
Let $\mathcal P_5=\bigoplus_{n=0}^{5}\mathcal H_n(S^3)$ with normalized
Haar inner product, so $\dim_\mathbb C\mathcal P_5=91$.
The evaluation map
\[
 Sf=\bigl(120^{-1/2}f(h^kg_r)\bigr)_{r,k}
 \in\mathbb C^{12}\otimes\mathbb C^{10}
\]
is an isometry.  If $f=\sum_{q=-5}^{5}f_q$ is its left-fibre Hopf-charge
resolution, $f_q(h_\alpha g)=e^{iq\alpha}f_q(g)$, then the unitary
length-ten discrete Fourier transform in $k$ gives
\[
 \boxed{\widehat S_{r,a}f=
 \frac1{\sqrt{12}}\sum_{q\equiv a\ (\mathrm{mod}\,10)}f_q(g_r),
 \quad a=0,\ldots,9.}
\]
The associated ten $12\times12$ evaluation Gram blocks $B_a$ are
orthogonal projections of exact ranks
\[
 \boxed{(\operatorname{rank}B_a)_{a=0}^{9}
 =(9,12,8,10,5,12,5,10,8,12).}
\]
In particular $B_1=B_5=B_9=I_{12}$, the full $120$-sample
Gram is an orthogonal projector of rank $91$ and nullity $29$,
and the $a=5$ block reconstructs the \emph{joint} twelve-dimensional
$q=5$ and $q=-5$ sector despite their aliasing on each single fibre.
\end{theorem}
\begin{proof}
The existing degree-eleven $2I$ spherical-design identity in
Theorem~\ref{p1b:thm:v349-single-degree12-alias} applies to
$\overline fg$ for $f,g\in\mathcal P_5$, because its ambient
polynomial degree is at most ten.  Hence $S^*S=I_{91}$.
Under the left-fibre convention of
Definition~\ref{p2:def:finite-diagonal-spin-weight-hopf-charge},
$\mathcal H_n$ carries charges $q=-n,-n+2,\ldots,n$, each of
multiplicity $n+1$.  Since $h^k$ acts on charge $q$ by
$\exp(2\pi ikq/10)$, Fourier orthogonality gives the displayed
formula with coefficient $10/(\sqrt{10}\sqrt{120})=1/\sqrt{12}$.
Every Fourier residue block is the evaluation of a pairwise
orthogonal subspace, and the full Fourier-transformed evaluation
remains an isometry, so $B_a=\widehat S_a\widehat S_a^*$ is an
orthogonal projector of rank
\[
 d_a=\sum_{\substack{|q|\le5\\q\equiv a\ (\mathrm{mod}\,10)}}
 \sum_{\substack{|q|\le n\le5\\n\equiv q\ (\mathrm{mod}\,2)}}(n+1).
\]
This yields the integer rank list above.  Each block lives on twelve
coset coordinates, so $d_a=12$ implies $B_a=I_{12}$.  The
$a=5$ residue contains only $q=5,-5$, each of dimension six;
its square isometric evaluation therefore separates those
subspaces globally.  Summing the ranks gives $91$, leaving
$120-91=29$ null directions.
\end{proof}

\begin{corollary}[Exact order-ten Hopf-charge filtering and first alias]
\label{p2d:cor:cc14-charge-filter}
On $\mathcal P_9=\bigoplus_{n=0}^9\mathcal H_n$, the cyclic average
$\Pi_{10}=10^{-1}\sum_{k=0}^{9}U_L(h)^k$ is exactly the continuous
left-Hopf charge-zero projector: $\Pi_{10}|_{\mathcal P_9}=P_{q=0}$.
For all $\mathcal P_m$ with $m\le4$, every individual Hopf charge
is resolved without fibre aliasing by the ten-point Fourier transform.
At degree $n=5$ the first pair of charges aliases, namely $q=\pm5$,
but the theorem gives its full global de-aliasing.
\end{corollary}
\begin{proof}
The finite cyclic projector retains exactly charges divisible by ten.
On $\mathcal P_9$ the available charges satisfy $|q|\le9$, so only
$q=0$ survives.  On $\mathcal P_m$ with $m\le4$, distinct charges
differ by at most eight and cannot be congruent modulo ten.
For $m=5$ the only distinct congruent pair in $[-5,5]$ is $-5,5$.
\end{proof}

\begin{remark}[Geometric gauge and actualization firewall]
A chosen $h$ and the coset representatives constitute a marked
comparison atlas, not a generated preferred Hopf direction.
The twelve base positions represent Hopf-bundle sampling locations;
nonzero-charge data transform under local frame changes, so the result
does not trivialize charged Hopf line bundles.  The $91$-dimensional
band is a \emph{scalar comparison} band, not the physical $15$-mode
retained carrier, the full fermionic Fock space, or a new
superselection assertion.  The source already establishes generic
positive cubature frame bounds; the exact order-ten Fourier block diagonalization and rank/alias
certificate is established for the stated marked comparison packet.
\end{remark}


\subsection{Seventy-two-base cubature and charge-separated scalar readout}
\label{p2d:sec:cc20-72-hopf-sampler}
The marked order-ten axis and scalar bands in this base rule are
part of the stated comparison input and need not coincide with the
marked packet used in the twelve-fibre theorem.

\begin{theorem}[Positive seventy-two-base degree-nine cubature]
\label{p2d:thm:cc20-72-base9}
Fix the marked order-ten binary-icosahedral Hopf fibre subgroup in the
120-vertex/600-facet degree-nineteen rule.  The vertex/facet nodes
fall into twelve and sixty distinct fibres of ten points, with
seventy-two distinct base directions on $S^2$.  The resulting base rule
is exact through spherical polynomial degree nine and has positive
weights $5/252$ per vertex-base and $4/315$ per facet-base.  The
full scalar band $\mathcal P_9(S^3)$ has dimension $385$, and the
unitary ten-point fibre DFT of the weighted 720-point evaluation map
splits into ten charge-residue blocks, of ranks
\[
 \boxed{(25,40,33,46,37,48,37,46,33,40),\qquad\sum=385.}
\]
Thus charge aliases on individual ten-point fibres are resolved by
the full set of seventy-two base positions.  This is an exact
\emph{complex-amplitude} statement, not an intensity-only tomography
or a physically preferred Hopf direction.
\end{theorem}
\begin{proof}
The existing positive 720-node degree-nineteen cubature is exact for
pullbacks of spherical polynomials of degree at most nine, since the
Hopf projection has quadratic quaternion coordinates.  The two
fibre masses are ten times the respective 720-node point weights.
The unitarity of the fibre DFT and isometry of the 720-node scalar
band give orthogonal residue blocks.  Counting the source Hopf
charges in $\bigoplus_{n=0}^9\mathcal H_n(S^3)$ gives the displayed
ranks, whose sum equals the complete band dimension.
\end{proof}
\begin{proposition}[Fixed-support impossibility of improving these rules]
\label{p2d:prop:cc20-fixed-support-no-go}
For a marked fivefold base direction let $H_6,H_{10}$ denote the
normalized $A_5$-invariant zonal harmonics taking orbit-average
values $(1,-5/16)$ and $(1,31/64)$ on the twelve and sixty base
orbits.  Then
\[
 q_{10}=17+11H_6-28H_{10}
\]
vanishes at every one of the seventy-two base points but has
normalized spherical integral $17$.  Hence \emph{no weights, even
signed}, on this fixed support can achieve degree-ten exactness.
Analogously $Q_{20}=17+11C_{12}-28C_{20}$ vanishes at every vertex and
facet centre of the parent 720-point support while its Haar integral
is $17$.  The degree-nineteen exactness on that support therefore
cannot be improved by reweighting alone.
\end{proposition}
\begin{proof}
The stated orbit values are the exact source Legendre and
binary-icosahedral character moments.  The linear polynomial
combination vanishes on both support types, whereas every
nonconstant harmonic integrates to zero.  Any purported improved
cubature would assign the identically zero sample sum to a function
with nonzero integral, a contradiction.
\end{proof}


\subsection{Conditional two-sided Hopf charge decomposition of 600-chamber memory}
\label{cc19:sec:two-hopf-charges}
This finite comparison is conditional on the 41-state representation
of Theorem~\ref{p1b:thm:cc16-minimal-41state-memory}; it supplies
no physical occurrence, preparation or decay-time operator.

\begin{theorem}[Sixteen nonzero joint charges and one universal two-pole return]
\label{cc19:thm:joint-memory-charges}
Let $h\in2I$ have quaternion order ten with $h^5=-1\in2T$,
acting from left and right on the 600 facet centres and on their
25 registers. The two induced register permutations are independent
5-cycles. Over $\mathbb C$, the register carrier has one character
for each $(a,b)\in\mathbb Z_5\times\mathbb Z_5$. The stationary
character is $(0,0)$, the eight one-sided characters have exactly one
vanishing index, and the sixteen correlated characters have $a,b\ne0$.
All sixteen correlated characters have the same exact comparison-return
factor
\[
\boxed{f_c(t)=\left(\frac12+\frac3{4\sqrt{21}}\right)
   e^{(-9+\sqrt{21})t/8}
 +\left(\frac12-\frac3{4\sqrt{21}}\right)
   e^{(-9-\sqrt{21})t/8}.}
\]
The existing minimal 41-state linear realization decomposes into
one stationary, eight additive and sixteen visible/hidden two-state
blocks. The charge residues on the order-ten fibres are even modulo
ten: $2a,2b\in\{2,4,6,8\}$. This is not a physical decay-clock claim.
\end{theorem}
\begin{proof}
The central element $h^5=-1$ lies in $2T$, hence acts trivially on
the five cosets of $2T$ in $2I$. The image of $h$ acts freely and
transitively on each five-coset label set, giving two commuting
regular $C_5$ actions. The Fourier decomposition is their tensor
product. The three register projectors have ranks $1,8,16$ and are
the indicator projections for the three stated Fourier subsets.
The source's correlated visible/hidden generator is
$\bigl(\begin{smallmatrix}-15/16&5\sqrt3/16\\5\sqrt3/16&-21/16\end{smallmatrix}\bigr)$.
Diagonalization gives eigenvalues $(-9\pm\sqrt{21})/8$ and the
stated visible spectral weights; tensoring with each of the sixteen
one-dimensional charge characters proves the block decomposition.
\end{proof}


\section{Spectral operators and downstream comparison interfaces}
\label{paper2d:sec:organized-exports}

\subsection{Quantitative spectral-owner and normalization transfers}
\label{p2d:sec:v2210-lit-owner-transfers}
The exact scalar fixed-band recovery and round spectra remain unchanged.
The following certificates address only stronger observable comparisons,
actual kinetic normalization and the meaning of a determinant readout.

\begin{theorem}[Observable-specific positive comparison bound]
\label{p2d:thm:v2210-lit-positive-recovery}
Let $\mathcal A,\mathcal B$ be supplied finite unital $C^*$-algebras and
let $\Phi:\mathcal A\to\mathcal B$, $\Psi:\mathcal B\to\mathcal A$ be
licensed unital completely positive comparison maps.  For a state
$\omega$ on $\mathcal A$, set $\omega_{\mathcal B}=\omega\circ\Psi$.
For every consumed observable $A$,
\[
 |\omega_{\mathcal B}(\Phi(A))-\omega(A)|
 \le\|\Psi\Phi(A)-A\|.
\]
For a consumed word $W=A_1\cdots A_n$, put
$W_\Phi=\Phi(A_1)\cdots\Phi(A_n)$.  Then
\[
 |\omega_{\mathcal B}(W_\Phi)-\omega(W)|
 \le\|W_\Phi-\Phi(W)\|+\|\Psi\Phi(W)-W\|.
\]
Thus reconstruction and multiplication errors are separate and jointly
sufficient for the stated expectation transfer.  Zero errors give exact
transfer without an all-resolution injective map into $\End(\mathcal H_1)$.
\end{theorem}
\begin{proof}
A state has norm one.  Subtract the two evaluated operators to prove the
first bound.  For the second, add and subtract $\omega\Psi\Phi(W)$ and
use that a unital completely positive map is contractive.  In finite
dimension the latter follows from its Schwarz inequality:
$\Psi(X)^\dagger\Psi(X)\le\Psi(X^\dagger X)\le\|X\|^2I$.
\end{proof}
These are the native fixed-observable hypotheses required before using
operator-system spectral recovery technology
\cite{LitGEOVanSuijlekom2020Truncation}.  The theorem proves the expectation
certificate; it does not assert a full quantum-metric convergence theorem
without its own uniform metric estimates.  For compressed comparisons,
Theorem~\ref{p2b:thm:v2210-lit-word-defect} computes the first error.
Actual source derivatives and moving comparison maps are included in the
differentiated word; an undifferentiated expectation bound is not a jet bound.

For the native two-dimensional center, the explicit sector-tagged comparison
of Paper~II-C,
Corollary~\ref{p2c:cor:v2211-tagged-intrinsic-populations}, evaluates both
errors in this theorem to zero for every central word.  Its maps are a
multiplicative classical tag embedding and an explicit unital completely
positive retraction.  This supplies a completed finite comparison on that
claim algebra.  It does not extend the certificate to noncentral source
words, repair the quaternion rule's quadratic aliasing, or provide an
injective all-resolution map into the frozen momentum carrier.

\begin{proposition}[A spatial gap transfers only through its actual kinetic pair]
\label{p2d:prop:v2210-lit-kinetic-gap}
On a specified retained physical source subspace $\mathcal K$, let
$Z>0$, $K=K^\dagger$, and $L\ge\lambda_*I$, with $\lambda_*>0$.
Suppose the actual inherited kinetic/source forms satisfy
\[
 Z\le z_*I,\qquad K\ge aL+bZ,
 \qquad z_*>0,\quad a\ge0.
\]
Then the normalized kinetic endomorphism obeys
\[
 Z^{-1/2}KZ^{-1/2}\ge
 \left(b+\frac{a\lambda_*}{z_*}\right)I.
\]
If the same physical clock/action identifies its eigenvalues with
$\omega^2$ and the right-hand side is positive, the frequency gap is at
least $\sqrt{b+a\lambda_*/z_*}$ in that clock normalization.
\end{proposition}
\begin{proof}
For $x\ne0$,
$x^\dagger Kx/(x^\dagger Zx)\ge
b+a\lambda_*\|x\|^2/(x^\dagger Zx)\ge b+a\lambda_*/z_*$.
Apply this Rayleigh-quotient estimate after $y=Z^{1/2}x$.
\end{proof}
The constants are bounds on the actual consumed forms, not new mass
parameters.  The gapped source subspace, $Z$, and the physical operator
pencil must be supplied; the round eigenvalue alone supplies none of them.

\begin{theorem}[Connection and mass terms preserve a supplied local Dirac cone]
\label{p2d:thm:v2210-lit-lower-order-cone}
On an already soldered local four-dimensional Lorentzian chart, let
\[
 \mathscr D=\gamma^ae_a^{\ \mu}(\partial_\mu+\mathcal A_\mu)+V,
 \qquad \{\gamma^a,\gamma^b\}=2\eta^{ab}I,
\]
where $\mathcal A_\mu,V$ act with differential order zero and the coframe
is nondegenerate.  The principal symbol, with the common scalar Fourier
factor suppressed, is
\[
 \sigma_1(\mathscr D)(k)=\gamma^ae_a^{\ \mu}k_\mu,
 \qquad \sigma_1(\mathscr D)(k)^2
 =g^{\mu\nu}k_\mu k_\nu I.
\]
On the rank-four complex spinor block its determinant is
$(g^{\mu\nu}k_\mu k_\nu)^2$.  Consequently these connection/mass
interactions preserve the supplied characteristic cone on this local domain.
\end{theorem}
\begin{proof}
Only the derivative term contributes to the first symbol.  Symmetry of
$k_\mu k_\nu$ removes the Clifford commutator and gives the square.
The two signed Clifford eigenvalues each have multiplicity two on the
rank-four block; taking their product gives the determinant polynomial,
including its null locus by continuity.
\end{proof}
This narrows the interacting-cone obligation for the stated lower-order
branch.  A higher-derivative or environment-Schur kinetic return must have
its own full principal-symbol test.  The theorem imports the existing
physical solder and Lorentzian real form; it does not derive them from a
positive compact Born metric or impose a noncompact unitary action on H1.

\begin{theorem}[Determinant scale response and finite part have different data]
\label{p2d:thm:v2210-lit-zeta-scale}
Let $L$ be a supplied positive operator on its nonzero-mode domain whose
spectral zeta function is regular at zero; this includes a finite positive
block and the usual regular elliptic zeta domains.  With a fixed reference
unit and $c>0$,
\[
 \zeta_{cL}(s)=c^{-s}\zeta_L(s),\qquad
 \log\det{}'_{\zeta}(cL)
 =-\zeta_L'(0)+\zeta_L(0)\log c.
\]
For the Gaussian determinant word $W_G(c)=\tfrac12\log\det{}'_{\zeta}(cL)$,
\[
 \frac{\mathrm dW_G}{\mathrm d\log c}=\tfrac12\zeta_L(0),
 \qquad W_G(1)=-\tfrac12\zeta_L'(0).
\]
\end{theorem}
\begin{proof}
The zeta identity follows on its convergent half-plane by rescaling every
nonzero eigenvalue and elsewhere by the same analytic continuation.
Differentiate at zero and use
$\log\det{}'_{\zeta}L=-\zeta_L'(0)$
\cite{Vassilevich2003HeatKernel}.  Ordinary differentiation in $\log c$
gives the last two identities.
\end{proof}
For a finite block $\zeta_L(0)$ is its nonzero rank, whereas
$-\zeta_L'(0)$ is the sum of logarithms of its nonzero eigenvalues.
For example $\operatorname{diag}(1,1)$ and
$\operatorname{diag}(1,e^d)$ have the same value $\zeta(0)=2$ and different
log determinants $0,d$.  A compiler whose defined output is the logarithmic
scale germ can therefore consume $\zeta(0)$ exactly.  A claimed finite
determinant value additionally consumes $\zeta'(0)$ and its actual
reference/subtraction convention.  Common operator ancestry is compatible
with both outputs and does not identify them.  This diagnostic changes
neither a separately fixed shell atom nor a hierarchy constant; a physical
use retains the precise determinant/germ word and complete action jets.


\subsection{Curvature canonical exchange as an imported response interface}
\label{p2d:sec:v229-curvature-polarization}
This interface imports the result of Paper~III-F, not a
new II-D derivation. On the regular local metric-$f_c(R)$ cap chart of
Theorem~\ref{p3f:thm:v223-cap-parent}, put $\Phi=f_c'(r)$ and
$\chi=f_c''(r)>0$. Corollary~\ref{p3f:cor:v260-canonical-exchange}
proves the exact canonical exchange, with a \emph{constant} $\eta\ne0$,
\[
 P_c=\eta r,\qquad Q=-\chi(r)p_\Phi/\eta,\qquad
 p_\Phi\,d\Phi=P_c\,dQ-d(P_cQ).
\]
The notation $P_c$ distinguishes this momentum coordinate from the
Einstein coefficient $P_E$ of the capped density. The phase action
retains the $-[P_cQ]$ endpoint change, the full pulled-back Hamiltonian
and all source/boundary terms of its actual owner; the resulting
quadratic term contains $Q^2/\chi^2$ and a $p_aQ$ cross term rather
than an unchanged $P_c^2$ flux.
Proposition~\ref{p3f:prop:v260-jordan-flux} derives the Jordan lapse
constraint and canonical boundary momentum from that same restricted
action and exhibits on-shell slices of equal curvature and induced
metric with opposite scalar momenta. Corollary~
\ref{p3f:cor:v260-conormal-jordan-distinction} evaluates the distinct
conormal source norm and its necessary/sufficient amplitude match.
Neither the coordinate exchange nor a fixed curvature profile proves
a physical incoming conormal amplitude or the occurrence of the cap
in the full common action; those identifications keep their explicit
III-E/III-F source-owner obligations.


\section{Conclusions and mathematical export limits}
Paper~II-D takes the round parent and the fixed momentum role
$\mathcal H_p\cong\mathcal H_1$, $m_p=1$ as constituted inputs.
The smooth parent spectra, the finite-operator calculations and the
reciprocal-position band comparisons have different domains; none is a
second momentum constructor. The scalar fixed-band recoveries, guarded
operator words and round-metric local approximation have their own
comparison maps and quantitative error conditions.

The central $\Spin(4)$ grading yields the geometric/spinorial split
$11+4$, even and odd endomorphism dimensions $137$ and $88$ and the
absence of a nonzero invariant odd bridge. Geometric parity and time
reversal are transport and fixed-locus statements: the unmarked parity
locus is $(r,r,\varrho_{\rm sp})$, whereas unmarked time reversal fixes
the full completion quotient. The marked Hopf connection gives an
integer-character atlas and $q\leftrightarrow-q$ transport, but no
selected physical axis, sign or charge; this is not a physical
charge-conjugation or CPT theorem.

The canonical Levi--Civita retained Hodge/spin square has sector
values $0$, $3/\Rmom^2$, $4/\Rmom^2$ and $9/(4\Rmom^2)$, with a unique
scalar ground. Its lower-order data matter: the same principal symbol
alone fixes neither the spectrum nor a physical vacuum or kinetic law.
The conditional Petersen, Hopf-sampling, cubature and finite-memory
comparisons likewise determine mathematical operator structure, not
an observed energy gap, mass, instrument or decay clock.
\begin{remark}[Downstream operator and recovery requirements]
The momentum carrier is already fixed upstream; later dynamics need not
re-select \(\mathcal H_p=\mathcal H_1\).  Whenever the same frozen object is
retained, Paper~II-A supplies \((\mathrm{FRZ}_p)\) by chronological
persistence rather than by a new branch assumption.  Reciprocal position uses the
derived mathematical tower of Paper~II-C and, for physical applications, an
actual frontier history or the direct condition
\((\mathrm{BAND}_x[K,t])\).  Recovery is stated by its comparison map,
topology, and remainder \(\varepsilon_{\rm rec}\); on the scalar fixed-band
and guard-banded word branches that remainder is zero.  Round metric
flatness is the separate explicit error \(\varepsilon_g(L,R_x)\) of
Proposition~\ref{p2d:prop:exact-round-flatness}.  Nonlinear terminal-shell
products, propagation, and dynamics require their own controlled remainder
or resolvent estimates.  A non-scalar or nonlinear recovery claim therefore
retains an explicit comparison map, domain, topology, and remainder.  A
physical spectral-energy claim separately retains its dispersion/operator
identification and calibration, while a continuum renormalized value retains
its subtraction scheme, reference scale, and allowed local counterterms.
These are downstream evolution/recovery/interpretation questions, not new
momentum-constitution problems.
\end{remark}

\paragraph{Physical scale and kinetic data not inferred.}
The structural kinematic data do not determine downstream numerical or
dynamical identifications:
\[
 \boxed{\begin{gathered}
 \mathcal H_p,\ m_p,\ \text{round spectral data},\ \text{Hopf/Chern data}\\
 \not\Rightarrow\\
 \Rmom^{\rm numerical},\ H_{\rm kin}^{\rm physical},\
 \text{particle mass/threshold},\\
 \text{same-owner field ancestry or a vacuum-light identification}.
 \end{gathered}}
\]
Any such statement requires its own generated dynamical bridge or calibration.

\paragraph{Downstream sharp-frontier recovery.}
Paper~III-F, Corollary~\ref{p3:cor:spectral-resolution-thresholds},
generates complete-band admission at
$a_L=\sqrt{L(L+2)}/\Lambda_F$ from an occurring sharp frontier and scale
history. Paper~VI-A, Proposition~\ref{p8:prop:large-mx-smooth-limit},
therefore discharges an independent exhaustive-history premise on the
nondecreasing unbounded expansion branch. Finite-band recovery only needs
the reached band (and the stated guard for a nonlinear word); its input
need not be recounted as a new recovery selector. Stronger interacting,
uniform or continuum claims retain their own estimates. Primitive
structure remains on compact momentum at $m_p=1$.


\clearpage

\printbibliography
\end{document}
